\thispagestyle{plain}
\ctsection{index vocum quæ in foro militari, civili sacroque obtinent}

\par{\dictfontsize\hangindent=0pt\noindent\textls[120]{\scshape a}\par}
\par{\dictfontsize\hangindent=1.5em\hangafter=1\noindent Abbé (\textit{d’un couvent}).—\textit{Gardien (chez les Capucins).} \lxlinks{Præfectus monachorum, monasterii rector, princeps monachorum; abbas.}\par}
\par{\dictfontsize\hangindent=1.5em\hangafter=1\noindent Abbesse (\textit{d’un couvent}); \textit{supérieure} (\textit{d’une simple congrégation}.) \lxlinks{Rectrix virginum DEO sacrarum, quæ collegio sacrarum virginum \textit{vel} virginibus DEO sacratis præest, antistes sacrarum virginum; abbatissa.}\par}
\par{\dictfontsize\hangindent=1.5em\hangafter=1\noindent Absolution \textit{du prêtre}. \lxlinks{Oratio, solemnis precatio, qua sacerdos Numinis vice reo peccatorum veniam impertit; absolutio.}\par}
\par{\dictfontsize\hangindent=1.5em\hangafter=1\noindent Absoudre \textit{qqn}; \textit{remettre les péchés à qqn}. \lxlinks{Vinculis criminum absolvere aliquem, culpæ expertem pronuntiare aliquem.}\par}
\par{\dictfontsize\hangindent=1.5em\hangafter=1\noindent Accise \textit{ou} Douane, \textit{impôt sur les denrées, sur la consommation}. \lxlinks{Portorium, vectigal eorum quæ consumuntur, tributum rerum ad victum et amictum pertinentium.}\par}
\par{\dictfontsize\hangindent=1.5em\hangafter=1\noindent Accise, \textit{impôt sur les cheminées}. \lxlinks{Exactio focorum.}\par}
\par{\dictfontsize\hangindent=1.5em\hangafter=1\noindent \textit{Caisse de l’}accise, \textit{des contributions indirectes}. \lxlinks{Fiscus vectigalium.}\par}
\par{\dictfontsize\hangindent=1.5em\hangafter=1\noindent \textit{Mettre l’}accise \textit{sur chaque mesure de vin}. \lxlinks{In singulas vini urnas nomine statam pecuniam exigere.}\par}
\par{\dictfontsize\hangindent=1.5em\hangafter=1\noindent Receveur \textit{de l’}accise, \textit{des contributions indirectes}. \lxlinks{Curator census e victu persolvi soliti.}\par}
\par{\dictfontsize\hangindent=1.5em\hangafter=1\noindent Accrédité (\textit{ambassadeur ou envoyé}); \textit{un ambassadeur muni de lettres de créance, de pleins pouvoirs}. \lxlinks{Legatus mandatis et plena agendi potestate instructus.}\par}
\par{\dictfontsize\hangindent=1.5em\hangafter=1\noindent Accuser \textit{qqn d’un crime capital}. \lxlinks{Judicio capitis arcessere aliquem, arguere aliquem rei capitalis, in judicium capitis vocare.}\par}
\par{\dictfontsize\hangindent=1.5em\hangafter=1\noindent Actions (\textit{finances}); \textit{mises} (\textit{fin.}). \lxlinks{Litteræ quibus cavetur de pecunia societati credita.}\par}
\par{\dictfontsize\hangindent=1.5em\hangafter=1\noindent Accuser \textit{qqn de sédition},.... \textit{de révolte}. \lxlinks{Actionem perduellionis intendere, majestatis postulare.}\par}
\par{\dictfontsize\hangindent=1.5em\hangafter=1\noindent \textit{Voir, à l’aide de l’}addition \textit{et de la soustraction, ce qui reste}. \lxlinks{Addendo deducendoque experiri quid reliqui sit.}\par}
\par{\dictfontsize\hangindent=1.5em\hangafter=1\noindent Adjudant, \textit{aide-de-camp}; \textit{officier d’ordonnance}. \lxlinks{Adjutor castrensis, vigilum præfecti adjutor.}\par}
\par{\dictfontsize\hangindent=1.5em\hangafter=1\noindent Administrateur \textit{de la caisse}, \textit{des fonds d’une ville}; \textit{trésorier}; \textit{payeur}. \lxlinks{Ærarii civici tribunus.}\par}
\par{\dictfontsize\hangindent=1.5em\hangafter=1\noindent \textit{N’avoir aucun rapport d’intérêts, d’}affaires \textit{avec}..., \textit{ne pas faire d’}affaires, \textit{le commerce}; \textit{ne pas entretenir, n’avoir pas de relations commerciales avec}......; \textit{ne pas traiter d’affaires avec}.... \lxlinks{Nullo commercii jure misceri.}\par}
\par{\dictfontsize\hangindent=1.5em\hangafter=1\noindent \textit{Avoir des rapports d’intérêts ou d’affaires avec qqn}; \textit{faire des} affaires, \textit{le commerce avec qqn}. \lxlinks{Res rationesque contrahere cum aliquo.}\par}
\par{\dictfontsize\hangindent=1.5em\hangafter=1\noindent Agent, \textit{chargé d’affaires}. \lxlinks{Curator negotiorum publicus, actor publicus.}\par}
\par{\dictfontsize\hangindent=1.5em\hangafter=1\noindent \textit{Général-adjudant}; aide-de-camp. \lxlinks{Adjutor ducis.}\par}
\par{\dictfontsize\hangindent=1.5em\hangafter=1\noindent L’aile \textit{droite}, \textit{l’aile gauche} (\textit{d’une armée}). \lxlinks{Ala, cornu, dextra, sinistra.}\par}
\par{\dictfontsize\hangindent=1.5em\hangafter=1\noindent \textit{Droit d’}aînesse. \lxlinks{Jus natu majoris, vel maximi.}\par}
\par{\dictfontsize\hangindent=1.5em\hangafter=1\noindent Ajournement \textit{d’une cause}. \lxlinks{Ampliatio judicii.}\par}
\par{\dictfontsize\hangindent=1.5em\hangafter=1\noindent \textit{Sonner l’}alarme. \lxlinks{Ad arma conclamare.}\par}
\par{\dictfontsize\hangindent=1.5em\hangafter=1\noindent \textit{Franc}-alleu, \textit{c. à d.}, \textit{un bien libre, franc de tous droits}. \lxlinks{Vectigalium onere immune.}\par}
\par{\dictfontsize\hangindent=1.5em\hangafter=1\noindent \textit{Bien} allodial. \lxlinks{Prædium \textit{vel} fundus a clientelari nexu liber.}\par}
\par{\dictfontsize\hangindent=1.5em\hangafter=1\noindent \textit{Bien} allodial, \textit{terre allodiale}. \lxlinks{Prædium a clientelarum nexu liberum.}\par}
\par{\dictfontsize\hangindent=1.5em\hangafter=1\noindent Almanach, \textit{calendrier}. \lxlinks{Fasti.}\par}
\par{\dictfontsize\hangindent=1.5em\hangafter=1\noindent \textit{Faiseur d’}almanachs. \lxlinks{Fastorum auctor.}\par}
\par{\dictfontsize\hangindent=1.5em\hangafter=1\noindent Alternativement. \lxlinks{Alternis vicibus.}\par}
\par{\dictfontsize\hangindent=1.5em\hangafter=1\noindent Ambassadeur, \textit{envoyé}, \textit{député}. \lxlinks{Legatus, orator, nuntius solemnis.}\par}
\par{\dictfontsize\hangindent=1.5em\hangafter=1\noindent Ambassadeur \textit{de premier ordre}. \lxlinks{Legatus primi ordinis.}\par}
\par{\dictfontsize\hangindent=1.5em\hangafter=1\noindent \textit{Résident}, ambassadeur, \textit{chargé d’affaires (d’un État) auprès d’une puissance étrangère}. \lxlinks{Qui regis \textit{aut} reipublicæ negotia agit.}\par}
\par{\dictfontsize\hangindent=1.5em\hangafter=1\noindent Amiral. \lxlinks{Præfectus classis, prætor navalis, qui classi præest. Summus classium præfectus. Dux militiæ navalis, rei maritimæ præfectus.}\par}
\par{\dictfontsize\hangindent=1.5em\hangafter=1\noindent \textit{Vaisseau} amiral. \lxlinks{Navis prætoria.}\par}
\par{\dictfontsize\hangindent=1.5em\hangafter=1\noindent Amirauté. \lxlinks{Summus rei navalis senatus.}\par}
\par{\dictfontsize\hangindent=1.5em\hangafter=1\noindent Amnistie. \lxlinks{Lex oblivionis, dictorum factorumque omnium abolitio; plebiscitum, ne qua præteritarum rerum mentio fiat.}\par}
\par{\dictfontsize\hangindent=1.5em\hangafter=1\noindent \textit{Être à l’}ancre. \lxlinks{In anchoris navem tenere, in anchoris exspectare.}\par}
\par{\dictfontsize\hangindent=1.5em\hangafter=1\noindent \textit{Lever l’}ancre. \lxlinks{Anchoras tollere.}\par}
\par{\dictfontsize\hangindent=1.5em\hangafter=1\noindent Angles \textit{d’un fort, d’un bastion}. \lxlinks{Rostrum propugnaculi.}\par}
\par{\dictfontsize\hangindent=1.5em\hangafter=1\noindent Annales. \lxlinks{Fasti.}\par}
\par{\dictfontsize\hangindent=1.5em\hangafter=1\noindent \textit{J’ai de l’}antipathie, \textit{j’éprouve de l’aversion pour un tel}. \lxlinks{Animo dissideo ab hoc.}\par}
\par{\dictfontsize\hangindent=1.5em\hangafter=1\noindent Apanage, \textit{revenus annuels destinés à l’entretien d’un noble}. \lxlinks{Redituum portio, principibus natu minoribus assignata.}\par}
\par{\dictfontsize\hangindent=1.5em\hangafter=1\noindent \textit{Seigneurs apanagés}; \textit{princes apanagés ou} apanagistes; \textit{grands (du royaume), hauts fonctionnaires apanagés}. \lxlinks{Privati principes.}\par}
\par{\dictfontsize\hangindent=1.5em\hangafter=1\noindent Appariteur, \textit{huissier}, \textit{garçon de bureau}. \lxlinks{Apparitor, minister magistratus.}\par}
\par{\dictfontsize\hangindent=1.5em\hangafter=1\noindent \textit{S’}approcher \textit{de l’ennemi}; \textit{s’avancer vers l’ennemi}. \lxlinks{Castra castris conferre.}\par}
\par{\dictfontsize\hangindent=1.5em\hangafter=1\noindent \textit{Contre}-approches. \lxlinks{Excursus, procursus obsessorum.}\par}
\par{\dictfontsize\hangindent=1.5em\hangafter=1\noindent Archives \textit{privées}. \lxlinks{Tabularia sanctiora.}\par}
\par{\dictfontsize\hangindent=1.5em\hangafter=1\noindent \textit{Les} archives, \textit{actes publics}. \textit{(Se dit des archives elles-mêmes et aussi du local où elles sont conservées.) Chartrier (d’un monastère).} \lxlinks{Chartophylacium, scrinium cum litteris, memoria principum publica, tabularium rerum gestarum, tabulæ publicæ, locus in quo tabulæ et litteræ publicæ servantur.}\par}
\par{\dictfontsize\hangindent=1.5em\hangafter=1\noindent \textit{Consigner dans les} archives, \textit{inscrire sur les registres publics}. \lxlinks{Ad tabulas referre.}\par}
\par{\dictfontsize\hangindent=1.5em\hangafter=1\noindent \textit{Tirer de la poudre des annales ou des} archives \textit{le souvenir de qqche}. \lxlinks{Ex annalium vetustate eruere alicujus rei memoriam.}\par}
\par{\dictfontsize\hangindent=1.5em\hangafter=1\noindent Archiviste; \textit{chartrier} (\textit{d’un monastère}). \lxlinks{Custos tabularum publicarum, tabularii præfectus.}\par}
\par{\dictfontsize\hangindent=1.5em\hangafter=1\noindent Argent \textit{monnayé}. \lxlinks{Argentum signatum.}\par}
\par{\dictfontsize\hangindent=1.5em\hangafter=1\noindent Argent \textit{en lingot}, \textit{argent non travaillé}. \lxlinks{Rude, infectum.}\par}
\par{\dictfontsize\hangindent=1.5em\hangafter=1\noindent Argent \textit{comptant}. \lxlinks{Pecunia præsens, numerata.}\par}
\par{\dictfontsize\hangindent=1.5em\hangafter=1\noindent \textit{Mettre de l’}argent \textit{en circulation}. \lxlinks{Pecuniam habere implicitam, quæ in foro versatur, quæ cohæret cum aliorum pecuniis.}\par}
\par{\dictfontsize\hangindent=1.5em\hangafter=1\noindent \textit{Donner quittance} (\textit{pour de l’}argent \textit{reçu}), \textit{acquitter une note, un paiement}. \lxlinks{Pecunias alicui expensas ferre.}\par}
\par{\dictfontsize\hangindent=1.5em\hangafter=1\noindent \textit{Manque d’}argent \textit{à la chambre du trésor}. \lxlinks{Incredibiles angustiæ pecuniæ publicæ.}\par}
\par{\dictfontsize\hangindent=1.5em\hangafter=1\noindent \textit{Prendre de l’}argent \textit{à intérêt}, \textit{emprunter de l’argent à intérêt}. \lxlinks{Fenori sumere, fenore accipere.}\par}
\par{\dictfontsize\hangindent=1.5em\hangafter=1\noindent Argument, \textit{preuve solide}. \lxlinks{Argumentum gravissimum et firmissimum.}\par}
\par{\dictfontsize\hangindent=1.5em\hangafter=1\noindent \textit{Les principaux} arguments. \lxlinks{Puncta argumentorum.}\par}
\par{\dictfontsize\hangindent=1.5em\hangafter=1\noindent \textit{Réfuter, renverser les} arguments. \lxlinks{Argumentis aliquid refellere.}\par}
\par{\dictfontsize\hangindent=1.5em\hangafter=1\noindent \textit{Patron de vaisseau}; armateur. \lxlinks{Navis dominus.}\par}
\par{\dictfontsize\hangindent=1.5em\hangafter=1\noindent \textit{Mettre bas les} armes, \textit{rendre les armes};—\textit{demander quartier}. \lxlinks{Armis positis, submissis, abjectis vitæ gratiam deposcere.}\par}
\par{\dictfontsize\hangindent=1.5em\hangafter=1\noindent \textit{Prononcer un} arrêt, \textit{rendre un jugement}. \lxlinks{Sententiam dicere.}\par}
\par{\dictfontsize\hangindent=1.5em\hangafter=1\noindent \textit{Quoiqu’il} arrive. \lxlinks{Utcumque cesserit, utut evenerit.}\par}
\par{\dictfontsize\hangindent=1.5em\hangafter=1\noindent Article de foi, \textit{symbole}. \lxlinks{Caput fidei Christianæ.}\par}
\par{\dictfontsize\hangindent=1.5em\hangafter=1\noindent \textit{Préparer des couronnes foudroyantes ou d’}artifice. \lxlinks{Malleolos facesque ad incendendam urbem comparare.}\par}
\par{\dictfontsize\hangindent=1.5em\hangafter=1\noindent \textit{Jeter, lancer des couronnes foudroyantes ou d’}artifice \textit{sur les assaillants}. \lxlinks{Piceis sulphureisque sertis subeuntes induere \textit{seu} in subeuntes jacere.}\par}
\par{\dictfontsize\hangindent=1.5em\hangafter=1\noindent Artificier (\textit{t. milit.}); \textit{l’artilleur qui met la mèche au canon pour faire partir le coup}; \textit{brigadier d’artillerie}; \textit{bombardier} (\textit{t. milit.}). \lxlinks{Librator tormentorum, mortariorum, ollarum ignitarum.}\par}
\par{\dictfontsize\hangindent=1.5em\hangafter=1\noindent Artillerie, \textit{pièces d’artillerie ou simplement pièces}; (\textit{ainsi on dit: des pièces de campagne, de siége, des pièces de 6, de 12}...); \textit{batterie}. \lxlinks{Res tormentaria.}\par}
\par{\dictfontsize\hangindent=1.5em\hangafter=1\noindent \textit{Garde-magasin d’}artillerie. \lxlinks{Armamentario præfectus.}\par}
\par{\dictfontsize\hangindent=1.5em\hangafter=1\noindent \textit{Colonel}; \textit{capitaine général}, \textit{officier d’}artillerie, \textit{capitaine}, \textit{officier du genie}. \lxlinks{Præfectus, tribunus cohortis tormentariæ, præfectus fabrum.}\par}
\par{\dictfontsize\hangindent=1.5em\hangafter=1\noindent \textit{Mener les soldats à l’}assaut, \textit{les faire monter à l’assaut}. \textit{Donner l’assaut.} \lxlinks{Ad assultum educere militem.}\par}
\par{\dictfontsize\hangindent=1.5em\hangafter=1\noindent Assiéger, \textit{mettre le siége devant}.... \lxlinks{Circumvallare, includere.}\par}
\par{\dictfontsize\hangindent=1.5em\hangafter=1\noindent \textit{Réduire les} assiégés \textit{par la famine}; \textit{affamer les assiégés}. \lxlinks{Macerare fame circumvallatos.}\par}
\par{\dictfontsize\hangindent=1.5em\hangafter=1\noindent \textit{Répandre la terreur et semer la destruction parmi les} assiégés, \textit{en tirant du canon et en lançant des bombes}. \lxlinks{Magna vi pilarum fragoreque tormentorum obsessis terrorem detrimentumque inferre, injicere.}\par}
\par{\dictfontsize\hangindent=1.5em\hangafter=1\noindent \textit{Comparaître au jour assigné}; \textit{comparaître en justice}, \textit{répondre à une} assignation. \lxlinks{Venire ad constitutum diem, vadimonium obire.}\par}
\par{\dictfontsize\hangindent=1.5em\hangafter=1\noindent \textit{Fixer le jour de la comparution en justice}; \textit{lancer contre qqn un mandat de comparution}; \textit{donner une} assignation, \textit{une citation}. \lxlinks{Vadimonium concipere.}\par}
\par{\dictfontsize\hangindent=1.5em\hangafter=1\noindent Attaque, \textit{charge}; \textit{assaut}; \textit{choc}. \lxlinks{Incursio, incursus.}\par}
\par{\dictfontsize\hangindent=1.5em\hangafter=1\noindent Attaque \textit{nocturne}. \lxlinks{Nocturna grassatio.}\par}
\par{\dictfontsize\hangindent=1.5em\hangafter=1\noindent \textit{Prendre}, attaquer \textit{l’ennemi en queue}. \lxlinks{Adversos proterere.}\par}
\par{\dictfontsize\hangindent=1.5em\hangafter=1\noindent Attaquer \textit{l’ennemi de front}. \lxlinks{In adversos hostes impetum facere, impressionem dare in adversos hostes.}\par}
\par{\dictfontsize\hangindent=1.5em\hangafter=1\noindent \textit{En} attester \textit{DIEU et les hommes}. \lxlinks{Testari DEUM hominesque.}\par}
\par{\dictfontsize\hangindent=1.5em\hangafter=1\noindent \textit{Droit de retraite}, \textit{droit d’}aubaine (\textit{dr.}); \textit{droit}, \textit{taxe de détraction}. (\textit{dr.}) \lxlinks{Census emigrationis causa persolvi solitus.}\par}
\par{\dictfontsize\hangindent=1.5em\hangafter=1\noindent \textit{Donner} audience \textit{à qqn}; \textit{admettre qqn à l’audience}. \lxlinks{Admittere aliquem, audire benigne attenteque.}\par}
\par{\dictfontsize\hangindent=1.5em\hangafter=1\noindent \textit{Ne pas recevoir d’}audience, \textit{n’être pas admis à l’audience}; \textit{se voir refuser une audience}. \lxlinks{Alloquio excludi, ad colloquium non admitti.}\par}
\par{\dictfontsize\hangindent=1.5em\hangafter=1\noindent \textit{Rendre la justice}, \textit{juger}; \textit{être jour d’}audience (\textit{en parlant d’un tribunal}). \lxlinks{Jus dicere, exercere.}\par}
\par{\dictfontsize\hangindent=1.5em\hangafter=1\noindent \textit{Être admis à l’}audience \textit{de congé}. \lxlinks{Ultimum convenire principem, eum supremum discessus causa alloqui.}\par}
\par{\dictfontsize\hangindent=1.5em\hangafter=1\noindent Auditeur \textit{général}, \textit{premier juge militaire}; \textit{grand prévôt}. \lxlinks{Prætor, judex castrensis supremus, quæsitor causarum capitalium, militarium supremus.}\par}
\par{\dictfontsize\hangindent=1.5em\hangafter=1\noindent Aumônier \textit{militaire}, \textit{aumônier de l’armée} (\textit{cath.});—\textit{ministre du camp} (\textit{protest.}). \lxlinks{Sacerdos castrensis, qui læsum supremis Mysteriis procurat militem.}\par}
\par{\dictfontsize\hangindent=1.5em\hangafter=1\noindent Avant-garde. \lxlinks{Antecessores, signa prima, acies prima.}\par}
\par{\dictfontsize\hangindent=1.5em\hangafter=1\noindent \textit{Commander, conduire l’}avant-garde. \textit{Avoir le commandement de l’avant-garde.} \lxlinks{Primo pilo præesse; primum pilum ducere.}\par}
\par{\dictfontsize\hangindent=1.5em\hangafter=1\noindent \textit{Sans vous incommoder}; \textit{c’est à votre} avantage. \lxlinks{Quod commodo tuo fiat.}\par}
\par{\dictfontsize\hangindent=1.5em\hangafter=1\noindent \textit{Temps de l’}Avent. \lxlinks{Hebdomades quæ anniversariam nascentis Christi memoriam præcurrunt.}\par}
\par{\dictfontsize\hangindent=0pt\noindent\textls[120]{\scshape b}\par}
\par{\dictfontsize\hangindent=1.5em\hangafter=1\noindent Bagages. \lxlinks{Impedimenta militaria.}\par}
\par{\dictfontsize\hangindent=1.5em\hangafter=1\noindent \textit{Piller}, \textit{enlever les} bagages. \lxlinks{Impedimenta diripere.}\par}
\par{\dictfontsize\hangindent=1.5em\hangafter=1\noindent \textit{Piller les} bagages \textit{des vivandiers, des cantiniers}; \textit{enlever les fourgons des vivandiers, des cantinières}. \lxlinks{Lixarum mercatorumque sarcinas intercipere.}\par}
\par{\dictfontsize\hangindent=1.5em\hangafter=1\noindent \textit{Ciel de lit}; baldaquin \textit{quelconque} (\textit{au dessus d’un trône, d’une table, d’un autel, etc.}). \lxlinks{Tegmen pensile, aulæa suspensa.}\par}
\par{\dictfontsize\hangindent=1.5em\hangafter=1\noindent Balle, \textit{grenade}. \lxlinks{Glans, ignis missilis, malleoli, pilæ ignivomæ ignita jacula.}\par}
\par{\dictfontsize\hangindent=1.5em\hangafter=1\noindent Balle \textit{de fusil}. \lxlinks{Glans plumbea.}\par}
\par{\dictfontsize\hangindent=1.5em\hangafter=1\noindent Banque. \lxlinks{Mensa publica argentaria. Ærarium mercatorum.}\par}
\par{\dictfontsize\hangindent=1.5em\hangafter=1\noindent Banque; \textit{comptoir d’escompte}. \lxlinks{Taberna argentaria.}\par}
\par{\dictfontsize\hangindent=1.5em\hangafter=1\noindent Banquier, \textit{changeur}. \lxlinks{Numularius, mensarius, argentarius.}\par}
\par{\dictfontsize\hangindent=1.5em\hangafter=1\noindent Banquier, changeur. \lxlinks{Numularius, mensarius, argentarius fenerator.}\par}
\par{\dictfontsize\hangindent=1.5em\hangafter=1\noindent \textit{Uni avec qqn par le lien du} baptême. \lxlinks{Baptismatis fœdere conjunctus.}\par}
\par{\dictfontsize\hangindent=1.5em\hangafter=1\noindent Baraques; \textit{baraquements}. \lxlinks{Tuguriola militum; raptim excitatæ militum domus. Tabernacula militum e trabulis.}\par}
\par{\dictfontsize\hangindent=1.5em\hangafter=1\noindent Barrière; \textit{barres}. \lxlinks{Cancelli lignei.}\par}
\par{\dictfontsize\hangindent=1.5em\hangafter=1\noindent Bataillon; \textit{cohorte}; \textit{compagnie}. \lxlinks{Cohors.}\par}
\par{\dictfontsize\hangindent=1.5em\hangafter=1\noindent \textit{Former le} bataillon \textit{carré, ou former le carré}. \lxlinks{Agmen quadratum e milite constituere.}\par}
\par{\dictfontsize\hangindent=1.5em\hangafter=1\noindent Batterie. \lxlinks{Agger, tormentorum suggestus.}\par}
\par{\dictfontsize\hangindent=1.5em\hangafter=1\noindent \textit{Dresser des} batteries. \lxlinks{Suggestus tormentarios constituere, ibique tormenta collocare.}\par}
\par{\dictfontsize\hangindent=1.5em\hangafter=1\noindent \textit{Dresser des} batteries. \lxlinks{Aggeres jacere, struere.}\par}
\par{\dictfontsize\hangindent=1.5em\hangafter=1\noindent Bastion. \lxlinks{Propugnaculum, munimentum mœnium latera sustentans.}\par}
\par{\dictfontsize\hangindent=1.5em\hangafter=1\noindent Bastion; \textit{fort}. \lxlinks{Propugnaculi agger, munimentum majus; propugnaculum.}\par}
\par{\dictfontsize\hangindent=1.5em\hangafter=1\noindent Bedeau, \textit{appariteur}. \lxlinks{Minister Academiæ publicus.}\par}
\par{\dictfontsize\hangindent=1.5em\hangafter=1\noindent \textit{Dire le} benedicite. \lxlinks{Mensam sacrare, mensam cibosque appositos precibus consecrare.}\par}
\par{\dictfontsize\hangindent=1.5em\hangafter=1\noindent Bibliothèque, \textit{tablettes}; \textit{rayons}, \textit{cases} (\textit{d’une bibliothèque}). \lxlinks{Foruli, plutei, armarium.}\par}
\par{\dictfontsize\hangindent=1.5em\hangafter=1\noindent Biens-fonds \textit{affectés à l’aîné d’une famille} (\textit{en France, dans l’Amérique espagnole, etc.}). \lxlinks{Bona, patrimonia, quæ caput stirpis sequuntur, quæ natu maximo obtingunt.}\par}
\par{\dictfontsize\hangindent=1.5em\hangafter=1\noindent \textit{Fonds de terre}, \textit{bien}, biens-fonds, \textit{terre}, \textit{propriété}, \textit{domaine}. \lxlinks{Fundus.}\par}
\par{\dictfontsize\hangindent=1.5em\hangafter=1\noindent Biens meubles; \textit{mobilier}. \lxlinks{Ea quæ moveri possunt, suppellex.}\par}
\par{\dictfontsize\hangindent=1.5em\hangafter=1\noindent \textit{Faire le} bilan; \textit{balancer un compte}; \textit{faire le compte des dépenses et des recettes}; \textit{faire le compte du bilan}. \lxlinks{Accepti expensique rationes subducere.}\par}
\par{\dictfontsize\hangindent=1.5em\hangafter=1\noindent Bill; \textit{projet de loi du parlement} (\textit{en Angleterre}). \lxlinks{Scriptum a supremo Angliæ Senatu Regi propositum, ut ei calculum addat vimque legis conferat; typus condendæ legis in Anglorum Comitiis, constitutio publica senatus Anglici.}\par}
\par{\dictfontsize\hangindent=1.5em\hangafter=1\noindent Billet. \lxlinks{Codicilli, epistola minuscula, ceræ pusillæ.}\par}
\par{\dictfontsize\hangindent=1.5em\hangafter=1\noindent Billet \textit{de logement} (\textit{milit.}). \lxlinks{Tessera hospitii militaris.}\par}
\par{\dictfontsize\hangindent=1.5em\hangafter=1\noindent Biscuit. \lxlinks{Panis nauticus, bis coctus.}\par}
\par{\dictfontsize\hangindent=1.5em\hangafter=1\noindent Biscuit \textit{de mer}. \lxlinks{Panis nauticus.}\par}
\par{\dictfontsize\hangindent=1.5em\hangafter=1\noindent Blanc-seing. \lxlinks{Charta vacua manu alicujus signata.}\par}
\par{\dictfontsize\hangindent=1.5em\hangafter=1\noindent Bloquer, \textit{assiéger une ville}; \textit{faire le blocus, le siége d’une ville}. \lxlinks{Assidere urbi, corona circumdare, mœnia exercitu circumvenire.}\par}
\par{\dictfontsize\hangindent=1.5em\hangafter=1\noindent Bombarder \textit{une ville, une place}. \lxlinks{Verberare urbem tormentis, injicere in urbem pilas ferreas nitrato pulvere sartas. Totius belli apparatu ad urbem agere.}\par}
\par{\dictfontsize\hangindent=1.5em\hangafter=1\noindent Bombes. \lxlinks{Malleoli, olla ignaria ferramentis varii generis \textit{vel} grandine ferrata conferta.}\par}
\par{\dictfontsize\hangindent=1.5em\hangafter=1\noindent Bombes; \textit{obus,—boulets de canon}. \lxlinks{Bolides igneæ, pilæ igniariæ, pilæ ferreæ nitrato pulvere sartæ. Candentes globi in sublime protrusi. Missilis ignis globi. Tædæ, malleoli illiti pice: qui modus adhibebatur in oppugnationibus veterum.}\par}
\par{\dictfontsize\hangindent=1.5em\hangafter=1\noindent Bombe \textit{remplie de matières inflammables} (\textit{t. de guerre}). \lxlinks{Malleoli.}\par}
\par{\dictfontsize\hangindent=1.5em\hangafter=1\noindent \textit{Agrandir l’empire, le royaume,—reculer les} bornes \textit{de l’empire. Étendre son règne, sa domination}. \lxlinks{Proferre fines, imperium.}\par}
\par{\dictfontsize\hangindent=1.5em\hangafter=1\noindent Bourse, \textit{édifice public où s’assemblent les négociants, les agents de change... pour travailler d’affaires}. \lxlinks{Basilica mercatorum, locus mercatorum conventui destinatus, ædificium mercatorum conventui agendisque deliberationibus exstructum.}\par}
\par{\dictfontsize\hangindent=1.5em\hangafter=1\noindent Boussole. \lxlinks{Pyxis nautica, arcula magneticam acum complexa.}\par}
\par{\dictfontsize\hangindent=1.5em\hangafter=1\noindent Brèche; \textit{dans les murs, dans les remparts}; \textit{crevasse};—\textit{remparts en ruine}. \lxlinks{Valli hiatus, disjecta mœnia, ruinæ monumentorum.}\par}
\par{\dictfontsize\hangindent=1.5em\hangafter=1\noindent \textit{Battre en} brèche (\textit{les fortifications}).—\textit{Faire, ouvrir une brèche.} \lxlinks{Crebris tormentorum ictibus mœnia convellere, labefactare; ruinam mœnium edere, impressionem facere in muros.}\par}
\par{\dictfontsize\hangindent=1.5em\hangafter=1\noindent \textit{Réparer les} brèches. \lxlinks{Muros quassatos raptim objectis saxis, vimineis fascibus reficere.}\par}
\par{\dictfontsize\hangindent=1.5em\hangafter=1\noindent \textit{Battre en} brèche \textit{les remparts}. \lxlinks{Tormentis muros quatere.}\par}
\par{\dictfontsize\hangindent=1.5em\hangafter=1\noindent Bref \textit{pontifical}. \lxlinks{Litteræ a Pontifice Romano publice scriptæ \textit{vel} scribi solitæ.}\par}
\par{\dictfontsize\hangindent=1.5em\hangafter=1\noindent Bréviaire. \lxlinks{Liber precationum et hymnorum sacerdotum \textit{vel} cleri usibus destinatus. Codex precum in horas obiri solitarum.}\par}
\par{\dictfontsize\hangindent=1.5em\hangafter=1\noindent Brigade.—\textit{Escouade.} \lxlinks{Caterva, manus delectorum militum.}\par}
\par{\dictfontsize\hangindent=1.5em\hangafter=1\noindent Brigadier; \textit{chef d’escouade}. \lxlinks{Præfectus alæ, ductor catervæ.}\par}
\par{\dictfontsize\hangindent=1.5em\hangafter=1\noindent Brûlot (\textit{t. de marine}). \lxlinks{Navis incendiaria, navis inferendis incendiis parata.}\par}
\par{\dictfontsize\hangindent=1.5em\hangafter=1\noindent Bulletin \textit{de victoire}. \lxlinks{Laureatæ litteræ.}\par}
\par{\dictfontsize\hangindent=0pt\noindent\textls[120]{\scshape c}\par}
\par{\dictfontsize\hangindent=1.5em\hangafter=1\noindent \textit{Couper le} câble (\textit{de l’ancre}). \lxlinks{Rudentes succidere.}\par}
\par{\dictfontsize\hangindent=1.5em\hangafter=1\noindent Cadet (\textit{d’une famille noble}). \lxlinks{Secundo tertiove loco a nobili familia ortus.}\par}
\par{\dictfontsize\hangindent=1.5em\hangafter=1\noindent Cadet (\textit{milit.}), \textit{jeune gentilhomme qui servait de soldat pour apprendre le métier}; \textit{soldat du corps des cadets}. \lxlinks{Miles e nobilium cohorte.}\par}
\par{\dictfontsize\hangindent=1.5em\hangafter=1\noindent Caisse \textit{militaire}. \lxlinks{Ærarium militare.}\par}
\par{\dictfontsize\hangindent=1.5em\hangafter=1\noindent Calendrier, \textit{almanach}. \lxlinks{Fasti.}\par}
\par{\dictfontsize\hangindent=1.5em\hangafter=1\noindent Camp \textit{retranché}. \lxlinks{Castra opere munitissima.}\par}
\par{\dictfontsize\hangindent=1.5em\hangafter=1\noindent \textit{Lever le} camp; \textit{décamper}; \textit{se mettre en marche}. \lxlinks{Castra movere ex aliquo loco.}\par}
\par{\dictfontsize\hangindent=1.5em\hangafter=1\noindent \textit{Gens de la campagne}, campagnards, \textit{villageois, paysans}. \lxlinks{Gentiles, populares.}\par}
\par{\dictfontsize\hangindent=1.5em\hangafter=1\noindent \textit{Entrer en} campagne. \lxlinks{Copias in campum deducere, copias extrahere ex hibernis, exercitum in expeditionem deducere.}\par}
\par{\dictfontsize\hangindent=1.5em\hangafter=1\noindent \textit{Entrer, se mettre en} campagne; \textit{quitter, sortir de ses quartiers d’hiver}. \lxlinks{Excire ex hibernis copias.}\par}
\par{\dictfontsize\hangindent=1.5em\hangafter=1\noindent Campement \textit{fixe, campement, camp, cantonnement}. \lxlinks{Castra stativa.}\par}
\par{\dictfontsize\hangindent=1.5em\hangafter=1\noindent Camp-\textit{volant}. \lxlinks{Manus expedita, agmen expeditum.}\par}
\par{\dictfontsize\hangindent=1.5em\hangafter=1\noindent Camp-\textit{volant}; \textit{un petit corps} (\textit{de cavalerie}). \lxlinks{Expedita delectorum manus. Equitum globus, manus, parva equitum turma.}\par}
\par{\dictfontsize\hangindent=1.5em\hangafter=1\noindent \textit{Former un} camp-volant, \textit{la phalange} (\textit{grecque}); \textit{former, organiser un corps d’élite}; \textit{faire un choix des plus braves soldats}. \lxlinks{Phalangem constituere, deligere pugnacissimos viros.}\par}
\par{\dictfontsize\hangindent=1.5em\hangafter=1\noindent \textit{Pièce}, canon (\textit{t. d’artillerie}). \lxlinks{Tormentum bellicum, igniarium, machina muralis.}\par}
\par{\dictfontsize\hangindent=1.5em\hangafter=1\noindent \textit{Gros} canon, \textit{canon de gros calibre, pièce de siége}. \lxlinks{Tormentum æneum majus. Tormentum majoris moduli. Tormentum diverberandis muris aptum.}\par}
\par{\dictfontsize\hangindent=1.5em\hangafter=1\noindent Canon \textit{démonté, cloué}. \lxlinks{Tormenta exarmata, clavis confixa.}\par}
\par{\dictfontsize\hangindent=1.5em\hangafter=1\noindent \textit{A la portée d’un coup de} canon. \lxlinks{Ad ictum tormenti, quantum tormentum jacere potest.}\par}
\par{\dictfontsize\hangindent=1.5em\hangafter=1\noindent \textit{Tirer du} canon, \textit{faire partir des coups de canon}. \lxlinks{Tormenta excutere, ignem tormentis admovere.}\par}
\par{\dictfontsize\hangindent=1.5em\hangafter=1\noindent Canonicat, \textit{prébende}. \lxlinks{Canonia; canonica.}\par}
\par{\dictfontsize\hangindent=1.5em\hangafter=1\noindent Canoniser, \textit{mettre au nombre des saints, élever sur les autels}. \lxlinks{Sanctorum ordinibus inferre, in numerum sanctorum referre.}\par}
\par{\dictfontsize\hangindent=1.5em\hangafter=1\noindent Cantinier. \lxlinks{Mercatores annonarii, mercatores castrenses.}\par}
\par{\dictfontsize\hangindent=1.5em\hangafter=1\noindent Capitaine. \lxlinks{Centurio.}\par}
\par{\dictfontsize\hangindent=1.5em\hangafter=1\noindent Capitaine \textit{général} (\textit{en Espagne}). \lxlinks{Supremus belli dux.}\par}
\par{\dictfontsize\hangindent=1.5em\hangafter=1\noindent Capitaine \textit{commandant la garde}. \lxlinks{Præfectus pretoriano militi.}\par}
\par{\dictfontsize\hangindent=1.5em\hangafter=1\noindent Capital. \lxlinks{Caput, sors.}\par}
\par{\dictfontsize\hangindent=1.5em\hangafter=1\noindent Capital \textit{improductif}. \lxlinks{Nummi vacui, pecunia nil lucri reddens.}\par}
\par{\dictfontsize\hangindent=1.5em\hangafter=1\noindent \textit{Augmentation du} capital. \lxlinks{Pecuniæ accessio.}\par}
\par{\dictfontsize\hangindent=1.5em\hangafter=1\noindent \textit{Décréter l’impôt personnel, la} capitation. \lxlinks{Imperare tributum in capita.}\par}
\par{\dictfontsize\hangindent=1.5em\hangafter=1\noindent \textit{Capitaux, fonds}, capital. \lxlinks{Pecuniæ in fœnore positæ.}\par}
\par{\dictfontsize\hangindent=1.5em\hangafter=1\noindent Capitulation \textit{de l’empereur romain}. \lxlinks{Lex regia. Pacta in quæ princeps regimen iniens fidem obstringit.}\par}
\par{\dictfontsize\hangindent=1.5em\hangafter=1\noindent Capituler, \textit{se rendre, traiter de la capitulation, de la reddition d’une place, en arrêter les conditions}. \lxlinks{Cum hoste convenire, conventa tractare.}\par}
\par{\dictfontsize\hangindent=1.5em\hangafter=1\noindent Caporal. \lxlinks{Decurio.}\par}
\par{\dictfontsize\hangindent=1.5em\hangafter=1\noindent Carabinier. \lxlinks{Sclopus equestris.}\par}
\par{\dictfontsize\hangindent=1.5em\hangafter=1\noindent Cardinal. \lxlinks{Purpuratus Pater; purpuratus Romanus, is quem rubens galerus ornat.}\par}
\par{\dictfontsize\hangindent=1.5em\hangafter=1\noindent \textit{Collége des} Cardinaux, \textit{le sacré Collége}. \lxlinks{Senatus sacer Pontificis Romani.}\par}
\par{\dictfontsize\hangindent=1.5em\hangafter=1\noindent \textit{Mardi-gras}; carnaval. \lxlinks{Dionysia, Bacchanalia, Liberalia, dies Baccho sacri, publica comessatione exigi soliti.}\par}
\par{\dictfontsize\hangindent=1.5em\hangafter=1\noindent \textit{Connaître son} caractère, \textit{son humeur, son tempérament}. \lxlinks{Suum nosse ingenium, genium.}\par}
\par{\dictfontsize\hangindent=1.5em\hangafter=1\noindent \textit{Ceux qui ont le même} caractère, \textit{le même tempérament, le même naturel, les mêmes penchants}. \lxlinks{Qui natura consentiunt, eodem ducuntur genio.}\par}
\par{\dictfontsize\hangindent=1.5em\hangafter=1\noindent \textit{Ce qui nous répugne}; \textit{ce qui répugne à notre} caractère, \textit{à notre tempérament}; \textit{ce qui est contraire à nos penchants, à nos inclinations}. \lxlinks{Quæ non faciunt ad nostrum genium.}\par}
\par{\dictfontsize\hangindent=1.5em\hangafter=1\noindent Carroussel; \textit{course aux chevaux}; \textit{course aux chars}; \textit{tournoi}. \lxlinks{Pugna equestris ludicra, de annulo ad metam proposito. Exercitatio equestris animi causa instituta, dimicatio ludicra, equestris.}\par}
\par{\dictfontsize\hangindent=1.5em\hangafter=1\noindent Cartel. \lxlinks{Litteræ provocatoriæ ad pugnam singularem.}\par}
\par{\dictfontsize\hangindent=1.5em\hangafter=1\noindent Casemates. \lxlinks{Cæsæ criptæ. Imæ criptæ ad latera propugnaculorum. Domunculæ militares occultæ.}\par}
\par{\dictfontsize\hangindent=1.5em\hangafter=1\noindent Casuel; \textit{émoluments}. \lxlinks{Reditus fortuiti, sportulæ.}\par}
\par{\dictfontsize\hangindent=1.5em\hangafter=1\noindent Catafalque. \lxlinks{Castrum doloris.}\par}
\par{\dictfontsize\hangindent=1.5em\hangafter=1\noindent Cathédrale. \lxlinks{Ædes sacra primaria, ædes princeps.}\par}
\par{\dictfontsize\hangindent=1.5em\hangafter=1\noindent Catéchiser; \textit{enseigner à qqn les éléments de la doctrine chrétienne}. \lxlinks{Prima rudimenta sanctioris doctrinæ aliquem docere, imbuere aliquem fidei mysteriis.}\par}
\par{\dictfontsize\hangindent=1.5em\hangafter=1\noindent Cause \textit{criminelle}; \textit{chambre criminelle}. \lxlinks{Judicium capitis, causa capitis, quæstio.}\par}
\par{\dictfontsize\hangindent=1.5em\hangafter=1\noindent \textit{Être} caution \textit{de qqche, se rendre caution de qqche, répondre de qqche, garantir qqche}. \lxlinks{In rem aliquam fidem suam interponere, sponsorem, se alicujus rei præstare.}\par}
\par{\dictfontsize\hangindent=1.5em\hangafter=1\noindent \textit{Garantie}, caution, \textit{cautionnement, droit, fidéjussion, acte, écrit, par lequel on répond de qqn ou de qqche}. \lxlinks{Formula præstandæ securitatis.}\par}
\par{\dictfontsize\hangindent=1.5em\hangafter=1\noindent Cavalerie \textit{légère}; (\textit{par ex. hussards}). \lxlinks{Levis armaturæ equites.}\par}
\par{\dictfontsize\hangindent=1.5em\hangafter=1\noindent \textit{Fête} centenaire \textit{de la Fondation de l’Université, de l’Académie}. \lxlinks{Secularis Academiæ natalis.}\par}
\par{\dictfontsize\hangindent=1.5em\hangafter=1\noindent \textit{Faire une charge de} cavalerie. \lxlinks{Equestrem procellam excitare.}\par}
\par{\dictfontsize\hangindent=1.5em\hangafter=1\noindent Cérémonial. \lxlinks{Rituum aulæ et solemnium consignatio.}\par}
\par{\dictfontsize\hangindent=1.5em\hangafter=1\noindent \textit{Maître des} cérémonies. \lxlinks{Magister rituum et morum aulicorum, \textit{vel} rituum sacrorum supremus moderator.}\par}
\par{\dictfontsize\hangindent=1.5em\hangafter=1\noindent \textit{Battre la} chamade, \textit{c. à d.}, \textit{donner le signal de la reddition}; \textit{proposer de se rendre}; \textit{faire des propositions de reddition}. \lxlinks{Proclamare ad deditionem, colloquium hostium poscere.}\par}
\par{\dictfontsize\hangindent=1.5em\hangafter=1\noindent \textit{Battre la chamade; battre le tambour de la caisse pour donner le signal de la reddition.} \lxlinks{Signum dare deditionis, deditionem offerre.}\par}
\par{\dictfontsize\hangindent=1.5em\hangafter=1\noindent Chambellan. \lxlinks{Regi \textit{vel} Principi ab interioribus cubiculi præfectus.}\par}
\par{\dictfontsize\hangindent=1.5em\hangafter=1\noindent \textit{Grand} chambellan. \lxlinks{Cubiculi regii \textit{vel} Principis præfectus.}\par}
\par{\dictfontsize\hangindent=1.5em\hangafter=1\noindent Chambre \textit{des députés} (\textit{anct. la place où se tenaient les assemblées générales du peuple romain}). \lxlinks{Comitium.}\par}
\par{\dictfontsize\hangindent=1.5em\hangafter=1\noindent \textit{Membres de la} chambre \textit{haute et membres de la} chambre \textit{basse} (\textit{Angl.}); \textit{membre du sénat et membre de la chambre des députés} (\textit{France}); \textit{membre de la chambre des seigneurs et membres de la chambre des députés} (\textit{Prusse}). \lxlinks{Patres superioris et inferioris curiæ.}\par}
\par{\dictfontsize\hangindent=1.5em\hangafter=1\noindent Chancelier. \lxlinks{Præses fori regii, scriba regis supremus et sanctior.}\par}
\par{\dictfontsize\hangindent=1.5em\hangafter=1\noindent Chapelet, \textit{rosaire}. \lxlinks{Corolla Mariana; corona calculorum precatoriorum.}\par}
\par{\dictfontsize\hangindent=1.5em\hangafter=1\noindent Chapelle. \lxlinks{Ædicula sacra, sacellum, sacrarium.}\par}
\par{\dictfontsize\hangindent=1.5em\hangafter=1\noindent \textit{Maître de} chapelle. \lxlinks{Præfectus chori musicorum \textit{vel} canentium.}\par}
\par{\dictfontsize\hangindent=1.5em\hangafter=1\noindent \textit{Le} chapitre \textit{d’une cathédrale, le corps des chanoines}. \lxlinks{Senatus sacerdotum.}\par}
\par{\dictfontsize\hangindent=1.5em\hangafter=1\noindent \textit{Faire une} charge \textit{de cavalerie} (\textit{t. de guerre}). \lxlinks{Equestrem procellam excitare.}\par}
\par{\dictfontsize\hangindent=1.5em\hangafter=1\noindent Chemin \textit{couvert}. \lxlinks{Via operta.}\par}
\par{\dictfontsize\hangindent=1.5em\hangafter=1\noindent Chevalier \textit{de Malte}. \lxlinks{Eques sacram in Turcas militiam professus, eques Rhodius \textit{seu} Melitensis.}\par}
\par{\dictfontsize\hangindent=1.5em\hangafter=1\noindent Chevaux \textit{de bât, de somme, de trait}. \lxlinks{Jumenta sarcinaria, clitellaria.}\par}
\par{\dictfontsize\hangindent=1.5em\hangafter=1\noindent Chevaux \textit{de frise} (\textit{t. de guer.}), \textit{pièces de bois garnies de pieux ferrés, pour défendre une brèche, arrêter la cavalerie, etc.} \lxlinks{Ericii militares, cervi, trabes aculeis confixæ, asperatæ, ante portam, castrorum valla, aciei principia defixæ.}\par}
\par{\dictfontsize\hangindent=1.5em\hangafter=1\noindent Chevaux \textit{de frise} (\textit{t. de fortif.}). \lxlinks{Sudes ligneæ ferreis cuspidibus horrentes, armatæ.}\par}
\par{\dictfontsize\hangindent=1.5em\hangafter=1\noindent Cheveux \textit{postiches, faux cheveux, perruque, coiffure}. \lxlinks{Fictitiæ comæ, supposititiæ suggestus comæ, galericulus, caliendrum.}\par}
\par{\dictfontsize\hangindent=1.5em\hangafter=1\noindent Chiffre. \lxlinks{Numeri Arabici.}\par}
\par{\dictfontsize\hangindent=1.5em\hangafter=1\noindent \textit{Écrire en} chiffres; \textit{chiffres}. \lxlinks{Notis arithmeticis et occultis uti.}\par}
\par{\dictfontsize\hangindent=1.5em\hangafter=1\noindent \textit{Le} chœur. \lxlinks{Statio musicorum.}\par}
\par{\dictfontsize\hangindent=1.5em\hangafter=1\noindent \textit{Ligne de} circonvallation. \lxlinks{Limes extimus, circuitus castrorum.}\par}
\par{\dictfontsize\hangindent=1.5em\hangafter=1\noindent \textit{Ouvrages}, retranchements, \textit{lignes de défense en avant du camp}. \lxlinks{Pro portis castrorum.}\par}
\par{\dictfontsize\hangindent=1.5em\hangafter=1\noindent \textit{Lettre} circulaire. \lxlinks{Libelli circum ordines missi.}\par}
\par{\dictfontsize\hangindent=1.5em\hangafter=1\noindent Lancer \} une lettre circulaire. Donner \} Publier \} \lxlinks{Provincias publicis litteris de re aliqua edocere.}\par}
\par{\dictfontsize\hangindent=1.5em\hangafter=1\noindent \textit{Modeler qqn en} cire. \lxlinks{In cereis figuris fingere aliquem, cera exprimere aliquem.}\par}
\par{\dictfontsize\hangindent=1.5em\hangafter=1\noindent \textit{Traduire}, citer \textit{qqn en justice}. \lxlinks{Deferre alicujus nomen ad senatum, in jus vocare, postulare legibus, reum peragere.}\par}
\par{\dictfontsize\hangindent=1.5em\hangafter=1\noindent \textit{Prêter le serment de fidélité à qqn}; \textit{se mettre sous le patronage de qqn}; \textit{se faire le} client \textit{de qqn}. \lxlinks{In fidem et clientelam alicujus se recipere, sacramentum apud aliquem dicere, clientelam profiteri.}\par}
\par{\dictfontsize\hangindent=1.5em\hangafter=1\noindent Coadjuteur \textit{d’un Évêque}. \lxlinks{Adjutor Episcopi, in partem sacrarum curarum assumptus.}\par}
\par{\dictfontsize\hangindent=1.5em\hangafter=1\noindent Collation (\textit{en parlant des charges, des dignités}). \lxlinks{Legitima collocatio sacri fundi.}\par}
\par{\dictfontsize\hangindent=1.5em\hangafter=1\noindent Collation (\textit{léger repas au lieu du souper}). \lxlinks{Modica et vespertina cænula, pomeridiana libatio.}\par}
\par{\dictfontsize\hangindent=1.5em\hangafter=1\noindent \textit{Droit de} collation. \lxlinks{Jus vocandi sacerdotem.}\par}
\par{\dictfontsize\hangindent=1.5em\hangafter=1\noindent Colonel; \textit{tribun} (\textit{militaire}). \lxlinks{Tribunus militum. Moderator legionis.}\par}
\par{\dictfontsize\hangindent=1.5em\hangafter=1\noindent Colonel, \textit{commandant des cadets, du corps des cadets}. \lxlinks{Præfectus principibus nobilissimæ juventutis.}\par}
\par{\dictfontsize\hangindent=1.5em\hangafter=1\noindent Colonel, \textit{capitaine général, officier d’}artillerie; \textit{capitaine, officier du génie}. \lxlinks{Præfectus, tribunus cohortis tormentariæ, præfectus fabrum.}\par}
\par{\dictfontsize\hangindent=1.5em\hangafter=1\noindent \textit{Lieutenant}-colonel. \lxlinks{Legionis legatus, tribuni militum vicarius, tribuni legatus.}\par}
\par{\dictfontsize\hangindent=1.5em\hangafter=1\noindent Commandant \textit{des postes, des sentinelles de nuit}. \lxlinks{Præfectus vigiliarum classis.}\par}
\par{\dictfontsize\hangindent=1.5em\hangafter=1\noindent \textit{Major de place}; commandant \textit{de place}. \lxlinks{Præfectus vigiliarum præsidiarius.}\par}
\par{\dictfontsize\hangindent=1.5em\hangafter=1\noindent \textit{Maintenir une discipline sévère parmi les troupes, avoir le} commandement \textit{rude}. \lxlinks{Acerbiore uti imperio, disciplinæ militaris rigidum exactorem se præbere. Terrore magis minisque quam benevolentia continere militem.}\par}
\par{\dictfontsize\hangindent=1.5em\hangafter=1\noindent \textit{Conférer à un général un pouvoir dictatorial, un} commandement \textit{absolu, la dictature}. \lxlinks{Summam imperii, imperium liberum militare, nullis legibus aut conditionibus adstrictum, alicui committere.}\par}
\par{\dictfontsize\hangindent=1.5em\hangafter=1\noindent Commanderies (\textit{en parlant de l’ordre teutonique et des chevaliers de Malte}). \lxlinks{Ordinis sacri præfecturæ quæ \textit{vel} ad Teutonicos \textit{aut} Melitenses equites spectant.}\par}
\par{\dictfontsize\hangindent=1.5em\hangafter=1\noindent \textit{Tour de} comédie, \textit{représenter une comédie}. \lxlinks{Peragere fabulam.}\par}
\par{\dictfontsize\hangindent=1.5em\hangafter=1\noindent Comédiens, \textit{histrions, saltimbanques}. \lxlinks{Actores fabularum, artifices scenici, histriones.}\par}
\par{\dictfontsize\hangindent=1.5em\hangafter=1\noindent Commissaire. \lxlinks{Recuperator principis, Regis a negotiis commissus, cognitor causarum, rei inquisitor, quæsitor.}\par}
\par{\dictfontsize\hangindent=1.5em\hangafter=1\noindent Commission, \textit{exécution} (\textit{d’un mandat}). \lxlinks{Procuratio rei mandatæ.}\par}
\par{\dictfontsize\hangindent=1.5em\hangafter=1\noindent Compagnie (\textit{des soldats}). \lxlinks{Centuria.}\par}
\par{\dictfontsize\hangindent=1.5em\hangafter=1\noindent \textit{Être de faible} complexion, \textit{avoir une santé délicate}. \lxlinks{Tenui aut nulla potius valetudine esse. Corpus morbis obnoxium.}\par}
\par{\dictfontsize\hangindent=1.5em\hangafter=1\noindent Complices. \lxlinks{Sceleris socii, facinoris administri.}\par}
\par{\dictfontsize\hangindent=1.5em\hangafter=1\noindent Compliment, \textit{révérence}. \lxlinks{Honor verborum, honorifica compellatio; verba officiis plena.}\par}
\par{\dictfontsize\hangindent=1.5em\hangafter=1\noindent Sans compliment, \textit{sans déguisement, loyalement}. \lxlinks{Sine fuco et fallaciis, citra honorem verborum.}\par}
\par{\dictfontsize\hangindent=1.5em\hangafter=1\noindent Compte \textit{borgne, inexact}. \lxlinks{Impeditæ.}\par}
\par{\dictfontsize\hangindent=1.5em\hangafter=1\noindent Compte \textit{juste, exact}. \lxlinks{Expeditæ.}\par}
\par{\dictfontsize\hangindent=1.5em\hangafter=1\noindent Compte \textit{des recettes et des défenses}. \lxlinks{Ratio accepti et expensi, ratio acceptorum et datorum.}\par}
\par{\dictfontsize\hangindent=1.5em\hangafter=1\noindent \textit{Régler un} compte; \textit{liquider}, (\textit{terme de finances}). \lxlinks{Rationes subducere, rationes dispungere, edere.}\par}
\par{\dictfontsize\hangindent=1.5em\hangafter=1\noindent \textit{Articles de} compte. \lxlinks{Æra.}\par}
\par{\dictfontsize\hangindent=1.5em\hangafter=1\noindent \textit{Mettre qqchose sur le} compte \textit{de qqn}; \textit{faire passer qqche en compte à qqn}. \lxlinks{Expensum alicui ferre.}\par}
\par{\dictfontsize\hangindent=1.5em\hangafter=1\noindent \textit{Résumé d’un} compte; \textit{compte par bref état}. \lxlinks{Breviarium rationum.}\par}
\par{\dictfontsize\hangindent=1.5em\hangafter=1\noindent Comtes \textit{qui sont princes du Saint-Empire}. \lxlinks{Comites Principum dignitate præfulgentes.}\par}
\par{\dictfontsize\hangindent=1.5em\hangafter=1\noindent Condamné. \lxlinks{Datus ad supplicium.}\par}
\par{\dictfontsize\hangindent=1.5em\hangafter=1\noindent Se confesser, \textit{aller à confesse}. \lxlinks{Animum noxis expiare, obire confessionis sacrum mysterium.}\par}
\par{\dictfontsize\hangindent=1.5em\hangafter=1\noindent Confesseur, \textit{directeur de conscience}; \textit{directeur spirituel}. \lxlinks{Qui alicui a confessionibus est, qui alicujus animum moderatur.}\par}
\par{\dictfontsize\hangindent=1.5em\hangafter=1\noindent \textit{Saisir qqche}; \textit{faire la saisie de qqche}; \textit{mettre arrêt sur qqche}; \textit{séquestrer}; \textit{mettre sous séquestre}; confisquer. \lxlinks{Manu prætoria alicujus bona detinere.}\par}
\par{\dictfontsize\hangindent=1.5em\hangafter=1\noindent Confisquer \textit{les biens}. \lxlinks{Facultes alicujus in publicum conferre, bona alicujus publicare, in publicum addicere, proscribere.}\par}
\par{\dictfontsize\hangindent=1.5em\hangafter=1\noindent \textit{Obtenir son} congé (\textit{du service militaire}). \lxlinks{Vacatione militiæ, missione donari.}\par}
\par{\dictfontsize\hangindent=1.5em\hangafter=1\noindent Connétable (\textit{chef des armées}). \lxlinks{Comes stabuli, summus rei bellicæ in Gallia præfectus, tribunus celerum.}\par}
\par{\dictfontsize\hangindent=1.5em\hangafter=1\noindent \textit{Mauvaise} conscience; \textit{conscience perverse}. \textit{Une conscience bourrelée de remords.} \lxlinks{Conscientia sceleris, facinoris, mens male sibi conscia, quæ agitatur et perterretur tædis furiarum ardentibus.}\par}
\par{\dictfontsize\hangindent=1.5em\hangafter=1\noindent \textit{Chercher, trouver, mettre sa consolation dans une bonne} conscience. \lxlinks{Conscientia bene factorum se solari.}\par}
\par{\dictfontsize\hangindent=1.5em\hangafter=1\noindent Être empêché de faire \} une chose de conscience. ne pouvoir pas faire \} ( Quand la conscience ne permet pas de faire une chose.) \lxlinks{Religione impediri.}\par}
\par{\dictfontsize\hangindent=1.5em\hangafter=1\noindent \textit{Résoudre des cas de} conscience; \textit{traiter des devoirs du chrétien}. \lxlinks{Quæstiones de virtutibus, moribus explicare, Officia Christiani hominis tradere, quæ sunt de finibus bonorum et malorum in re Christiana explicare.}\par}
\par{\dictfontsize\hangindent=1.5em\hangafter=1\noindent \textit{N’avoir point de} conscience, \textit{être un homme sans conscience, sans foi ni loi}. \lxlinks{Religione nihil sibi ducere. Homo sine ulla religione et fide.}\par}
\par{\dictfontsize\hangindent=1.5em\hangafter=1\noindent \textit{Je puis faire cela en sûreté de conscience, ou en bonne} conscience, \textit{ou en toute conscience}. \lxlinks{Salvo officio, salva fide id mihi integrum est.}\par}
\par{\dictfontsize\hangindent=1.5em\hangafter=1\noindent \textit{Membre du conseil ecclésiastique, membre du} conseil de fabrique. \lxlinks{Assessor sacri senatus.}\par}
\par{\dictfontsize\hangindent=1.5em\hangafter=1\noindent \textit{Conseil ecclésiastique}, conseil de fabrique (\textit{cathol.}), \textit{consistoire} (\textit{protest.}). \lxlinks{Senatus sacer consessus, \textit{vel} consilium sancti ordinis.}\par}
\par{\dictfontsize\hangindent=1.5em\hangafter=1\noindent Conseiller \textit{d’État; membre du conseil d’État}; \textit{membre du conseil privé}. \lxlinks{Qui principi a consiliis est in rebus quas publicas appellamus. Vir e sacerrimo principis consilio.}\par}
\par{\dictfontsize\hangindent=1.5em\hangafter=1\noindent Conseiller \textit{intime}. \lxlinks{Consiliorum regiorum particeps, qui est in sanctiore principis consilio.}\par}
\par{\dictfontsize\hangindent=1.5em\hangafter=1\noindent \textit{Membre du} consistoire. \lxlinks{Assessor sacri senatus.}\par}
\par{\dictfontsize\hangindent=1.5em\hangafter=1\noindent \textit{Décret, décision du} consistoire. \lxlinks{Decretum senatus sacri.}\par}
\par{\dictfontsize\hangindent=1.5em\hangafter=1\noindent Constitution (\textit{ensemble des lois fondamentales d’un état}). \lxlinks{Leges imperii quibus salus reipublicæ continetur, nititur.}\par}
\par{\dictfontsize\hangindent=1.5em\hangafter=1\noindent Constitution \textit{robuste}; \textit{forte constitution}; \textit{santé solide}. \lxlinks{Corpus natura solidum in multos labores duratum. Corporis firmitas.}\par}
\par{\dictfontsize\hangindent=1.5em\hangafter=1\noindent \textit{État du ciel}; \textit{aspect des étoiles} (\textit{en termes d’astrologie}), constellation. \lxlinks{Astrorum cœli affectio, siderum positus.}\par}
\par{\dictfontsize\hangindent=1.5em\hangafter=1\noindent Constitution \textit{militaire}. \lxlinks{Belli apparatus.}\par}
\par{\dictfontsize\hangindent=1.5em\hangafter=1\noindent \textit{Écrire l’histoire} contemporaine. \lxlinks{Scriptis res suorum temporum mandare, celebrare, illustrare.}\par}
\par{\dictfontsize\hangindent=1.5em\hangafter=1\noindent Contexte. \lxlinks{Nexus orationis et forma, verborum ambitus, connexa et cohærens inter se orationis series.}\par}
\par{\dictfontsize\hangindent=1.5em\hangafter=1\noindent Contrebande, \textit{marchandise prohibée}. \lxlinks{Merces prohibitæ, interdictæ.}\par}
\par{\dictfontsize\hangindent=1.5em\hangafter=1\noindent Contre-batterie. \lxlinks{Suggesta tormentorum hostibus opposita.}\par}
\par{\dictfontsize\hangindent=1.5em\hangafter=1\noindent \textit{Inspecteur}; contre-maître, \textit{surveillant, intendant}. \lxlinks{Redemptor operis, manceps operarum.}\par}
\par{\dictfontsize\hangindent=1.5em\hangafter=1\noindent Contributions, \textit{impôts}; \textit{charge}. \lxlinks{Census bonorum, tributum.}\par}
\par{\dictfontsize\hangindent=1.5em\hangafter=1\noindent Contrôleur. \lxlinks{Rationum custos et vindex, redituum consignator.}\par}
\par{\dictfontsize\hangindent=1.5em\hangafter=1\noindent Controverses, \textit{points controverses}. \lxlinks{Fidei capita, quæ in disceptationem vocantur.}\par}
\par{\dictfontsize\hangindent=1.5em\hangafter=1\noindent Convalescence. \lxlinks{Pristina valetudo, confirmata a morbo valetudo.}\par}
\par{\dictfontsize\hangindent=1.5em\hangafter=1\noindent Convention, \textit{traité pour} l’échange \textit{des transfuges et des prisonniers}. \lxlinks{Pactio de transfugis et captivis permutandis.}\par}
\par{\dictfontsize\hangindent=1.5em\hangafter=1\noindent Convoquer \textit{la diète, la chambre}. \lxlinks{Comitia edere, indicere.}\par}
\par{\dictfontsize\hangindent=1.5em\hangafter=1\noindent Copie \textit{de lettre}. \lxlinks{Descriptio, exemplum litterarum.}\par}
\par{\dictfontsize\hangindent=1.5em\hangafter=1\noindent \textit{Un exemplaire, une copie figurée}; \textit{une} copie \textit{conforme}. \lxlinks{Exemplum recognitum.}\par}
\par{\dictfontsize\hangindent=1.5em\hangafter=1\noindent Copiste. \lxlinks{Homo, puer a manu, librarius.}\par}
\par{\dictfontsize\hangindent=1.5em\hangafter=1\noindent \textit{Ceux qui sont du même corps de métier, de la même} Corporation. \lxlinks{Curiales.}\par}
\par{\dictfontsize\hangindent=1.5em\hangafter=1\noindent Corps, \textit{armée de réserve}. \lxlinks{Manus militum ad subitos belli casus seposita.}\par}
\par{\dictfontsize\hangindent=1.5em\hangafter=1\noindent Corsaire (\textit{t. de mar.}). \lxlinks{Navis prædatoria.}\par}
\par{\dictfontsize\hangindent=1.5em\hangafter=1\noindent Corvette; \textit{yacht}. \lxlinks{Navis liburnica.}\par}
\par{\dictfontsize\hangindent=1.5em\hangafter=1\noindent Cotisation \textit{volontaire}; \textit{secours volontaire en argent}; \textit{subside, subvention volontaire}. \lxlinks{Collatio gratuita.}\par}
\par{\dictfontsize\hangindent=1.5em\hangafter=1\noindent Couleuvrine. \lxlinks{Tormentum colubrinum, colubrina.}\par}
\par{\dictfontsize\hangindent=1.5em\hangafter=1\noindent Couleuvrine (\textit{artill.}). \lxlinks{Tormenta longiora a colubris nominata.}\par}
\par{\dictfontsize\hangindent=1.5em\hangafter=1\noindent \textit{On fait jouer les} couleuvrines; \textit{on a tiré des couleuvrines}. \lxlinks{Colubrinis verberatæ sunt ædes, hostis colubrinas exercuit, misit in tecta.}\par}
\par{\dictfontsize\hangindent=1.5em\hangafter=1\noindent Couper \textit{la retraite à qqn}; couper \textit{les vivres à une ville, etc.}...; \textit{empêcher le transport des vivres,—des provisions} (\textit{pour une forteresse, une armée}). \lxlinks{Reditu aliquem excludere. Commeatu privare, intercludere, prohibere.}\par}
\par{\dictfontsize\hangindent=1.5em\hangafter=1\noindent \textit{Par} coups, \textit{de taille}. \lxlinks{Cæsim.}\par}
\par{\dictfontsize\hangindent=1.5em\hangafter=1\noindent Cour; \textit{suite d’un prince, d’un roi}; \textit{courtisans}. \lxlinks{Comitatus aulicus.}\par}
\par{\dictfontsize\hangindent=1.5em\hangafter=1\noindent \textit{Juge de la} cour. \lxlinks{Judex curiæ provincialis.}\par}
\par{\dictfontsize\hangindent=1.5em\hangafter=1\noindent Cour \textit{d’Appel}. \lxlinks{Summus senatus provocationum.}\par}
\par{\dictfontsize\hangindent=1.5em\hangafter=1\noindent \textit{Employé à la chambre, à la} cour des comptes; \textit{directeur, président de la chambre, de la cour des comptes}; \textit{trésorier, administrateur de la caisse du roi, d’un prince}. \lxlinks{Curiæ fisci adscriptus.}\par}
\par{\dictfontsize\hangindent=1.5em\hangafter=1\noindent \textit{Secrétaire de la chambre des domaines, de la} cour des comptes. \lxlinks{Census curator.}\par}
\par{\dictfontsize\hangindent=1.5em\hangafter=1\noindent \textit{Rendre ses comptes à la} cour. \lxlinks{Ad ærarium rationes referre.}\par}
\par{\dictfontsize\hangindent=1.5em\hangafter=1\noindent \textit{Marcher}, courir \textit{en grande hâte, précipitamment, à toutes jambes}; \textit{hâter son pas}. \lxlinks{Cursum incitare, iter accelerare, celeri itinere aliquo contendere, non observata, stata quiete festinare iter.}\par}
\par{\dictfontsize\hangindent=1.5em\hangafter=1\noindent Couronne \textit{foudroyante ou d’artifice} (\textit{guer.}). \lxlinks{Malleoli et ignes, tædæ pice et sulphure mistæ, picea sulphureaque serta.}\par}
\par{\dictfontsize\hangindent=1.5em\hangafter=1\noindent Courtine (\textit{mur entre deux bastions et qui en joint les flancs}). \lxlinks{Injecti muri loricæ, longitudo \textit{seu} ductus murorum lorica mœnium.}\par}
\par{\dictfontsize\hangindent=1.5em\hangafter=1\noindent \textit{Vils} courtisans. \lxlinks{Tineæ soricesque palatii.}\par}
\par{\dictfontsize\hangindent=1.5em\hangafter=1\noindent Couvrir d’une flotte, \textit{de vaisseaux, les côtes maritimes}. \lxlinks{Ora maritima classem disponere.}\par}
\par{\dictfontsize\hangindent=1.5em\hangafter=1\noindent \textit{Tout} crédit \textit{lui est enlevé}. \lxlinks{Omnium rerum fides ei abrogatur.}\par}
\par{\dictfontsize\hangindent=1.5em\hangafter=1\noindent On lui enlève \} tout crédit, Il perd \} Avoir perdu son ( tout) crédit. \lxlinks{Fidem amittere.}\par}
\par{\dictfontsize\hangindent=1.5em\hangafter=1\noindent \textit{Être en} crédit \textit{auprès de qqn}; \textit{avoir bonne réputation}.\textit{Avoir perdu son} (\textit{tout}) crédit. \lxlinks{Fidem amittere.}\par}
\par{\dictfontsize\hangindent=1.5em\hangafter=1\noindent \textit{Il n’y a plus de} crédit; \textit{il n’y a pas de confiance}. \lxlinks{Concidit fides, fides cessit e foro.}\par}
\par{\dictfontsize\hangindent=1.5em\hangafter=1\noindent Croiser (\textit{sur la mer}), \textit{établir des croisières} (\textit{t. de marine}), \textit{courir} (\textit{les côtes}). \lxlinks{Maritimos prædones consectando mare tutum reddere.}\par}
\par{\dictfontsize\hangindent=1.5em\hangafter=1\noindent \textit{Courir la mer}, croiser \textit{sur la mer en tout sens}; \textit{couvrir la mer de ses croisières}. \lxlinks{Classem dimittere quoquoversus. Maria, sinus omnes scrutari navibus.}\par}
\par{\dictfontsize\hangindent=1.5em\hangafter=1\noindent \textit{Protéger une province, un pays contre les ravages de la guerre}; couvrir une province. \lxlinks{Defendere provinciam a depopulatione.}\par}
\par{\dictfontsize\hangindent=1.5em\hangafter=1\noindent Couvrir \textit{les derrières de l’armée par une forte arrière-garde}. \lxlinks{Cingere ultimum agmen valida manu.}\par}
\par{\dictfontsize\hangindent=1.5em\hangafter=1\noindent \textit{La Sainte} Croix. \lxlinks{Lignum salutare, in quo omnium noxas sua unius morte expiavit Christus.}\par}
\par{\dictfontsize\hangindent=1.5em\hangafter=1\noindent \textit{Faire le signe de la} Croix, \textit{se signer}. \lxlinks{Salutari se signo munire, signum crucis digito in fronte effingere, ducere.}\par}
\par{\dictfontsize\hangindent=1.5em\hangafter=1\noindent Cuirassier; \textit{en général tous les soldats qui portent la cuirasse, tels que les carabiniers}. \lxlinks{Milites cataphracti, equites loricati quibus tegumenta ferreis laminis serie inter se connexis.}\par}
\par{\dictfontsize\hangindent=0pt\noindent\textls[120]{\scshape d}\par}
\par{\dictfontsize\hangindent=1.5em\hangafter=1\noindent Date \textit{d’une lettre}. \lxlinks{Dies epistolæ adscripta.}\par}
\par{\dictfontsize\hangindent=1.5em\hangafter=1\noindent \textit{Lettre} datée \textit{du même jour}; \textit{lettre de même date}. \lxlinks{Litteræ eodem datæ tempore.}\par}
\par{\dictfontsize\hangindent=1.5em\hangafter=1\noindent Davier, \textit{tenailles, pince} (\textit{instr. de chirurgie}). \lxlinks{Forcipes.}\par}
\par{\dictfontsize\hangindent=1.5em\hangafter=1\noindent Débarquer \textit{les soldats}. \lxlinks{Milites in terram exponere.}\par}
\par{\dictfontsize\hangindent=1.5em\hangafter=1\noindent \textit{Dissiper ses biens en} débauches. \lxlinks{Luxu opes perdere, prodige profundere.}\par}
\par{\dictfontsize\hangindent=1.5em\hangafter=1\noindent Décret, \textit{ordre du Sacré-Collége}. \lxlinks{Decretum sacri senatus.}\par}
\par{\dictfontsize\hangindent=1.5em\hangafter=1\noindent Défalquer;—\textit{soustraire}. \lxlinks{Detrahere de summa.}\par}
\par{\dictfontsize\hangindent=1.5em\hangafter=1\noindent \textit{Craindre les} défilés. \lxlinks{Angustias itineris vereri.}\par}
\par{\dictfontsize\hangindent=1.5em\hangafter=1\noindent \textit{Occuper les} défilés. \lxlinks{Viarum, montium fauces milite insidere, præsidia in angustiis collocare, arctas vias milite et operibus munire.}\par}
\par{\dictfontsize\hangindent=1.5em\hangafter=1\noindent \textit{Enfermer l’armée ennemie dans un} défilé, \textit{dans une gorge de montagne}; \textit{acculer l’armée ennemie contre une montagne, un défilé}. \lxlinks{Adversarios locorum angustiis constringere.}\par}
\par{\dictfontsize\hangindent=1.5em\hangafter=1\noindent Définir \textit{un mot}; \textit{en fixer le sens, l’acception}. \lxlinks{Notionem verbi certis limitibus circumscribere, definire, arctare.}\par}
\par{\dictfontsize\hangindent=1.5em\hangafter=1\noindent Dégrader; \textit{casser}. \lxlinks{Exauctorare, sacro gradu dejicere.}\par}
\par{\dictfontsize\hangindent=1.5em\hangafter=1\noindent Dégréer \textit{un vaisseau} (\textit{terme de marine}). \lxlinks{Exarmare navem.}\par}
\par{\dictfontsize\hangindent=1.5em\hangafter=1\noindent Démocratie. \lxlinks{Imperium populi, \textit{seu} populare.}\par}
\par{\dictfontsize\hangindent=1.5em\hangafter=1\noindent Départ (\textit{des troupes}); \textit{décampement}. \lxlinks{Discessus militum; profectio; militare iter.}\par}
\par{\dictfontsize\hangindent=1.5em\hangafter=1\noindent \textit{Annoncer le} départ. \lxlinks{Iter pronuntiare.}\par}
\par{\dictfontsize\hangindent=1.5em\hangafter=1\noindent \textit{Donner les ordres du} départ. \lxlinks{Vasa militari more conclamare; signum dare colligendi vasa.}\par}
\par{\dictfontsize\hangindent=1.5em\hangafter=1\noindent Dépendances. \lxlinks{Jura accessionum.}\par}
\par{\dictfontsize\hangindent=1.5em\hangafter=1\noindent Députés \textit{d’un pays}. \lxlinks{Status provincialis, collationum provincialium præfecti.}\par}
\par{\dictfontsize\hangindent=1.5em\hangafter=1\noindent Desserrer \textit{les rangs}. \lxlinks{Manipulos laxare, diducere.}\par}
\par{\dictfontsize\hangindent=1.5em\hangafter=1\noindent Dessert; \textit{friandises, confitures, pâtisseries, tout ce qu’on sert à table au dessert}. \lxlinks{Mensa secunda, bellaria.}\par}
\par{\dictfontsize\hangindent=1.5em\hangafter=1\noindent Détachement \textit{de soldats}. \lxlinks{Manus militum, globus armatorum, delecti milites.}\par}
\par{\dictfontsize\hangindent=1.5em\hangafter=1\noindent Diamètre; \textit{ligne diamétrale}. \lxlinks{Linea dimetiens, recta, ducta ab extremitate una circuli ad alteram per centrum.}\par}
\par{\dictfontsize\hangindent=1.5em\hangafter=1\noindent Action ou \} en diffamation. Procès \} \lxlinks{Actio injuriarum.}\par}
\par{\dictfontsize\hangindent=1.5em\hangafter=1\noindent \textit{Avoir un procès en} diffamation. \lxlinks{De fama dimicare.}\par}
\par{\dictfontsize\hangindent=1.5em\hangafter=1\noindent \textit{Il est d’humeur} difficile, \textit{d’un caractère difficile, d’une grande inégalité d’humeur}. \lxlinks{Inest dissimilitudo morum.}\par}
\par{\dictfontsize\hangindent=1.5em\hangafter=1\noindent Digue; \textit{jetée}; \textit{chaussée}. \lxlinks{Moles fluctibus opposita.}\par}
\par{\dictfontsize\hangindent=1.5em\hangafter=1\noindent \textit{Arrêter le débordement d’un fleuve, d’un torrent, par une} digue; endiguer \textit{un fleuve, un torrent}. \lxlinks{Fluvium extra ripas effusum coercere.}\par}
\par{\dictfontsize\hangindent=1.5em\hangafter=1\noindent \textit{Articles de} discipline \textit{militaire}. \lxlinks{Leges militariæ, castrenses.}\par}
\par{\dictfontsize\hangindent=1.5em\hangafter=1\noindent \textit{Il parle tantôt d’une chose, tantôt d’une autre, sans suite}; \textit{ses} discours \textit{n’ont ni queue ni tête}; \textit{il fait des coq-à-l’âne}; \textit{il dit des extravagances}. \lxlinks{Sermo sibi non constat.}\par}
\par{\dictfontsize\hangindent=1.5em\hangafter=1\noindent \textit{Se rendre à} discrétion (\textit{se dit d’une armée ou d’une ville prise}). \lxlinks{Armis depositis ad victoris clementiam confugere, se suaque omnia alterius potestati permittere, nullis conditionibus admissis urbem \textit{vel} provinciam sub suam potestatem redigere.}\par}
\par{\dictfontsize\hangindent=1.5em\hangafter=1\noindent Dispenser, \textit{exempter, accorder la dispense, l’exemption d’une loi}. \lxlinks{Legibus solvere, legis gratiam facere.}\par}
\par{\dictfontsize\hangindent=1.5em\hangafter=1\noindent \textit{Nous sommes à une lieue, à la} distance \textit{d’une lieue l’un de l’autre}. \lxlinks{Uno lapide distamus.}\par}
\par{\dictfontsize\hangindent=1.5em\hangafter=1\noindent \textit{Être en délire}, divaguer, \textit{battre la campagne}. \lxlinks{Aliena loqui.}\par}
\par{\dictfontsize\hangindent=1.5em\hangafter=1\noindent Douane \textit{ou Accise, impôt sur les denrées, sur la consommation}. \lxlinks{Portorium, vectigal eorum quæ consumuntur, tributum rerum ad victum et amictum pertinentium.}\par}
\par{\dictfontsize\hangindent=1.5em\hangafter=1\noindent \textit{Dragons.} \lxlinks{Dimachi. Desultorii milites.}\par}
\par{\dictfontsize\hangindent=1.5em\hangafter=1\noindent Drapeaux, \textit{étendards, enseignes, bannières}. \lxlinks{Signa militaria.}\par}
\par{\dictfontsize\hangindent=1.5em\hangafter=1\noindent Déserter, \textit{abandonner le} drapeau, \textit{ou simplement déserter}. \lxlinks{A signis discedere, signa relinquere.}\par}
\par{\dictfontsize\hangindent=1.5em\hangafter=1\noindent \textit{Disposer les} drapeaux. \lxlinks{Signa statuere.}\par}
\par{\dictfontsize\hangindent=1.5em\hangafter=1\noindent \textit{Prêter serment sous le} drapeau. \lxlinks{Sacramento se illis obstringere. De signis tutandis sequendisque perpetuo sacramentum sacerrimum dicere.}\par}
\par{\dictfontsize\hangindent=1.5em\hangafter=1\noindent Droit \textit{de nommer au bénéfice}. \lxlinks{Jus conferendi sacram dignitatem.}\par}
\par{\dictfontsize\hangindent=1.5em\hangafter=1\noindent \textit{Se battre en} duel. \lxlinks{Ferro cum aliquo decernere.}\par}
\par{\dictfontsize\hangindent=0pt\noindent\textls[120]{\scshape e}\par}
\par{\dictfontsize\hangindent=1.5em\hangafter=1\noindent \textit{Souscrire à l’}échange \textit{des prisonniers}; \textit{traiter de l’échange des prisonniers}. \lxlinks{Captivorum commutationem recipere.}\par}
\par{\dictfontsize\hangindent=1.5em\hangafter=1\noindent Écho. \lxlinks{Vox repercussa, sonus repercussus. Imago verbi.}\par}
\par{\dictfontsize\hangindent=1.5em\hangafter=1\noindent \textit{Faire} écho. \lxlinks{Vocem multiplicato sono reddere.}\par}
\par{\dictfontsize\hangindent=1.5em\hangafter=1\noindent \textit{Les montagnes font} écho. \lxlinks{Saxa et solitudines voci respondent.}\par}
\par{\dictfontsize\hangindent=1.5em\hangafter=1\noindent \textit{Coureur} (\textit{guer.}); \textit{batteur d’estrade} (\textit{guer.}); \textit{guérilla}, éclaireur (\textit{guer.}). \lxlinks{Excursor, procursator.}\par}
\par{\dictfontsize\hangindent=1.5em\hangafter=1\noindent \textit{Ouvrir une} école; \textit{ouvrir un cours} (\textit{de littérature}). \lxlinks{Ludum litterarium aperire.}\par}
\par{\dictfontsize\hangindent=1.5em\hangafter=1\noindent \textit{L’}écriture \textit{sainte}; \textit{les saintes écritures}. \lxlinks{Sacræ litteræ; sacer codex.}\par}
\par{\dictfontsize\hangindent=1.5em\hangafter=1\noindent \textit{Registre} d’écrou. \lxlinks{Commentarii custodiarum.}\par}
\par{\dictfontsize\hangindent=1.5em\hangafter=1\noindent \textit{Princes} électeurs; \textit{princes de l’empire de l’Allemagne qui avaient droit de vote dans l’élection de l’empereur}. \lxlinks{Principes Imperii qui suffragio suo imperatores creant.}\par}
\par{\dictfontsize\hangindent=1.5em\hangafter=1\noindent \textit{Jour des} élections; \textit{jour de l’élection}. \lxlinks{Dies comitialis.}\par}
\par{\dictfontsize\hangindent=1.5em\hangafter=1\noindent \textit{Prince} électoral. \lxlinks{Heres Electoratus, in spem Imperii genitus, princeps juventutis.}\par}
\par{\dictfontsize\hangindent=1.5em\hangafter=1\noindent Élite \textit{de l’armée}; \textit{troupes} levées \textit{pour défendre une province}. \lxlinks{Delecta manus, delecti milites.}\par}
\par{\dictfontsize\hangindent=1.5em\hangafter=1\noindent Emprunter \textit{de l’argent pour payer ses dettes}. \lxlinks{Versuram facere.}\par}
\par{\dictfontsize\hangindent=1.5em\hangafter=1\noindent \textit{Mettre, vendre qqche à l’enchère, aux enchères};—\textit{vendre à} l’encan. \lxlinks{Auctionem constituere.}\par}
\par{\dictfontsize\hangindent=1.5em\hangafter=1\noindent \textit{Vendre à l’}encan; \textit{mettre à l’enchère}. \lxlinks{Bona voci præconis subjicere, hastæ subjicere.}\par}
\par{\dictfontsize\hangindent=1.5em\hangafter=1\noindent Encre. \lxlinks{Atramentum scriptorium, librarium.}\par}
\par{\dictfontsize\hangindent=1.5em\hangafter=1\noindent Encre \textit{d’imprimerie}. \lxlinks{Atramentum fuligineum.}\par}
\par{\dictfontsize\hangindent=1.5em\hangafter=1\noindent \textit{Arrêter le débordement d’un fleuve, d’un torrent, par une} digue; endiguer \textit{un fleuve, un torrent}. \lxlinks{Fluvium extra ripas effusum coercere.}\par}
\par{\dictfontsize\hangindent=1.5em\hangafter=1\noindent Enfoncer \textit{le flanc de l’armée ennemie}; \textit{prendre l’ennemi en flanc}. \lxlinks{In latus hostium invehi.}\par}
\par{\dictfontsize\hangindent=1.5em\hangafter=1\noindent Enfoncer \textit{les escadrons, les rangs de la cavalerie}. \lxlinks{Se inter equitum turmas inferre.}\par}
\par{\dictfontsize\hangindent=1.5em\hangafter=1\noindent Enfoncer \textit{le front de l’armée ennemie}. \lxlinks{Principia aciei disjicere.}\par}
\par{\dictfontsize\hangindent=1.5em\hangafter=1\noindent \textit{Tenir parole, remplir sa promesse, remplir un} engagement, \textit{dégager sa parole}. \lxlinks{Liberare fidem.}\par}
\par{\dictfontsize\hangindent=1.5em\hangafter=1\noindent \textit{Celui qui a pris des} engagements, \textit{doit y être fidèle, doit les remplir}; \textit{celui qui a reçu un fidei-commis, doit y être fidèle, doit en être l’exécuteur fidèle}. \lxlinks{Qui fiduciam accepit, debet præstare et fidem.}\par}
\par{\dictfontsize\hangindent=1.5em\hangafter=1\noindent \textit{Serment militaire}; \textit{enrôlement}, engagement. \lxlinks{Sacramentum militare.}\par}
\par{\dictfontsize\hangindent=1.5em\hangafter=1\noindent \textit{S’}engager \textit{par écrit}. \lxlinks{Scripto fidem obstringere.}\par}
\par{\dictfontsize\hangindent=1.5em\hangafter=1\noindent \textit{Faire diversion, porter la guerre, le combat sur un point éloigné}; \textit{attaquer l’}ennemi \textit{sur un point afin de le surprendre sur un autre}; \textit{diviser ses forces par une attaque imprévue, simulée}. \lxlinks{Distinere hostem, alio hostis curas avocare, consiliis hostis occurrere.}\par}
\par{\dictfontsize\hangindent=1.5em\hangafter=1\noindent \textit{Charger l’}ennemi, \textit{fondre sur l’ennemi, les escadrons serrés}. \lxlinks{Confertis turmis in hostem impetum facere.}\par}
\par{\dictfontsize\hangindent=1.5em\hangafter=1\noindent \textit{Couper la retraite et les vivres à l’}ennemi. \lxlinks{Receptu commeatuque hostem excludere.}\par}
\par{\dictfontsize\hangindent=1.5em\hangafter=1\noindent \textit{S’élancer, fondre, se précipiter avec impétuosité sur l’}ennemi. \lxlinks{In hostem se concitare.}\par}
\par{\dictfontsize\hangindent=1.5em\hangafter=1\noindent \textit{Offrir le bataille à l’}ennemi; \textit{marcher sur l’ennemi}. \lxlinks{Hosti pugnæ copiam offerre; ad certamen hostem extrahere.}\par}
\par{\dictfontsize\hangindent=1.5em\hangafter=1\noindent \textit{Entourer, enfermer l’}ennemi \textit{d’une ligne de circonvallation}. \lxlinks{Vallo, continuo castrorum aggere, sepimento crebrisque castellis hostem cingere.}\par}
\par{\dictfontsize\hangindent=1.5em\hangafter=1\noindent \textit{Être de connivence avec l’}ennemi. \lxlinks{Litteris cum hoste collatisve consiliis ei belli opportunitates prodere. Clandestina cum hostibus consilia conferre.}\par}
\par{\dictfontsize\hangindent=1.5em\hangafter=1\noindent \textit{Se frayer un passage à travers les rangs} ennemis, \textit{l’armée ennemie}. \lxlinks{Per medios hostes perrumpere, corporibus viam aperire.}\par}
\par{\dictfontsize\hangindent=1.5em\hangafter=1\noindent \textit{Inscrire},—enregistrer; \textit{porter sur le livre, sur les registres}; \textit{insérer dans un acte}. \lxlinks{Referre in tabulas, in commentarium, in codicem.}\par}
\par{\dictfontsize\hangindent=1.5em\hangafter=1\noindent Enrôler, \textit{recruter des soldats}. \lxlinks{Copias comparare, scribere.}\par}
\par{\dictfontsize\hangindent=1.5em\hangafter=1\noindent Entonner. \lxlinks{Voce præire.}\par}
\par{\dictfontsize\hangindent=1.5em\hangafter=1\noindent Entonner; \textit{entonner le chant, la prière}. \lxlinks{Voce præire, preces præire, sacros modulos.}\par}
\par{\dictfontsize\hangindent=1.5em\hangafter=1\noindent Épitaphe; \textit{inscription tumulaire}. \lxlinks{Postumæ laudes marmori \textit{vel} æri insculptæ.}\par}
\par{\dictfontsize\hangindent=1.5em\hangafter=1\noindent \textit{Maintenir l’}équilibre \textit{entre les différents États}. \lxlinks{Æqualitatem potentiæ principes inter tueri, potentiam ita moderari ut ne quis principum reliquos viribus opprimat.}\par}
\par{\dictfontsize\hangindent=1.5em\hangafter=1\noindent Ermitage. \lxlinks{Secessus amœnus, ad vitandam populi frequentiam delectus \textit{vel} positus.}\par}
\par{\dictfontsize\hangindent=1.5em\hangafter=1\noindent \textit{Fondre sur l’ennemi en rangs serrés, en} escadrons \textit{serrés}. \lxlinks{Consertis turmis in hostem impetum dare.}\par}
\par{\dictfontsize\hangindent=1.5em\hangafter=1\noindent \textit{Escarmoucher}; \textit{faire une} escarmouche, \textit{engager une échauffourée, un combat de tirailleurs}. \lxlinks{Prœlium minutum, velitarem pugnam conserere.}\par}
\par{\dictfontsize\hangindent=1.5em\hangafter=1\noindent Escarpe. \lxlinks{Murus fossæ interior.}\par}
\par{\dictfontsize\hangindent=1.5em\hangafter=1\noindent \textit{Contre}-escarpe. \lxlinks{Lorica, acclive opus, fossa, viæ coopertæ, fossæ pars mœnibus adversæ.}\par}
\par{\dictfontsize\hangindent=1.5em\hangafter=1\noindent Escorte; \textit{garde}; \textit{convoi}. \lxlinks{Manus præsidiaria.}\par}
\par{\dictfontsize\hangindent=1.5em\hangafter=1\noindent Escouade. \lxlinks{Manipulus, decuria, contubernium.}\par}
\par{\dictfontsize\hangindent=1.5em\hangafter=1\noindent Espion; \textit{émissaire, mouchard, mouche}. \lxlinks{Speculator, explorator, emissarius.}\par}
\par{\dictfontsize\hangindent=1.5em\hangafter=1\noindent \textit{Vers l’}est, \textit{vers le levant, à l’est, au levant}. \lxlinks{Ad orientem solem.}\par}
\par{\dictfontsize\hangindent=1.5em\hangafter=1\noindent \textit{Être en grande} estime \textit{auprès de qqn, être bien vu de qqn, être estimé de quelqu’un}. \lxlinks{Esse apud aliquem bono loco.}\par}
\par{\dictfontsize\hangindent=1.5em\hangafter=1\noindent Établir \textit{de nouvelles lois somptuaires}. \lxlinks{Novas leges de vestitus ratione temperanda condere.}\par}
\par{\dictfontsize\hangindent=1.5em\hangafter=1\noindent \textit{Faire double} étape, \textit{brûler l’étape}. \lxlinks{Iter duplicare.}\par}
\par{\dictfontsize\hangindent=1.5em\hangafter=1\noindent \textit{Faire une} étape \textit{ordinaire}. \lxlinks{Iter justum conficere.}\par}
\par{\dictfontsize\hangindent=1.5em\hangafter=1\noindent État-major, \textit{corps de généraux}. \lxlinks{Principes militum, delecti ordinum militarium, spectatissimi tribuni.}\par}
\par{\dictfontsize\hangindent=1.5em\hangafter=1\noindent \textit{Chef de l’}état-major. \lxlinks{Summus castrorum metator.}\par}
\par{\dictfontsize\hangindent=1.5em\hangafter=1\noindent États \textit{de l’empire, les états}. \lxlinks{Imperii ordines.}\par}
\par{\dictfontsize\hangindent=1.5em\hangafter=1\noindent \textit{Commentaire, exposition, interprétation, explication de l’}Évangile. \lxlinks{Explanatio evangelica, commentarius evangelicus, vitæ actorumque Christi diffusior explanatio.}\par}
\par{\dictfontsize\hangindent=1.5em\hangafter=1\noindent \textit{Relever qqn de} l’excommunication; \textit{lever l’interdit prononcé contre qqn}. \lxlinks{Pacem superum alicui reddere, communioni cœtus sacri restituere aliquem.}\par}
\par{\dictfontsize\hangindent=1.5em\hangafter=1\noindent Excommunier \textit{qqn}; \textit{anathématiser qqn} (\textit{Égl. cath.}). \textit{Lancer, fulminer l’excommunication, l’anathème.} \lxlinks{Devovere aliquem, exsecrari diro carmine, interdicere communione cœtus sacri, interdicere sacris, ejicere e cœtu sacro, usu sacrorum prohibere, exsecratione devovere; dirarum formulam promulgare, ferre in aliquem.}\par}
\par{\dictfontsize\hangindent=1.5em\hangafter=1\noindent Exécuter \textit{un criminel}. \lxlinks{De criminis reo supplicium sumere.}\par}
\par{\dictfontsize\hangindent=1.5em\hangafter=1\noindent Exécuter \textit{la peine encourue}. \lxlinks{Pœnam commissam persequi.}\par}
\par{\dictfontsize\hangindent=1.5em\hangafter=1\noindent Exercer \textit{les matelots}. \lxlinks{Nautas exercere.}\par}
\par{\dictfontsize\hangindent=1.5em\hangafter=1\noindent Exercice; \textit{exercer les soldats}; \textit{faire l’exercice aux soldats}. \lxlinks{Meditatio campestris. Decursio militaris. Componere solvereque ordines edocetur miles. Decursiones militares, voluptaria prœliorum simulacra.}\par}
\par{\dictfontsize\hangindent=1.5em\hangafter=1\noindent \textit{Je le sais, je l’ai appris par} expérience; \textit{j’en ai été instruit par expérience}. \lxlinks{Usu, experiendo didici.}\par}
\par{\dictfontsize\hangindent=1.5em\hangafter=1\noindent \textit{S’}expliquer; \textit{s’exprimer}; \textit{rendre, exprimer sa pensée}. \lxlinks{Exprimere sensa dicendo. Sensa mentis explicare, sensus suos alicui aperire. Promere, quæ in animo latent.}\par}
\par{\dictfontsize\hangindent=1.5em\hangafter=1\noindent Expressions, \textit{termes}; \textit{expressions choisies}; \textit{termes sonores}. \lxlinks{Exquisita verba et sonantia.}\par}
\par{\dictfontsize\hangindent=1.5em\hangafter=1\noindent \textit{Basse} extraction; \textit{un homme de vile origine, de honteuse extraction}. \lxlinks{Dedecus natalium. Homo cum probro natus.}\par}
\par{\dictfontsize\hangindent=1.5em\hangafter=1\noindent Extrait, \textit{abrégé, précis, sommaire} (\textit{d’un livre}). \lxlinks{Breviarium, compendium, commentarius brevis.}\par}
\par{\dictfontsize\hangindent=1.5em\hangafter=1\noindent \textit{Tout hasarder, aller jusqu’au bout, pousser les choses jusqu’à la dernière} extrémité. \lxlinks{Extrema omnia experiri, pati rem ad supremum discrimen adduci: novissimum casum exspectare.}\par}
\par{\dictfontsize\hangindent=0pt\noindent\textls[120]{\scshape f}\par}
\par{\dictfontsize\hangindent=1.5em\hangafter=1\noindent Façon, \textit{main d’œuvre}. \lxlinks{Manus pretium.}\par}
\par{\dictfontsize\hangindent=1.5em\hangafter=1\noindent Faire \{ un legs pieux; \{ une fondation pieuse; \{ un legs, des legs pour des œuvres pies. \lxlinks{Dicare aliquid sacris usibus.}\par}
\par{\dictfontsize\hangindent=1.5em\hangafter=1\noindent \textit{Faire} faillite, \textit{banqueroute}. \lxlinks{Cedere foro, cedere bonis.}\par}
\par{\dictfontsize\hangindent=1.5em\hangafter=1\noindent \textit{Homme} fantasque, \textit{agité de vaines terreurs}. \lxlinks{Homo fanaticus, homo manibus suis malisque umbris actus, \textit{vel} tuta vereri solitus, terroribus exagitatus.}\par}
\par{\dictfontsize\hangindent=1.5em\hangafter=1\noindent Fascines; \textit{fascinage}. \lxlinks{Fasces viminei, fasces fruticum, crates militares, sarmentorum fasciculi.}\par}
\par{\dictfontsize\hangindent=1.5em\hangafter=1\noindent \textit{Faire des} fascines; \textit{faire du fascinage}. \lxlinks{Sarmenta virgultaque colligere.}\par}
\par{\dictfontsize\hangindent=1.5em\hangafter=1\noindent \textit{Ouvrier en} fascines, \textit{fagoteur}; \textit{gabionneur}. \lxlinks{Lignator.}\par}
\par{\dictfontsize\hangindent=1.5em\hangafter=1\noindent \textit{Préparer des} fascines \textit{pour garnir les approches (les tranchées) et les gabions}. \lxlinks{Materiam cædere crateis et pluteis faciendis, struendis, plectendis.}\par}
\par{\dictfontsize\hangindent=1.5em\hangafter=1\noindent \textit{Un copiste infidèle}; faussaire, \textit{c. à d.}, \textit{qui altère un acte authentique, qui fait un faux, ou une fausse signature}. \lxlinks{Librarius mendosus.}\par}
\par{\dictfontsize\hangindent=1.5em\hangafter=1\noindent Fausse \textit{monnaie}. \lxlinks{Adulterini nummi.}\par}
\par{\dictfontsize\hangindent=1.5em\hangafter=1\noindent \textit{Sans perfidie, sans duplicité, sans fraude, sans} fausseté. \lxlinks{Sine fuco.}\par}
\par{\dictfontsize\hangindent=1.5em\hangafter=1\noindent Faux. \lxlinks{Fallax, insidiosus, veteratorius.}\par}
\par{\dictfontsize\hangindent=1.5em\hangafter=1\noindent Faux \textit{cachet}. \lxlinks{Signum adulterinum.}\par}
\par{\dictfontsize\hangindent=1.5em\hangafter=1\noindent \textit{Cacher une haine profonde sous des paroles mielleuses; être un homme} faux, \textit{un traître}. \lxlinks{Odium pectore abditum multis verborum lenociniis obtegere.}\par}
\par{\dictfontsize\hangindent=1.5em\hangafter=1\noindent Fécial, \textit{héraut d’armes qui déclarait la guerre}. \lxlinks{Fecialis.}\par}
\par{\dictfontsize\hangindent=1.5em\hangafter=1\noindent \textit{Droit} fécial. \lxlinks{Jus feciale.}\par}
\par{\dictfontsize\hangindent=1.5em\hangafter=1\noindent Félonie; \textit{action de félon}; \textit{perfidie d’un vassal}. \lxlinks{Crimen in beneficii auctorem commissum.}\par}
\par{\dictfontsize\hangindent=1.5em\hangafter=1\noindent Fête; \textit{solennité}. \lxlinks{Solemnis, festus dies.}\par}
\par{\dictfontsize\hangindent=1.5em\hangafter=1\noindent Fief. \lxlinks{Prædium beneficiarium.}\par}
\par{\dictfontsize\hangindent=1.5em\hangafter=1\noindent \textit{Demander} (\textit{une terre}) \textit{en} fief. \lxlinks{Inaugurationem beneficiorum ab aliquo petere.}\par}
\par{\dictfontsize\hangindent=1.5em\hangafter=1\noindent \textit{Recevoir} (\textit{une terre}) \textit{en} fief. \lxlinks{Ritu beneficiario inaugurari.}\par}
\par{\dictfontsize\hangindent=1.5em\hangafter=1\noindent Figures \textit{de Rhétorique, style figuré}. \lxlinks{Lumina orationis, ornamenta verborum et sententiarum.}\par}
\par{\dictfontsize\hangindent=1.5em\hangafter=1\noindent Fixer, \textit{faire la fixation des dépens, fixer l’amende qui doit être payée par le condamné dans une affaire d’exaction ou de péculat}; \textit{estimer le prix d’une chose disputée, ou régler les frais à payer par celui qui perdra sa cause}. \lxlinks{Litem æstimare.}\par}
\par{\dictfontsize\hangindent=1.5em\hangafter=1\noindent Flancs \textit{d’une armée rangée en bataille, ailes}. (\textit{On dit l’aile droite, l’aile gauche, en parlant d’une armée rangée en bataille.}) \lxlinks{Latera legionis, agminis; alæ.}\par}
\par{\dictfontsize\hangindent=1.5em\hangafter=1\noindent Flanc \textit{d’un bastion, d’un fort}. \lxlinks{Latus propugnaculi.}\par}
\par{\dictfontsize\hangindent=1.5em\hangafter=1\noindent Flancs \textit{dans les ouvrages de fortifications}. \lxlinks{Propugnacula in ambitu, recentiore structura.}\par}
\par{\dictfontsize\hangindent=1.5em\hangafter=1\noindent \textit{Attaquer les} flancs \textit{dégarnis de l’armée ennemie}. \lxlinks{Latere aperto hostem aggredi.}\par}
\par{\dictfontsize\hangindent=1.5em\hangafter=1\noindent \textit{Garnir de cavalerie les} flancs \textit{de l’armée}; \textit{protéger, couvrir les flancs}. \lxlinks{Latera equite munire.}\par}
\par{\dictfontsize\hangindent=1.5em\hangafter=1\noindent \textit{Pour que l’ennemi ne puisse pas prendre, attaquer ses troupes en} flanc. \lxlinks{Ne hostes a lateribus suos pugnantes circumvenire possent.}\par}
\par{\dictfontsize\hangindent=1.5em\hangafter=1\noindent \textit{Serrer l’ennemi par les deux} flancs. \lxlinks{Ab utroque instare, utrumque urgere.}\par}
\par{\dictfontsize\hangindent=1.5em\hangafter=1\noindent \textit{Descendre le} fleuve; \textit{se laisser aller au courant de l’eau, aller en aval, à vau-l’eau}. \lxlinks{Vehi secundo flumine.}\par}
\par{\dictfontsize\hangindent=1.5em\hangafter=1\noindent \textit{Remonter le} fleuve; \textit{naviguer contre le courant, aller d’amont}. \lxlinks{Adversum flumen navis subit.}\par}
\par{\dictfontsize\hangindent=1.5em\hangafter=1\noindent Flux \textit{et reflux,—flot et jusant}. \lxlinks{Maritimi æstus maximi. Maritimorum æstuum accessus et recessus.}\par}
\par{\dictfontsize\hangindent=1.5em\hangafter=1\noindent Fonder \textit{une ville};—\textit{établir une colonie}; \textit{coloniser}. \lxlinks{Coloniam deducere, coloniam constituere in aliquo loco.}\par}
\par{\dictfontsize\hangindent=1.5em\hangafter=1\noindent Force \textit{d’un terme, d’un mot}. \lxlinks{Vis verbi.}\par}
\par{\dictfontsize\hangindent=1.5em\hangafter=1\noindent \textit{De ses propres} forces. \lxlinks{Suo, proprio Marte.}\par}
\par{\dictfontsize\hangindent=1.5em\hangafter=1\noindent Formalités; \textit{étiquette} (\textit{en usage à la cour}). \lxlinks{Formulæ loquendi in aulis principum usurpari solitæ.}\par}
\par{\dictfontsize\hangindent=1.5em\hangafter=1\noindent Formulaire; \textit{formule, forme, formalités}; \textit{définition}—(\textit{t. de droit})—\textit{règle, principe, loi, formule} (\textit{au figuré}). \lxlinks{Formula.}\par}
\par{\dictfontsize\hangindent=1.5em\hangafter=1\noindent \textit{Jurer, prêter serment d’après la} formule \textit{prescrite}. \lxlinks{Jurare in certa verba.}\par}
\par{\dictfontsize\hangindent=1.5em\hangafter=1\noindent Fortifier \textit{une montagne}. \lxlinks{Montem opere circumvenire.}\par}
\par{\dictfontsize\hangindent=1.5em\hangafter=1\noindent Fortifier \textit{une ville}.—\textit{Bloquer, cerner une ville, entourer d’ouvrages} (\textit{une ville qu’on assiége}). \lxlinks{Oppidum vallo castellisque circummunire \textit{seu} claudere, oppidum circumvallare, vallo urbem circummunire, vallum in oppidi circuitu ducere. Oppidum operibus claudere.}\par}
\par{\dictfontsize\hangindent=1.5em\hangafter=1\noindent Fossé, \textit{fosse, excavation, creux}. \lxlinks{Fossa.}\par}
\par{\dictfontsize\hangindent=1.5em\hangafter=1\noindent \textit{Bord du} fossé, \textit{de la fosse}. \lxlinks{Labrum, margo fossæ.}\par}
\par{\dictfontsize\hangindent=1.5em\hangafter=1\noindent \textit{Chariot de bagages}; fourgon. \lxlinks{Currus sarcinis onustus.}\par}
\par{\dictfontsize\hangindent=1.5em\hangafter=1\noindent \textit{Fourrage.} \lxlinks{Pabulum, frumentatio.}\par}
\par{\dictfontsize\hangindent=1.5em\hangafter=1\noindent \textit{Être de piquet au} fourrage. \lxlinks{Pabulatores tueri.}\par}
\par{\dictfontsize\hangindent=1.5em\hangafter=1\noindent \textit{Envoyer la cavalerie} fourrager, \textit{au} fourrage. \lxlinks{Equites dimittere pabulandi causa, milites frumentatum, pabulatum mittere.}\par}
\par{\dictfontsize\hangindent=1.5em\hangafter=1\noindent Fourrier. \lxlinks{Scriba cohortis, legionis.}\par}
\par{\dictfontsize\hangindent=1.5em\hangafter=1\noindent Fourrier. \lxlinks{Curatores sclopetarii tribuni \textit{seu} centurionis.}\par}
\par{\dictfontsize\hangindent=1.5em\hangafter=1\noindent Fourrier \textit{de la cour}. \lxlinks{Curator hospitiorum aulicorum.}\par}
\par{\dictfontsize\hangindent=1.5em\hangafter=1\noindent Fourrier, \textit{sous-officier qui marque les logements, fait la répartition des vivres}. \lxlinks{Hospitiorum designator.}\par}
\par{\dictfontsize\hangindent=1.5em\hangafter=1\noindent \textit{Sergent}-fourrier (\textit{pour l’infanterie}); \textit{maréchal-des-logis} (\textit{pour la cavalerie}). \lxlinks{Designator, curator hospitiorum militarium.}\par}
\par{\dictfontsize\hangindent=1.5em\hangafter=1\noindent Frais de transport; \textit{voiture}; \textit{naulage}; \textit{bachotage}. \lxlinks{Portorium.}\par}
\par{\dictfontsize\hangindent=1.5em\hangafter=1\noindent \textit{Faire} fusiller \textit{qqn}. \lxlinks{Commilitonum glandibus sternendum aliquem objicere.}\par}
\par{\dictfontsize\hangindent=1.5em\hangafter=1\noindent Fusée. \lxlinks{Ignis artificiosus, contestandæ publicæ lætitiæ accensus, recta linea in altum assurgens.}\par}
\par{\dictfontsize\hangindent=1.5em\hangafter=1\noindent \textit{Décharger} (\textit{un} fusil). \lxlinks{Exonerare.}\par}
\par{\dictfontsize\hangindent=1.5em\hangafter=1\noindent \textit{Coup de} fusil. \lxlinks{Ictus sclopi.}\par}
\par{\dictfontsize\hangindent=1.5em\hangafter=1\noindent \textit{Recevoir un coup de} fusil; \textit{être tué d’un coup de fusil}. \lxlinks{Glande plumbea prosterni, transverberari.}\par}
\par{\dictfontsize\hangindent=1.5em\hangafter=1\noindent Fusils; \textit{carabines}; \textit{mousquets}. \lxlinks{Sclopi igniarii majores, ferreæ fistulæ, bombardæ ductiles. Sclopus igniarius.}\par}
\par{\dictfontsize\hangindent=1.5em\hangafter=1\noindent \textit{Charger} (\textit{un} fusil). \lxlinks{Fistulam ferream pulvere nitrato et glande onerare.}\par}
\par{\dictfontsize\hangindent=1.5em\hangafter=1\noindent \textit{Balle de} fusil. \lxlinks{Glans plumbea.}\par}
\par{\dictfontsize\hangindent=1.5em\hangafter=1\noindent \textit{Il est éloigné de la portée d’un} fusil. \lxlinks{Tanto spatio distat, quantum globus e ferrea fistula emissus conficit.}\par}
\par{\dictfontsize\hangindent=1.5em\hangafter=1\noindent \textit{Fantassin}; fusilier; \textit{simple soldat, mousquetaire}. \lxlinks{Levis armaturæ pedes, miles manipularis.}\par}
\par{\dictfontsize\hangindent=0pt\noindent\textls[120]{\scshape g}\par}
\par{\dictfontsize\hangindent=1.5em\hangafter=1\noindent Gabions. \lxlinks{Corbes terra farti, vimineæ cistæ.}\par}
\par{\dictfontsize\hangindent=1.5em\hangafter=1\noindent \textit{Soustraire les} gages \textit{à qqn, priver qqn de ses gages}; \textit{rabattre sur les gages de qqn}. \lxlinks{Stipendium alicujus avertere, stipendio aliquem fraudare.}\par}
\par{\dictfontsize\hangindent=1.5em\hangafter=1\noindent \textit{Salarier, payer qqn; donner, payer à qqn ses} gages, \textit{son salaire, ses appointements, son traitement}. \lxlinks{Stipendium, attributam pensionem, officii pretium alicui constituere, numerare, ferre.}\par}
\par{\dictfontsize\hangindent=1.5em\hangafter=1\noindent Gala (\textit{fêtes, festin à la cour}). \lxlinks{Pompa solemnis aulæ, celebritas aulæ procerum.}\par}
\par{\dictfontsize\hangindent=1.5em\hangafter=1\noindent Galant, \textit{homme poli, de bonnes manières, d’une éducation distinguée}. \lxlinks{Vir omnibus rebus ornatus, ad omnem humanitatem excultus.}\par}
\par{\dictfontsize\hangindent=1.5em\hangafter=1\noindent Galanterie; \textit{politesse, douceur, affabilité dans nos relations avec autrui}. \lxlinks{Comitas, lepor, urbanæ munditiæ.}\par}
\par{\dictfontsize\hangindent=1.5em\hangafter=1\noindent Galerie \textit{dans laquelle les Athlètes s’exerçaient pendant le mauvais temps}. \lxlinks{Tecta ambulatiuncula.}\par}
\par{\dictfontsize\hangindent=1.5em\hangafter=1\noindent Galeries, \textit{portique} (\textit{galerie dont le comble est soutenu par des colonnes, des arcades}), \textit{xyste}. \lxlinks{Intercolumnia ambulationis, xystus, porticus.}\par}
\par{\dictfontsize\hangindent=1.5em\hangafter=1\noindent Garantie \textit{du gouvernement}; \textit{protection de l’État} (\textit{t. de droit}). \lxlinks{Fides publica.}\par}
\par{\dictfontsize\hangindent=1.5em\hangafter=1\noindent Garde \textit{noble}. \lxlinks{Equestris ordinis juvenes, excubias circa cubiculum principis vice militum agentes.}\par}
\par{\dictfontsize\hangindent=1.5em\hangafter=1\noindent Garde \textit{du roi}. \lxlinks{Apparitor regius.}\par}
\par{\dictfontsize\hangindent=1.5em\hangafter=1\noindent Garde \textit{urbaine, garde municipale}; \textit{milice de la ville}. \lxlinks{Urbanæ cohortes.}\par}
\par{\dictfontsize\hangindent=1.5em\hangafter=1\noindent Garde, \textit{cohorte prétorienne}; \textit{garde particulière du général, du conseil, de l’empereur}. \lxlinks{Milites prætoriani, stipatores.}\par}
\par{\dictfontsize\hangindent=1.5em\hangafter=1\noindent \textit{Avant}-garde. \lxlinks{Antecursores ex militia levi.}\par}
\par{\dictfontsize\hangindent=1.5em\hangafter=1\noindent \textit{Préfet du prétoire}; \textit{commandant de la} garde. \lxlinks{Præfectus prætorio \textit{vel} prætoriano militi.}\par}
\par{\dictfontsize\hangindent=1.5em\hangafter=1\noindent \textit{Cohorte prétorienne}, garde \textit{du corps, garde royale, impériale}. \lxlinks{Cohors prætoria, custodiæ corporis regii assueti milites.}\par}
\par{\dictfontsize\hangindent=1.5em\hangafter=1\noindent \textit{Caporaux} (\textit{ceux qui sont exempts de monter la} garde). \textit{Editores tesseræ} (\textit{ceux qui donnent le mot d’ordre} (\textit{milit.})). \lxlinks{Exempti a statione, immunes stationum. Editores tesseræ.}\par}
\par{\dictfontsize\hangindent=1.5em\hangafter=1\noindent \textit{Mettre} garnison \textit{dans une place; occuper une place, une ville}. \lxlinks{Præsidiis locum tenere.}\par}
\par{\dictfontsize\hangindent=1.5em\hangafter=1\noindent \textit{Mettre l’armée en} garnison, \textit{en quartier}. \lxlinks{Exercitum per oppida dispertire.}\par}
\par{\dictfontsize\hangindent=1.5em\hangafter=1\noindent \textit{Renforcer la} garnison. \lxlinks{Præsidium majoribus copiis firmare.}\par}
\par{\dictfontsize\hangindent=1.5em\hangafter=1\noindent Général \textit{de brigade}. \lxlinks{Summus vigiliarum præfectus.}\par}
\par{\dictfontsize\hangindent=1.5em\hangafter=1\noindent Général \textit{de cavalerie}. \lxlinks{Summus equitum præfectus.}\par}
\par{\dictfontsize\hangindent=1.5em\hangafter=1\noindent \textit{Lieutenant}-général, \textit{général de division}. \lxlinks{Legatus, summi ducis vicarius, qui legati locum obtinet, qui secundum locum imperii tenet.}\par}
\par{\dictfontsize\hangindent=1.5em\hangafter=1\noindent Général \textit{en chef, généralissime, maréchal}; \textit{feld-maréchal} (\textit{allem. et autr.}). \lxlinks{Supremus castrorum præfectus.}\par}
\par{\dictfontsize\hangindent=1.5em\hangafter=1\noindent Général \textit{en chef de l’artillerie, grand-maître de l’artillerie—lieutenant-général d’artillerie} (\textit{jadis}). \lxlinks{Rei tormentariæ ad exercitum pertinentis præfectus summus.}\par}
\par{\dictfontsize\hangindent=1.5em\hangafter=1\noindent Général \textit{d’infanterie}. \lxlinks{Peditum summus præfectus.}\par}
\par{\dictfontsize\hangindent=1.5em\hangafter=1\noindent \textit{Vaguemestre} général. \lxlinks{Supremus rei vehicularis in exercitu præfectus.}\par}
\par{\dictfontsize\hangindent=1.5em\hangafter=1\noindent \textit{Commissaire} général \textit{des armées}; \textit{intendant militaire en chef}. \lxlinks{Supremus exercitus curator.}\par}
\par{\dictfontsize\hangindent=1.5em\hangafter=1\noindent \textit{Aide-de-camp} (\textit{d’un} général). \lxlinks{Adjutor ducis \textit{seu} præfecti, qui mandata ductoris perfert.}\par}
\par{\dictfontsize\hangindent=1.5em\hangafter=1\noindent Généalogie. \lxlinks{Origo familiarum, ex qua clarorum virorum propagines possint cognosci. Familiam alicujus inde a primordiis ad nostram ætatem recensere.}\par}
\par{\dictfontsize\hangindent=1.5em\hangafter=1\noindent Généralissime. \lxlinks{Imperator summus, ad quem unum summa imperii redit.}\par}
\par{\dictfontsize\hangindent=1.5em\hangafter=1\noindent \textit{Ingénieur militaire}; \textit{officier du} génie. \lxlinks{Castrorum metator, architectus militaris, artis muniendi magister, machinator operum bellicorum, machinali scientia clarus.}\par}
\par{\dictfontsize\hangindent=1.5em\hangafter=1\noindent \textit{Un} gentilhomme \textit{d’ancienne extraction, de vieille souche}. \lxlinks{Avitæ nobilitatis, multarum imaginum.}\par}
\par{\dictfontsize\hangindent=1.5em\hangafter=1\noindent \textit{Un} gentilhomme \textit{de fraîche date}. \lxlinks{Homo novus.}\par}
\par{\dictfontsize\hangindent=1.5em\hangafter=1\noindent Géographie. \lxlinks{Sedes locorum, regionum descriptio.}\par}
\par{\dictfontsize\hangindent=1.5em\hangafter=1\noindent Géôlier. \lxlinks{Triumvir carceris.}\par}
\par{\dictfontsize\hangindent=1.5em\hangafter=1\noindent Géométrie. \lxlinks{Studium dimetiendæ terræ.}\par}
\par{\dictfontsize\hangindent=1.5em\hangafter=1\noindent Giberne. \lxlinks{Theca sclopetaria pulveris pyrii, theca qua sclopi fartura custoditur.}\par}
\par{\dictfontsize\hangindent=1.5em\hangafter=1\noindent \textit{Batterie de} gobelets. \lxlinks{Caliculi quibus agyrtæ circumstantium oculos perstringunt.}\par}
\par{\dictfontsize\hangindent=1.5em\hangafter=1\noindent Golfe. \lxlinks{Sinus.}\par}
\par{\dictfontsize\hangindent=1.5em\hangafter=1\noindent Gondole; \textit{barque} (\textit{vénitienne}). \lxlinks{Navicula Venetorum, navis velox, biremis Veneta, Venetorum phaselus.}\par}
\par{\dictfontsize\hangindent=1.5em\hangafter=1\noindent \textit{Soldats du train}, goujats. \lxlinks{Calones.}\par}
\par{\dictfontsize\hangindent=1.5em\hangafter=1\noindent Gouverneur; \textit{président} (\textit{d’une province en Allemagne}). \lxlinks{Præfectus præsidii, provinciæ rector.}\par}
\par{\dictfontsize\hangindent=1.5em\hangafter=1\noindent \textit{Vice}-gouverneur; \textit{vice-président} (\textit{Allemagne}). \lxlinks{Provinciæ legatus, præsidio præfecti vicarius.}\par}
\par{\dictfontsize\hangindent=1.5em\hangafter=1\noindent Grands (\textit{d’Espagne}). \lxlinks{Optimates Hispanorum primi ordinis, primæ admissionis.}\par}
\par{\dictfontsize\hangindent=1.5em\hangafter=1\noindent Gradué. \lxlinks{Honoribus Academicis ornatus.}\par}
\par{\dictfontsize\hangindent=1.5em\hangafter=1\noindent Gratification. \lxlinks{Congiarium.}\par}
\par{\dictfontsize\hangindent=1.5em\hangafter=1\noindent Grenadier; \textit{fusilier}. \lxlinks{Miles qui glandes pulvere nitrato fartas manu ejaculatur.}\par}
\par{\dictfontsize\hangindent=1.5em\hangafter=1\noindent Grotte. \lxlinks{Crypta hortorum artificiosa, apta vitandis æstivis caloribus.}\par}
\par{\dictfontsize\hangindent=1.5em\hangafter=1\noindent Guerre \textit{sans caractère public}; \textit{guerre qui a pour motif le droit du plus fort}, \textit{c. à d.}, \textit{une guerre illégitime, illégale, injuste}. \lxlinks{Bella privata, privata auctoritate indicta.}\par}
\par{\dictfontsize\hangindent=1.5em\hangafter=1\noindent \textit{Faire une} guerre \textit{défensive}. \lxlinks{Bellum illatum repellere, defendere, propulsare.}\par}
\par{\dictfontsize\hangindent=1.5em\hangafter=1\noindent \textit{Bagages}; \textit{matériel de} guerre, \textit{munitions}. \lxlinks{Instrumentum militare, instrumentum militiæ.}\par}
\par{\dictfontsize\hangindent=1.5em\hangafter=1\noindent \textit{Motifs de la} guerre; \textit{son utilité}. \lxlinks{Ratio belli, rei militaris.}\par}
\par{\dictfontsize\hangindent=1.5em\hangafter=1\noindent \textit{Théâtre de la} guerre. \lxlinks{Belli sedes, bellum hic vel ibi constitit, collatum est.}\par}
\par{\dictfontsize\hangindent=0pt\noindent\textls[120]{\scshape h}\par}
\par{\dictfontsize\hangindent=1.5em\hangafter=1\noindent Hallebarde. \lxlinks{Hastata securis. Securis hastæ præfixa.}\par}
\par{\dictfontsize\hangindent=1.5em\hangafter=1\noindent \textit{Nombre}, harmonie \textit{d’un discours}. \lxlinks{Oratio ad numerum cadens.}\par}
\par{\dictfontsize\hangindent=1.5em\hangafter=1\noindent \textit{Orateur qui a la phrase} harmonieuse (\textit{pour le style et pour le débit}). \lxlinks{Cujus aures perfecto et completo verborum numero gaudent.}\par}
\par{\dictfontsize\hangindent=1.5em\hangafter=1\noindent Haste (\textit{lance, javelot sans fer}). \lxlinks{Hasta pura.}\par}
\par{\dictfontsize\hangindent=1.5em\hangafter=1\noindent Heiduques. \lxlinks{Pedites Hungari. Milites provinciales. Hungariæ pedites.}\par}
\par{\dictfontsize\hangindent=1.5em\hangafter=1\noindent \textit{Publier de tous côtés par la voix des} hérauts. \lxlinks{Præconibus circummissis edicere.}\par}
\par{\dictfontsize\hangindent=1.5em\hangafter=1\noindent \textit{Royaume} héréditaire; \textit{états héréditaires}. \lxlinks{Regnum patrium, regnum quod nascendi jure, hereditate obvenit.}\par}
\par{\dictfontsize\hangindent=1.5em\hangafter=1\noindent Héritier \textit{naturel, légitime, ab intestat}. \lxlinks{Ab intestato, cui nascendi jura hereditatem addicunt.}\par}
\par{\dictfontsize\hangindent=1.5em\hangafter=1\noindent Hommage \textit{lige}; \textit{devoir de vassal}. \lxlinks{Jus beneficiarium, nexus beneficiarius.}\par}
\par{\dictfontsize\hangindent=1.5em\hangafter=1\noindent Homme d’affaires \textit{de qqn, son gérant, régisseur, procureur, agent}. \lxlinks{Curator rei privatæ, qui procurat alterius rationes et negotia.}\par}
\par{\dictfontsize\hangindent=1.5em\hangafter=1\noindent \textit{Des} hommes, \textit{des gens de qualité, de haut parage}. \lxlinks{Viri amplissimi.}\par}
\par{\dictfontsize\hangindent=1.5em\hangafter=1\noindent \textit{Mon} honneur \textit{est en jeu, y est engagé}; \textit{il s’agit de mon honneur}. \lxlinks{Agitur de honore meo.}\par}
\par{\dictfontsize\hangindent=1.5em\hangafter=1\noindent \textit{Point d’}honneur. \lxlinks{Summa honoris.}\par}
\par{\dictfontsize\hangindent=1.5em\hangafter=1\noindent Hôpital; \textit{hôtel-Dieu}; lazaret. \lxlinks{Valetudinarium publicum.}\par}
\par{\dictfontsize\hangindent=1.5em\hangafter=1\noindent \textit{Tireur d’}horoscope. \lxlinks{Qui notato ortu, futura hominibus edicit, fatorum est astris vanus interpres.}\par}
\par{\dictfontsize\hangindent=1.5em\hangafter=1\noindent \textit{Dresser ou faire tirer, faire, dresser} l’horoscope \textit{de qqn}. \lxlinks{Divinatio natalitia; prædicere, divinare quo quis fato natus sit.}\par}
\par{\dictfontsize\hangindent=1.5em\hangafter=1\noindent Huguenot. \lxlinks{Religioni reformatæ initiatus.}\par}
\par{\dictfontsize\hangindent=1.5em\hangafter=1\noindent Huissier. \lxlinks{Accensus.}\par}
\par{\dictfontsize\hangindent=0pt\noindent\textls[120]{\scshape i}\par}
\par{\dictfontsize\hangindent=1.5em\hangafter=1\noindent Ici, \textit{là il y a des mines fort riches}. \lxlinks{Metalla exercentur hic vel ibi magno proventu.}\par}
\par{\dictfontsize\hangindent=1.5em\hangafter=1\noindent \textit{Faire des courses}, incursions \textit{sur le territoire ennemi. Faire la petite guerre, la guerre de guérilla, la guerre en partisan.} \lxlinks{Excursiones in hostilem agrum facere. Excursiones. Excursionibus se locupletare.}\par}
\par{\dictfontsize\hangindent=1.5em\hangafter=1\noindent \textit{Compagnie des} Indes \textit{orientales}. \lxlinks{Societas Indiæ orientalis, collegium mercatorum, qui orientalis Indiæ merces permutant.}\par}
\par{\dictfontsize\hangindent=1.5em\hangafter=1\noindent \textit{Je reste indifférent, cela me laisse} indifférent, \textit{dans une indifférence complète}; \textit{ce n’est pour moi, ni un sujet de joie, ni un sujet de crainte}. \lxlinks{Non est quod gaudeam metuamve.}\par}
\par{\dictfontsize\hangindent=1.5em\hangafter=1\noindent \textit{Choses} indifférentes. \lxlinks{Res neutræ.}\par}
\par{\dictfontsize\hangindent=1.5em\hangafter=1\noindent Indulgence \textit{papale}. \lxlinks{Noxarum condonationes a supremo Sacrorum Antistite concessæ.}\par}
\par{\dictfontsize\hangindent=1.5em\hangafter=1\noindent \textit{Sous}-inféoder. \lxlinks{Prædium beneficiarium transferre in alium eodem jure.}\par}
\par{\dictfontsize\hangindent=1.5em\hangafter=1\noindent Influence \textit{des astres, de la lune}. \lxlinks{Vis stellarum. Lunæ tractus.}\par}
\par{\dictfontsize\hangindent=1.5em\hangafter=1\noindent \textit{Astrologie, science prétendue de l’influence des astres; divination.} \lxlinks{Astrologia divinatrix.}\par}
\par{\dictfontsize\hangindent=1.5em\hangafter=1\noindent Ingénieur \textit{militaire}; \textit{officier du génie}. \lxlinks{Castrorum metator, architectus militaris, artis muniendi magister, machinator operum bellicorum, machinali scientia clarus.}\par}
\par{\dictfontsize\hangindent=1.5em\hangafter=1\noindent Ingénieur \textit{civil}. \lxlinks{Machinali scientia clarus; machinator.}\par}
\par{\dictfontsize\hangindent=1.5em\hangafter=1\noindent Ingénieur \textit{des ponts et chaussées}. \lxlinks{Architectus pontibus viisque curandis præpositus; viarum pontiumque curatores.}\par}
\par{\dictfontsize\hangindent=1.5em\hangafter=1\noindent Inquisition, \textit{tribunal de l’inquisition}. \lxlinks{Fidei quæsitorum tribunal.}\par}
\par{\dictfontsize\hangindent=1.5em\hangafter=1\noindent \textit{Prendre une} inscription \textit{à l’université}. \lxlinks{Initiatio academica. Albo studiosorum inserere. Studiis initiare aliquem, ritibus academicis.}\par}
\par{\dictfontsize\hangindent=1.5em\hangafter=1\noindent \textit{Intendant}, inspecteur \textit{des contributions} (\textit{indirectes}). \lxlinks{Annonæ, redituum regiorum curator.}\par}
\par{\dictfontsize\hangindent=1.5em\hangafter=1\noindent Inspiration \textit{divine}; \textit{souffle d’en haut}. \lxlinks{Afflatus divinus, Numinis impulsus.}\par}
\par{\dictfontsize\hangindent=1.5em\hangafter=1\noindent Installation (\textit{en parlant des personnes}). \lxlinks{Legitima collocatio sacri census.}\par}
\par{\dictfontsize\hangindent=1.5em\hangafter=1\noindent \textit{Installer qqn.} \lxlinks{Inaugurare aliquem.}\par}
\par{\dictfontsize\hangindent=1.5em\hangafter=1\noindent \textit{Instituer qqn son} légataire, \textit{héritier universel}. \lxlinks{Ex asse heredem scribere.}\par}
\par{\dictfontsize\hangindent=1.5em\hangafter=1\noindent Intendant \textit{militaire en chef}. \lxlinks{Primarius negotiorum bellicorum curator.}\par}
\par{\dictfontsize\hangindent=1.5em\hangafter=1\noindent Intendant \textit{des ouvriers}. \lxlinks{Præfectus fabrum.}\par}
\par{\dictfontsize\hangindent=1.5em\hangafter=1\noindent Intérêt \textit{exorbitant, usuraire}. \lxlinks{Iniquissimum fenus.}\par}
\par{\dictfontsize\hangindent=1.5em\hangafter=1\noindent \textit{S’intéresser, prendre} intérêt \textit{au bien public}. \lxlinks{Consulere communi saluti, dignitati. In commune consulere.}\par}
\par{\dictfontsize\hangindent=1.5em\hangafter=1\noindent Intérêts. \lxlinks{Fenus, usura, pecuniæ reditus.}\par}
\par{\dictfontsize\hangindent=1.5em\hangafter=1\noindent \textit{Faire de son mieux pour tous les deux}; \textit{servir les} intérêts \textit{de l’un et de l’autre}. \lxlinks{Utriusque facere causa.}\par}
\par{\dictfontsize\hangindent=1.5em\hangafter=1\noindent \textit{Chercher le bien de qqn; aider qqn de ses conseils}; \textit{prendre à cœur les} intérêts \textit{de qqn}. \lxlinks{Consulere alicui; consilio, opera, studio adesse alicui.}\par}
\par{\dictfontsize\hangindent=1.5em\hangafter=1\noindent \textit{Gros} intérêts. \lxlinks{Gravissimæ usuræ.}\par}
\par{\dictfontsize\hangindent=1.5em\hangafter=1\noindent Interrègne. \lxlinks{Interregnum.}\par}
\par{\dictfontsize\hangindent=1.5em\hangafter=1\noindent Inventorier; \textit{dresser un inventaire}. \lxlinks{Commentarium bonorum conscribere.}\par}
\par{\dictfontsize\hangindent=1.5em\hangafter=1\noindent Invention, \textit{découverte}. \lxlinks{Inventum callidum, cogitatum acutum.}\par}
\par{\dictfontsize\hangindent=1.5em\hangafter=1\noindent Investiture (\textit{moyen-âge}). \lxlinks{Solemnia inaugurationis in munus sacrum.}\par}
\par{\dictfontsize\hangindent=1.5em\hangafter=1\noindent Imposer \textit{des marchandises étrangères}; \textit{mettre} l’accise \textit{sur les importations}. \lxlinks{Peregrinarum mercium portoria constituere.}\par}
\par{\dictfontsize\hangindent=1.5em\hangafter=1\noindent Impôt \textit{sur les consommations}. \lxlinks{Vectigal eorum quæ consumuntur, census rerum venalium.}\par}
\par{\dictfontsize\hangindent=1.5em\hangafter=1\noindent Impôt \textit{pour la guerre contre la Turquie}. \lxlinks{Pecunia in bellum Turcicum collata.}\par}
\par{\dictfontsize\hangindent=1.5em\hangafter=1\noindent Impôt \textit{ou droit sur le vin}. \lxlinks{Exactio vinaria.}\par}
\par{\dictfontsize\hangindent=1.5em\hangafter=1\noindent Impôt \textit{personnel, capitation}. \lxlinks{Exactio in capita.}\par}
\par{\dictfontsize\hangindent=1.5em\hangafter=1\noindent \textit{Accise}, impôt \textit{sur le vin}. \lxlinks{Exactio vinaria.}\par}
\par{\dictfontsize\hangindent=1.5em\hangafter=1\noindent \textit{Être exempt} d’impôts. \lxlinks{Omnium onerum nancisci immunitatem.}\par}
\par{\dictfontsize\hangindent=1.5em\hangafter=1\noindent Improvisateur.—\textit{Improviser, parler, décider sur-le-champ, sans préparation.} \lxlinks{Extemporalis, vir subita expedire paratus. Subita proferre, dicere, disceptare.}\par}
\par{\dictfontsize\hangindent=1.5em\hangafter=1\noindent Irrégulier. \lxlinks{Declinans a regula, præter regulam.}\par}
\par{\dictfontsize\hangindent=0pt\noindent\textls[120]{\scshape j}\par}
\par{\dictfontsize\hangindent=1.5em\hangafter=1\noindent Janissaires. \lxlinks{Turcarum pedites statarii, pedites prætoriani Turcarum Imperatoris.}\par}
\par{\dictfontsize\hangindent=1.5em\hangafter=1\noindent Journal, \textit{gazette}. \lxlinks{Ephemerides rerum novarum, commentarii rerum publicarum.}\par}
\par{\dictfontsize\hangindent=1.5em\hangafter=1\noindent Journal; \textit{livre-journal}. \lxlinks{Libellus in quem quotidiana acta referuntur.}\par}
\par{\dictfontsize\hangindent=1.5em\hangafter=1\noindent \textit{Livre-journal}; journal. \lxlinks{Commentarius, annales.}\par}
\par{\dictfontsize\hangindent=1.5em\hangafter=1\noindent Juge \textit{militaire}; auditeur, \textit{grand prévôt. Juges du camp, ceux qui dans les combats judiciaires, les joûtes... étaient chargés de veiller à ce que tout se passât selon l’usage}. \lxlinks{Judex militaris, castrensis, quæsitor causarum militarium.}\par}
\par{\dictfontsize\hangindent=1.5em\hangafter=1\noindent Jugement \textit{dernier}. \lxlinks{Decretorius ille dies generis humani, dies novissimus, quo Christus ad judicia exercenda descendet.}\par}
\par{\dictfontsize\hangindent=1.5em\hangafter=1\noindent \textit{De ce pas, je vous ferai citer en} justice. \lxlinks{Faciam ut a prætore tibi statim dies dicatur.}\par}
\par{\dictfontsize\hangindent=0pt\noindent\textls[120]{\scshape l}\par}
\par{\dictfontsize\hangindent=1.5em\hangafter=1\noindent Lance \textit{armée d’un fer}. \lxlinks{Præpilata.}\par}
\par{\dictfontsize\hangindent=1.5em\hangafter=1\noindent Langue \textit{maternelle}. \lxlinks{Sermo patrius.}\par}
\par{\dictfontsize\hangindent=1.5em\hangafter=1\noindent L’égalité. \lxlinks{Observatio legum, juris.}\par}
\par{\dictfontsize\hangindent=1.5em\hangafter=1\noindent \textit{Hôpital, hôtel-Dieu}; Lazaret. \lxlinks{Valetudinarium publicum.}\par}
\par{\dictfontsize\hangindent=1.5em\hangafter=1\noindent Légitime (\textit{subst. fém.}): \textit{portion que la loi assure aux enfants sur les biens du père et de la mère} (\textit{t. de droit}). \lxlinks{Pars bonorum quæ hereditate inter filios dividitur.}\par}
\par{\dictfontsize\hangindent=1.5em\hangafter=1\noindent Lettre \textit{de change}. \lxlinks{Litteræ argentariæ.}\par}
\par{\dictfontsize\hangindent=1.5em\hangafter=1\noindent Lettre \textit{de créance}. \lxlinks{Litteræ fidei publicæ, auctoritas Legati, litteræ quæ fidem legationis adstruunt.}\par}
\par{\dictfontsize\hangindent=1.5em\hangafter=1\noindent Lettre \textit{d’investiture}. \lxlinks{Litteræ beneficiariæ.}\par}
\par{\dictfontsize\hangindent=1.5em\hangafter=1\noindent Lettres \textit{de représentation}. \lxlinks{Litteræ commendatitiæ recens electi.}\par}
\par{\dictfontsize\hangindent=1.5em\hangafter=1\noindent Lettre \textit{perfide, contenant la perte de celui qui en est le porteur}. \lxlinks{Litteræ Bellerophonteæ.}\par}
\par{\dictfontsize\hangindent=1.5em\hangafter=1\noindent \textit{Adresser une} lettre \textit{à qqn}. \lxlinks{Inscribere alicui perferendas litteras.}\par}
\par{\dictfontsize\hangindent=1.5em\hangafter=1\noindent \textit{Élite de l’armée}; \textit{troupes} levées \textit{pour défendre une province}. \lxlinks{Delecta manus, delecti milites.}\par}
\par{\dictfontsize\hangindent=1.5em\hangafter=1\noindent Lexique; \textit{dictionnaire}; \textit{vocabulaire}. \lxlinks{Liber vocum notiones complexus, libri qui dictionum sensa collegerunt.}\par}
\par{\dictfontsize\hangindent=1.5em\hangafter=1\noindent Licencié. \lxlinks{Primos a doctore honores consecutus.}\par}
\par{\dictfontsize\hangindent=1.5em\hangafter=1\noindent Licencier, \textit{renvoyer du service, congédier}. (\textit{milit.}) \lxlinks{Exauctorare.}\par}
\par{\dictfontsize\hangindent=1.5em\hangafter=1\noindent Lieutenant. \lxlinks{Centurionis vicarius, procenturio.}\par}
\par{\dictfontsize\hangindent=1.5em\hangafter=1\noindent Lignes \textit{de communication} (\textit{t. de fortif.}). \lxlinks{Brachium fossarum, linea communis, quæ vulgo communionis.}\par}
\par{\dictfontsize\hangindent=1.5em\hangafter=1\noindent \textit{Entourer avec une armée} toute la ligne \textit{des fortifications}. \lxlinks{Totum opus militum corona cingere.}\par}
\par{\dictfontsize\hangindent=1.5em\hangafter=1\noindent \textit{Garnir de troupes toute la} ligne \textit{des fortifications, des retranchements}. \lxlinks{Totum opus militum corona cingere.}\par}
\par{\dictfontsize\hangindent=1.5em\hangafter=1\noindent \textit{Ouvrir, établir des} lignes \textit{de communication entre le château et le camp}. \lxlinks{Castellum brachiis cum opere castrorum conjungere.}\par}
\par{\dictfontsize\hangindent=1.5em\hangafter=1\noindent Livre \textit{des recettes et des dépenses}. \lxlinks{Codex accepti et expensi.}\par}
\par{\dictfontsize\hangindent=1.5em\hangafter=1\noindent Recueil \textit{de notes, mémoire, mémorial}; livre \textit{de raison ou grand livre}, \textit{c. à d.}, \textit{un registre où les négociants portent tout leur compte par doit et avoir}. \lxlinks{Electorum commentarius.}\par}
\par{\dictfontsize\hangindent=1.5em\hangafter=1\noindent Loterie. \lxlinks{Ludus, qui sorte facta constat, urna sortis fortuitæ.}\par}
\par{\dictfontsize\hangindent=1.5em\hangafter=1\noindent Demi-lune (\textit{fortific.}), lunettes (\textit{fortif.}). \lxlinks{Lunatum munimentum, propugnaculum.}\par}
\par{\dictfontsize\hangindent=0pt\noindent\textls[120]{\scshape m}\par}
\par{\dictfontsize\hangindent=1.5em\hangafter=1\noindent \textit{De ma} (\textit{ta, sa}) \textit{propre} main. \lxlinks{Manu mea, tua, sua.}\par}
\par{\dictfontsize\hangindent=1.5em\hangafter=1\noindent \textit{Mettre une} maison \textit{en vente}. \lxlinks{Ædes venales inscribere.}\par}
\par{\dictfontsize\hangindent=1.5em\hangafter=1\noindent Majeur. \lxlinks{Aliena tutela emancipatus, liberam rerum suarum administrationem adeptus.}\par}
\par{\dictfontsize\hangindent=1.5em\hangafter=1\noindent \textit{Fidéicommis.} Majorat, \textit{celui qui possède un} majorat. \textit{Fideicommissaire.} \lxlinks{Heres fiduciarius.}\par}
\par{\dictfontsize\hangindent=1.5em\hangafter=1\noindent Mandat \textit{d’arrét}. \lxlinks{Litteræ facinorosi comprehendendi causa missæ.}\par}
\par{\dictfontsize\hangindent=1.5em\hangafter=1\noindent \textit{Champ de Mars; place d’armes; champ de} manœuvre. \textit{Rendez-vous des troupes.} \lxlinks{Campus martius, diribitorium.}\par}
\par{\dictfontsize\hangindent=1.5em\hangafter=1\noindent Manufacture; \textit{fabrique}. \lxlinks{Officina operum manu factorum.}\par}
\par{\dictfontsize\hangindent=1.5em\hangafter=1\noindent Marche. \lxlinks{Iter.}\par}
\par{\dictfontsize\hangindent=1.5em\hangafter=1\noindent Marche \textit{forcée}. \lxlinks{Magnum iter, magnum itineris spatium.}\par}
\par{\dictfontsize\hangindent=1.5em\hangafter=1\noindent \textit{Battre, sonner la} marche; \textit{sonner la retraite}. \lxlinks{Receptui canere, signum dare, signum buccina dare, tuba revocare milites.}\par}
\par{\dictfontsize\hangindent=1.5em\hangafter=1\noindent \textit{Se mettre en} marche, \textit{en route; partir}. \lxlinks{Iter inire, aggredi.}\par}
\par{\dictfontsize\hangindent=1.5em\hangafter=1\noindent \textit{Observer la} marche \textit{de l’armée ennemie}. \lxlinks{Itinera hostium observare, explorare.}\par}
\par{\dictfontsize\hangindent=1.5em\hangafter=1\noindent \textit{Contre}-marche. \lxlinks{Itinera transversa.}\par}
\par{\dictfontsize\hangindent=1.5em\hangafter=1\noindent \textit{Feld}-maréchal (\textit{en Allemagne et en Autriche}), \textit{maréchal} (\textit{en France}), \textit{généralissime}; \textit{général en chef. Commandant en chef.} \lxlinks{Summus castrorum præfectus.}\par}
\par{\dictfontsize\hangindent=1.5em\hangafter=1\noindent Maréchalat; \textit{Feld-maréchalat} (\textit{Allem. et Aut.}). \lxlinks{Supremi castrorum præfecti vicarius.}\par}
\par{\dictfontsize\hangindent=1.5em\hangafter=1\noindent \textit{Donner la bénédiction nuptiale à}; marier. \lxlinks{Ritibus sacris jungere connubii fœdera, solemni precatione nuptias lustrare.}\par}
\par{\dictfontsize\hangindent=1.5em\hangafter=1\noindent Matricule. \lxlinks{Album, formula censendi.}\par}
\par{\dictfontsize\hangindent=1.5em\hangafter=1\noindent Maxime, \textit{règle, principe}. \lxlinks{Regula.}\par}
\par{\dictfontsize\hangindent=1.5em\hangafter=1\noindent Mausolée; \textit{un monument élevé pour témoigner de la douleur publique}. \lxlinks{Mausolæum; monumentum exsequiarum causa, contestando publico dolori positum.}\par}
\par{\dictfontsize\hangindent=1.5em\hangafter=1\noindent Mèche. \lxlinks{Funis e fomite. Incendiarii funes.}\par}
\par{\dictfontsize\hangindent=1.5em\hangafter=1\noindent Mécontent, \textit{malcontent, de mauvaise humeur}. \lxlinks{Animus alienus et offensus.}\par}
\par{\dictfontsize\hangindent=1.5em\hangafter=1\noindent \textit{Médaille commémorative}; médaille. \lxlinks{Nummus in memoriam rei singularis, facti, relatæ victoriæ cusus, nummus vetustatis gratia asservatus.}\par}
\par{\dictfontsize\hangindent=1.5em\hangafter=1\noindent Médaillon. \lxlinks{Nummus modi majoris, maximi moduli.}\par}
\par{\dictfontsize\hangindent=1.5em\hangafter=1\noindent \textit{Exercer la médecine}; \textit{être} médecin. \lxlinks{Medicinam exercere.}\par}
\par{\dictfontsize\hangindent=1.5em\hangafter=1\noindent \textit{Nourriture obtenue à force de prières, en demandant l’aumône}; mendicité. \lxlinks{Victus precarius.}\par}
\par{\dictfontsize\hangindent=1.5em\hangafter=1\noindent \textit{Dire la} messe, \textit{célébrer le saint Sacrifice}. \lxlinks{Rem divinam facere, sacris operari.}\par}
\par{\dictfontsize\hangindent=1.5em\hangafter=1\noindent \textit{Dire la} messe \textit{pour les défunts, la messe des morts}; \textit{chanter une messe de requiem, un requiem}; \textit{célébrer un office pour les morts, faire le service} (\textit{funèbre}) \textit{pour} (\textit{tel défunt}). \lxlinks{Sacrum piaculare facere, expiare manes defunctorum.}\par}
\par{\dictfontsize\hangindent=1.5em\hangafter=1\noindent \textit{Ouïr, entendre la} messe, \textit{assister à la Messe, au saint Sacrifice}. \lxlinks{Rebus divinis interesse.}\par}
\par{\dictfontsize\hangindent=1.5em\hangafter=1\noindent Mesure (\textit{musiq.}); \textit{cadence}. \lxlinks{Modus vocum, numerus in cantu.}\par}
\par{\dictfontsize\hangindent=1.5em\hangafter=1\noindent Meurtrières (\textit{ouvertures dans les remparts pour tirer à couvert}). \lxlinks{Gula propugnaculi.}\par}
\par{\dictfontsize\hangindent=1.5em\hangafter=1\noindent Milices. \lxlinks{Milites provinciales, armati populares, raptim ad arma convocati.}\par}
\par{\dictfontsize\hangindent=1.5em\hangafter=1\noindent Mine. \lxlinks{Cuniculus. Furnulus. Sinuosi mæandri, pulvere pyrio farti.}\par}
\par{\dictfontsize\hangindent=1.5em\hangafter=1\noindent \textit{Une} mine \textit{qu’on a fait jouer, sauter}. \lxlinks{Cuniculi accensi.}\par}
\par{\dictfontsize\hangindent=1.5em\hangafter=1\noindent Miner \textit{un ouvrage} (\textit{de fortification}), \textit{le faire sauter}. \lxlinks{Cuniculis aggerem subtrahere.}\par}
\par{\dictfontsize\hangindent=1.5em\hangafter=1\noindent Minéral; \textit{fossile}. \lxlinks{Fossilia.}\par}
\par{\dictfontsize\hangindent=1.5em\hangafter=1\noindent Mines \textit{de fer}. \lxlinks{Secturæ ferrariæ, fodinæ.}\par}
\par{\dictfontsize\hangindent=1.5em\hangafter=1\noindent \textit{Les} mines \textit{firent une large brèche aux remparts}. \lxlinks{Cuniculo suffossa mœnia ingens nudarunt spatium.}\par}
\par{\dictfontsize\hangindent=1.5em\hangafter=1\noindent \textit{Faire des} mines. \lxlinks{Cuniculos rectos agere.}\par}
\par{\dictfontsize\hangindent=1.5em\hangafter=1\noindent \textit{Faire des brèches aux remparts à l’aide des} mines. \lxlinks{Cuniculis muros subruere.}\par}
\par{\dictfontsize\hangindent=1.5em\hangafter=1\noindent Ici, \textit{là, il y a des} mines \textit{fort riches}. \lxlinks{Metalla exercentur hic vel ibi magno proventu.}\par}
\par{\dictfontsize\hangindent=1.5em\hangafter=1\noindent \textit{Revenu des} mines. \lxlinks{Pecunia quæ ex metallis redit.}\par}
\par{\dictfontsize\hangindent=1.5em\hangafter=1\noindent \textit{Contre}-mines. \lxlinks{Cuniculi adversi, oppositi. Transversus meatus.}\par}
\par{\dictfontsize\hangindent=1.5em\hangafter=1\noindent Mineur. \lxlinks{Nondum rerum suarum liberum usum adeptus.}\par}
\par{\dictfontsize\hangindent=1.5em\hangafter=1\noindent Missel. \lxlinks{Liber sacer ritualis, liber liturgiæ.}\par}
\par{\dictfontsize\hangindent=1.5em\hangafter=1\noindent \textit{Coiffé d’une} mitre. \lxlinks{Infulatus, mitratus, sacris infulis conspicuus.}\par}
\par{\dictfontsize\hangindent=1.5em\hangafter=1\noindent Modèle, \textit{exemple}. \lxlinks{Exemplar; modus formaque.}\par}
\par{\dictfontsize\hangindent=1.5em\hangafter=1\noindent \textit{Être fidèle à son} modèle. \lxlinks{Ab exemplari pendere.}\par}
\par{\dictfontsize\hangindent=1.5em\hangafter=1\noindent \textit{Avoir de bonnes} mœurs, \textit{avoir une conduite irréprochable}. \lxlinks{Moribus se probare.}\par}
\par{\dictfontsize\hangindent=1.5em\hangafter=1\noindent Mœurs \textit{légères, conduite légère, facile}. \lxlinks{Mores facillimi.}\par}
\par{\dictfontsize\hangindent=1.5em\hangafter=1\noindent Monarchie. \lxlinks{Dominatus unius; potentia singularis.}\par}
\par{\dictfontsize\hangindent=1.5em\hangafter=1\noindent \textit{Directeur de la} monnaie. \lxlinks{Monetæ censor.}\par}
\par{\dictfontsize\hangindent=1.5em\hangafter=1\noindent Montre. \lxlinks{Horologium portatile, rotatum minus.}\par}
\par{\dictfontsize\hangindent=1.5em\hangafter=1\noindent \textit{Lois morales, de la} morale. \lxlinks{Res morum, morum præcepta.}\par}
\par{\dictfontsize\hangindent=1.5em\hangafter=1\noindent \textit{Celui qui est} mort \textit{civilement, qui a perdu ses droits de citoyen}; \textit{qui est frappé de mort civile}. \lxlinks{Ærarius, capite census.}\par}
\par{\dictfontsize\hangindent=1.5em\hangafter=1\noindent Mortier (\textit{artillerie}), \textit{pièce d’artillerie pour lancer des bombes}. \lxlinks{Mortarium ignitum, bellicum, tormentum excitando incendio, succutiendis ædibus.}\par}
\par{\dictfontsize\hangindent=0pt\noindent\textls[120]{\scshape n}\par}
\par{\dictfontsize\hangindent=1.5em\hangafter=1\noindent \textit{De} naissance \textit{honnête, de famille noble}. \lxlinks{Honesto loco natus.}\par}
\par{\dictfontsize\hangindent=1.5em\hangafter=1\noindent \textit{D’illustre} naissance. \lxlinks{Natus haud obscuro loco, summo loco.}\par}
\par{\dictfontsize\hangindent=1.5em\hangafter=1\noindent Naturaliser. \lxlinks{Civitate donare peregrinum.}\par}
\par{\dictfontsize\hangindent=1.5em\hangafter=1\noindent Négociateur \textit{de la paix}; médiateur \textit{de la paix}. \lxlinks{Conciliator pacis, disceptator de pace.}\par}
\par{\dictfontsize\hangindent=1.5em\hangafter=1\noindent Neutralité. \lxlinks{Animus a studio partium alienus, neutri parti addictus animus.}\par}
\par{\dictfontsize\hangindent=1.5em\hangafter=1\noindent Neutre. \lxlinks{Medius, neutrius partis, medium partibus se præstare.}\par}
\par{\dictfontsize\hangindent=1.5em\hangafter=1\noindent \textit{Vers le} nord, \textit{du côté du nord}. \lxlinks{In septentriones spectare.}\par}
\par{\dictfontsize\hangindent=1.5em\hangafter=1\noindent Notaire. \lxlinks{Qui præscriptoris Cæsarei jura adeptus est, obsignator testamentorum publica auctoritate præditus.}\par}
\par{\dictfontsize\hangindent=0pt\noindent\textls[120]{\scshape o}\par}
\par{\dictfontsize\hangindent=1.5em\hangafter=1\noindent Observer \textit{la marche des ennemis}. \lxlinks{Itinera hostium observare, explorare.}\par}
\par{\dictfontsize\hangindent=1.5em\hangafter=1\noindent Officiers \textit{supérieurs}. \lxlinks{Præfecti belli superiori honore, primorum ordinum ductores.}\par}
\par{\dictfontsize\hangindent=1.5em\hangafter=1\noindent \textit{Sous}-officiers. \lxlinks{Præfecti militum inferiori dignitate.}\par}
\par{\dictfontsize\hangindent=1.5em\hangafter=1\noindent Opéras. \lxlinks{Ludi scenici qui canendo peraguntur.}\par}
\par{\dictfontsize\hangindent=1.5em\hangafter=1\noindent Opuscule, \textit{brochure, petit traité}. \lxlinks{Libellus exiguus.}\par}
\par{\dictfontsize\hangindent=1.5em\hangafter=1\noindent \textit{Le Pater, l’}oraison dominicale. \lxlinks{Formula precandi a Christo instituta.}\par}
\par{\dictfontsize\hangindent=1.5em\hangafter=1\noindent Oratoire; \textit{chapelle}. \lxlinks{Ædicula, sacellum, sacrarium, locus privatæ pietati sepositus, dicatus.}\par}
\par{\dictfontsize\hangindent=1.5em\hangafter=1\noindent \textit{Décret impérial; ordonnance du roi}; ordonnance \textit{royale. Capitulaires de Charlemagne.} \lxlinks{Sanctio imperatoria, lex regia, lex in quam fidem suam obstringunt Cæsares.}\par}
\par{\dictfontsize\hangindent=1.5em\hangafter=1\noindent Ordre, \textit{commandement}. \lxlinks{Imperium, imperata, præscriptum.}\par}
\par{\dictfontsize\hangindent=1.5em\hangafter=1\noindent \textit{Agir d’après les} ordres, \textit{le commandement}. \lxlinks{Agere ad imperium.}\par}
\par{\dictfontsize\hangindent=1.5em\hangafter=1\noindent \textit{Donner des} ordres \textit{secrets, des instructions secrètes aux officiers supérieurs}. \lxlinks{Tabellas signatas præfectis dare.}\par}
\par{\dictfontsize\hangindent=1.5em\hangafter=1\noindent \textit{Mot d’}ordre; \textit{mot de ralliement}. \lxlinks{Tessera militaris.}\par}
\par{\dictfontsize\hangindent=1.5em\hangafter=1\noindent \textit{Chercher le mot d’}ordre. \lxlinks{Tesseram vigiliarum petere.}\par}
\par{\dictfontsize\hangindent=1.5em\hangafter=1\noindent \textit{Donner le mot d’}ordre. \lxlinks{Signum edere. Tesseram vigiliarum promulgare.}\par}
\par{\dictfontsize\hangindent=1.5em\hangafter=1\noindent \textit{Obéir aux} ordres, \textit{au commandement}. \lxlinks{Imperio parere.}\par}
\par{\dictfontsize\hangindent=1.5em\hangafter=1\noindent \textit{Sans l’}ordre, \textit{contre les ordres du général}. \lxlinks{Injussu ducis.}\par}
\par{\dictfontsize\hangindent=1.5em\hangafter=1\noindent Orgue \textit{portatif}. \lxlinks{Organum pneumaticum duobus follibus instructum, quod in locum quemcumque transferri potest.}\par}
\par{\dictfontsize\hangindent=1.5em\hangafter=1\noindent Ouvrages \textit{à corne} (\textit{fortific.}). \lxlinks{Cornuta opera.}\par}
\par{\dictfontsize\hangindent=1.5em\hangafter=1\noindent Ouvrages \textit{extérieurs} (\textit{m.}), \textit{dehors} (\textit{m. plur.}), \textit{pièces détachées pour les fortifications}; \textit{forts avancés}. \lxlinks{Munimenta exteriora, quæ manu adjuncta sunt, opera munitioni præstructa, moles insularum specie a vallo divulsæ et aqua circumfusæ.}\par}
\par{\dictfontsize\hangindent=0pt\noindent\textls[120]{\scshape p}\par}
\par{\dictfontsize\hangindent=1.5em\hangafter=1\noindent Pacha (\textit{gouverneur d’une province turque}). \lxlinks{Purpuratus v. g. Bosnensis. Administrator provinciæ apud Turcas.}\par}
\par{\dictfontsize\hangindent=1.5em\hangafter=1\noindent Pain \textit{de munition}. \lxlinks{Panis castrensis, secundarius.}\par}
\par{\dictfontsize\hangindent=1.5em\hangafter=1\noindent Pairs. \lxlinks{Proceres regni, qui jure suffragandi in regno gaudent.}\par}
\par{\dictfontsize\hangindent=1.5em\hangafter=1\noindent Palissades. \lxlinks{Valli, sudes munimenti præpilatæ, cervi. Sudium sepimentum. Palorum vallum.}\par}
\par{\dictfontsize\hangindent=1.5em\hangafter=1\noindent Palissades \textit{de pieux pointus} (\textit{t. de fortif.}). \lxlinks{Furcati stipites ad hostis ascensum prohibendum, conferti sudium pectines.}\par}
\par{\dictfontsize\hangindent=1.5em\hangafter=1\noindent \textit{Couper, préparer des pieux pour faire une} palissade. \lxlinks{Comparare.}\par}
\par{\dictfontsize\hangindent=1.5em\hangafter=1\noindent \textit{Faire une} palissade; \textit{planter des pieux pour faire une palissade}. \lxlinks{Sudes, stipites præacutos defigere.}\par}
\par{\dictfontsize\hangindent=1.5em\hangafter=1\noindent Pape; \textit{souverain-pontife}; \textit{saint-père}. \lxlinks{Summus pontifex; papa. Romanæ Ecclesiæ supremus pastor.}\par}
\par{\dictfontsize\hangindent=1.5em\hangafter=1\noindent \textit{Recevoir la bénédiction du} Pape. \lxlinks{A pontifice solemni precatione lustrari.}\par}
\par{\dictfontsize\hangindent=1.5em\hangafter=1\noindent \textit{Recevoir un Bref du} pape, \textit{du souverain-pontife}. \lxlinks{Litteræ a Pontifice Romano publice scriptæ.}\par}
\par{\dictfontsize\hangindent=1.5em\hangafter=1\noindent Papiers \textit{sans valeur}. \lxlinks{Scriptum nullius pretii.}\par}
\par{\dictfontsize\hangindent=1.5em\hangafter=1\noindent \textit{Un} paragraphe, \textit{un article, un chapitre démesurément long}. \lxlinks{Nimium quam longa verborum continuatio.}\par}
\par{\dictfontsize\hangindent=1.5em\hangafter=1\noindent Parapet. \lxlinks{Lorica, lorica post circitorum viam.}\par}
\par{\dictfontsize\hangindent=1.5em\hangafter=1\noindent Parapet. \lxlinks{Collis jaculatorius, mœnibus superstructus, colles interpositi in extrema mœnium planitie, valla prominentia. Tumuli propugnaculo impositi.}\par}
\par{\dictfontsize\hangindent=1.5em\hangafter=1\noindent \textit{Élever des} parapets \textit{sur les remparts}; \textit{garnir les remparts de parapets}. \lxlinks{Pluteos vallo addere.}\par}
\par{\dictfontsize\hangindent=1.5em\hangafter=1\noindent Parer, \textit{détourner un} coup; \textit{éviter un coup en se détournant}. \lxlinks{Declinatione quadam alicujus petitionem effugere.}\par}
\par{\dictfontsize\hangindent=1.5em\hangafter=1\noindent \textit{Actes du} parlement \textit{en Angleterre}. \lxlinks{Tabulæ publicæ supremæ curiæ in Anglia.}\par}
\par{\dictfontsize\hangindent=1.5em\hangafter=1\noindent \textit{Héraut}; \textit{envoyé} parlementaire (\textit{qui porte un caducée}). \lxlinks{Caduceator.}\par}
\par{\dictfontsize\hangindent=1.5em\hangafter=1\noindent \textit{Messager de paix.}—\textit{Héraut} parlementaire (\textit{qui porte un caducée}). \lxlinks{Caduceator, pacis internuntius.}\par}
\par{\dictfontsize\hangindent=1.5em\hangafter=1\noindent Parquet; \textit{plancher parqueté}. \lxlinks{Pavimentum tesselatum.}\par}
\par{\dictfontsize\hangindent=1.5em\hangafter=1\noindent Parrain. \lxlinks{Baptismatis testis, pater lustricus, sponsor fidei, sacræ initiationis arbiter.}\par}
\par{\dictfontsize\hangindent=1.5em\hangafter=1\noindent \textit{Choisir un} parrain. \lxlinks{Adhibere aliquem sponsorem fidei.}\par}
\par{\dictfontsize\hangindent=1.5em\hangafter=1\noindent \textit{Servir de} parrain \textit{à un enfant}; \textit{tenir un enfant sur les fonts de baptême}. \lxlinks{Ad sacrum fontem sponsor fui.}\par}
\par{\dictfontsize\hangindent=1.5em\hangafter=1\noindent Pasquinade; \textit{libelle}. \lxlinks{Libellus famosus; criminationibus confertus liber.}\par}
\par{\dictfontsize\hangindent=1.5em\hangafter=1\noindent \textit{Défendre le} passage \textit{de.... à qqn}. \lxlinks{Itinere prohibere aliquem.}\par}
\par{\dictfontsize\hangindent=1.5em\hangafter=1\noindent \textit{Délivrer un} passeport \textit{à quelqu’un}. \lxlinks{Commeatu, publicis litteris instruere aliquem.}\par}
\par{\dictfontsize\hangindent=1.5em\hangafter=1\noindent \textit{Refuser un} passeport \textit{à quelqu’un}. \lxlinks{Prohibere commeatu, aliquem intercludere.}\par}
\par{\dictfontsize\hangindent=1.5em\hangafter=1\noindent Passion, \textit{affection, inclination}. \lxlinks{Prava animi affectio.}\par}
\par{\dictfontsize\hangindent=1.5em\hangafter=1\noindent \textit{La} passion \textit{de N.-S.-J.-C.} \lxlinks{Historia Christi patientis.}\par}
\par{\dictfontsize\hangindent=1.5em\hangafter=1\noindent \textit{Réciter le} pater \textit{pour qu’un autre le répète mot pour mot}. \lxlinks{Præire verba orationis Dominicæ alicui.}\par}
\par{\dictfontsize\hangindent=1.5em\hangafter=1\noindent Patrouille; \textit{guet}. \lxlinks{Excubitores, circitores.}\par}
\par{\dictfontsize\hangindent=1.5em\hangafter=1\noindent \textit{Piquet}; \textit{garde du camp}, patrouille \textit{à cheval}. \lxlinks{Milites, equites pro castris excubantes.}\par}
\par{\dictfontsize\hangindent=1.5em\hangafter=1\noindent Patrouiller, \textit{faire la patrouille}. \lxlinks{Vigilias intentiori cura lustrare.}\par}
\par{\dictfontsize\hangindent=1.5em\hangafter=1\noindent Pause. \lxlinks{Respiratio in arte musica.}\par}
\par{\dictfontsize\hangindent=1.5em\hangafter=1\noindent \textit{Presser le} payement. \lxlinks{Consectari debitum.}\par}
\par{\dictfontsize\hangindent=1.5em\hangafter=1\noindent Pélerinage. \lxlinks{Peregrinatio religionis causa.}\par}
\par{\dictfontsize\hangindent=1.5em\hangafter=1\noindent \textit{Le jour de} pénitence \textit{et de} prières publiques \textit{est prescrit}. \lxlinks{Solemnes, publicæ supplicationes ad deposcendam superum pacem indictæ sunt.}\par}
\par{\dictfontsize\hangindent=1.5em\hangafter=1\noindent Pensionnaire. \lxlinks{A pensionibus reipublicæ.}\par}
\par{\dictfontsize\hangindent=1.5em\hangafter=1\noindent \textit{Grand}-pensionnaire, \textit{c. à d.}, \textit{principal chargé d’affaires en Hollande}. \lxlinks{Procurator reipublicæ in Hollandia.}\par}
\par{\dictfontsize\hangindent=1.5em\hangafter=1\noindent Percepteur, \textit{receveur}. \lxlinks{Coactor, exactor.}\par}
\par{\dictfontsize\hangindent=1.5em\hangafter=1\noindent Période (\textit{litt.}). \lxlinks{Ambitus verborum, circuitus orationis, continuatio verborum.}\par}
\par{\dictfontsize\hangindent=1.5em\hangafter=1\noindent \textit{Appliquer des} pétards \textit{contre une porte}; \textit{visser des pétards à une porte} (\textit{guer.}). \lxlinks{Pylocastrum admovere, applicare, injicere portæ.}\par}
\par{\dictfontsize\hangindent=1.5em\hangafter=1\noindent \textit{Enfoncer, entamer une} phalange; \textit{rompre les lignes ennemies}. \lxlinks{Perfringere, disjicere phalangem.}\par}
\par{\dictfontsize\hangindent=1.5em\hangafter=1\noindent \textit{Arrière-bouche}; pharynx, \textit{gosier}; \textit{gorge, cou}; \textit{gueule}. \lxlinks{Gula.}\par}
\par{\dictfontsize\hangindent=1.5em\hangafter=1\noindent Phénomène, \textit{apparition céleste}. \lxlinks{Apparitio sideris.}\par}
\par{\dictfontsize\hangindent=1.5em\hangafter=1\noindent Pièces de campagne. \lxlinks{Tormentum castrense, tormenta legionum.}\par}
\par{\dictfontsize\hangindent=1.5em\hangafter=1\noindent Piété (\textit{filiale}); \textit{affection pour ses proches}; \textit{respect} (\textit{pour ses maîtres}); \textit{amour} (\textit{de la patrie}). \lxlinks{Pietas.}\par}
\par{\dictfontsize\hangindent=1.5em\hangafter=1\noindent \textit{Respect}; piété. \lxlinks{Religio.}\par}
\par{\dictfontsize\hangindent=1.5em\hangafter=1\noindent \textit{Livrer une ville au} pillage; \textit{mettre une ville à sac}. \lxlinks{Vitam fortunasque civium direptioni militum objicere.}\par}
\par{\dictfontsize\hangindent=1.5em\hangafter=1\noindent \textit{Demi}-pique; \textit{pique de vélite} (\textit{infanterie légère chez les Romains}). \lxlinks{Velitaris.}\par}
\par{\dictfontsize\hangindent=1.5em\hangafter=1\noindent \textit{Être de} piquet \textit{au fourrage};—\textit{mettre de piquet au fourrage}. \lxlinks{Populares tueri.}\par}
\par{\dictfontsize\hangindent=1.5em\hangafter=1\noindent Pistolet,—\textit{revolver}. \lxlinks{Bombarda minor, equestris; manualis fistula.}\par}
\par{\dictfontsize\hangindent=1.5em\hangafter=1\noindent Place \textit{frontière}. \lxlinks{Urbs quæ fines tuetur, propugnaculum imperii defendendis limitibus appositum.}\par}
\par{\dictfontsize\hangindent=1.5em\hangafter=1\noindent Plaider; \textit{plaider une cause}. \lxlinks{In foro versari, causam agere, causæ alicujus patronum exsistere.}\par}
\par{\dictfontsize\hangindent=1.5em\hangafter=1\noindent \textit{Profession d’avocat}; \textit{assistance d’un défenseur}; plaidoirie. \lxlinks{Opera forensis, commentationes causarum.}\par}
\par{\dictfontsize\hangindent=1.5em\hangafter=1\noindent Plénipotentiaire (\textit{d’un prince, d’une nation}). \lxlinks{Rerum agendarum summa potestate instructus.}\par}
\par{\dictfontsize\hangindent=1.5em\hangafter=1\noindent \textit{Celui qui fait jouer les machines de guerre}, (\textit{de nos jours on dirait} pointeur \textit{d’artillerie}) \textit{ou d’une manière plus générale: officier du génie et d’artillerie}. \lxlinks{Is qui tormenta intendit, machinarum bellicarum librator.}\par}
\par{\dictfontsize\hangindent=1.5em\hangafter=1\noindent \textit{Un tonneau, un baril de} poix. \lxlinks{Cupa tæda et pice conferta.}\par}
\par{\dictfontsize\hangindent=1.5em\hangafter=1\noindent \textit{Règlement de} police; \textit{loi somptuaire} (\textit{chez les Romains}). \lxlinks{Lex sumptuaria, ordo disciplinæ civilis.}\par}
\par{\dictfontsize\hangindent=1.5em\hangafter=1\noindent Politique. \lxlinks{Prudentia civilis.}\par}
\par{\dictfontsize\hangindent=1.5em\hangafter=1\noindent \textit{Excellent} politique. \lxlinks{Pacis artium gnarus.}\par}
\par{\dictfontsize\hangindent=1.5em\hangafter=1\noindent Pompe \textit{à feu, à incendie}. \lxlinks{Machina delendis ædium incendiis deserviens, deputata.}\par}
\par{\dictfontsize\hangindent=1.5em\hangafter=1\noindent Ponton. \lxlinks{Rates junctæ trajectui aptæ, ad transeundum instructæ.}\par}
\par{\dictfontsize\hangindent=1.5em\hangafter=1\noindent Porcelaine. \lxlinks{Vasa fictilia Sinensium.}\par}
\par{\dictfontsize\hangindent=1.5em\hangafter=1\noindent Portier, \textit{concierge}. \lxlinks{Atriensis.}\par}
\par{\dictfontsize\hangindent=1.5em\hangafter=1\noindent Portraits \textit{des aïeux}; \textit{portraits de famille}. \lxlinks{Patrum et avorum picta effigies.}\par}
\par{\dictfontsize\hangindent=1.5em\hangafter=1\noindent \textit{Reconnaître la} position, \textit{la situation d’une ville}. \lxlinks{Situm urbis perspicere.}\par}
\par{\dictfontsize\hangindent=1.5em\hangafter=1\noindent \textit{Ne pas maintenir une} position; \textit{lâcher pied}. \lxlinks{Locum non tenere.}\par}
\par{\dictfontsize\hangindent=1.5em\hangafter=1\noindent \textit{Prendre} position \textit{sur une hauteur, une colline}. \lxlinks{Capere collem.}\par}
\par{\dictfontsize\hangindent=1.5em\hangafter=1\noindent Possesseurs \textit{d’un bien héréditaire}. \lxlinks{Possessores prædii hereditarii.}\par}
\par{\dictfontsize\hangindent=1.5em\hangafter=1\noindent \textit{Obtenir qqche à titre de} possession \textit{absolue}. \lxlinks{Aliquid titulo supremi dominii, neminique obnoxio jure consequi.}\par}
\par{\dictfontsize\hangindent=1.5em\hangafter=1\noindent \textit{Chevaux de} poste. \lxlinks{Equi publici, meritorii, cursorii, veredarii.}\par}
\par{\dictfontsize\hangindent=1.5em\hangafter=1\noindent \textit{Grand-maître des} postes. \lxlinks{Cursuum publicorum supremus moderator.}\par}
\par{\dictfontsize\hangindent=1.5em\hangafter=1\noindent \textit{Environner l’ennemi d’une ligne de} postes, \textit{et l’assiéger en qqe sorte}. \lxlinks{Cingere hostem stationibus in modum obsidii.}\par}
\par{\dictfontsize\hangindent=1.5em\hangafter=1\noindent Post-scriptum. \lxlinks{Transversus versiculus extremæ epistolæ.}\par}
\par{\dictfontsize\hangindent=1.5em\hangafter=1\noindent Poudre. \lxlinks{Pulvis pyrius, nitratus, tormentarius.}\par}
\par{\dictfontsize\hangindent=1.5em\hangafter=1\noindent Praticable, \textit{faisable, possible}. \lxlinks{Quod fieri potest, in usum deduci potest.}\par}
\par{\dictfontsize\hangindent=1.5em\hangafter=1\noindent Prébende; \textit{bénéfice}. \lxlinks{Reditus et census ex bonis sacris.}\par}
\par{\dictfontsize\hangindent=1.5em\hangafter=1\noindent \textit{Gouverneur}, précepteur \textit{de jeunes gens de qualité, de bonne famille}. \lxlinks{Moderator juventutis nobilis, qui liberalibus juventutis studiis dux et moderator præest.}\par}
\par{\dictfontsize\hangindent=1.5em\hangafter=1\noindent Prêcher; \textit{faire un sermon}; \textit{annoncer la parole de Dieu}. \lxlinks{De rebus divinis ad populum concionem dicere, publice interpretari res divinas.}\par}
\par{\dictfontsize\hangindent=1.5em\hangafter=1\noindent Prédicateur. \lxlinks{Præco verbi divini.}\par}
\par{\dictfontsize\hangindent=1.5em\hangafter=1\noindent Préliminaires \textit{de paix, bases d’un traité de paix}. \lxlinks{Prima lineamenta pacis. Rudis pacis typus.}\par}
\par{\dictfontsize\hangindent=1.5em\hangafter=1\noindent \textit{Se mettre en devoir de faire le siége} (\textit{d’une ville}), \textit{faire les} préparatifs \textit{du siége}. \lxlinks{Quæ ad oppugnationem oppidi pertinent comportare.}\par}
\par{\dictfontsize\hangindent=1.5em\hangafter=1\noindent Prescription; \textit{surannation} (\textit{de}), \textit{usucapion} (\textit{t. de dr. rom.}). \lxlinks{Acquisitio seu adjectio dominii per continuationem temporis lege definiti. Usucapio.}\par}
\par{\dictfontsize\hangindent=1.5em\hangafter=1\noindent \textit{Dispute sur la} préséance. \lxlinks{Certamen honoris.}\par}
\par{\dictfontsize\hangindent=1.5em\hangafter=1\noindent \textit{Être en dispute, disputer sur la} préséance. \lxlinks{In honoris contentionem incidere.}\par}
\par{\dictfontsize\hangindent=1.5em\hangafter=1\noindent Président \textit{du conseil privé}. \lxlinks{Senatus sanctioris princeps.}\par}
\par{\dictfontsize\hangindent=1.5em\hangafter=1\noindent \textit{Être} président, \textit{avoir la présidence de la chambre des députés}. \lxlinks{Tota comitia suis humeris, curis sustinere.}\par}
\par{\dictfontsize\hangindent=1.5em\hangafter=1\noindent Prétendant. \lxlinks{Æmulus regni.}\par}
\par{\dictfontsize\hangindent=1.5em\hangafter=1\noindent \textit{Élever des} prétentions \textit{sur d’anciens droits}. \lxlinks{Jura antiqua revocare.}\par}
\par{\dictfontsize\hangindent=1.5em\hangafter=1\noindent Prêter \textit{de l’argent à intérêts}; \textit{placer de l’argent}. \lxlinks{Fenori pecuniam dare.}\par}
\par{\dictfontsize\hangindent=1.5em\hangafter=1\noindent Prêter, \textit{avancer de l’argent à qqn}. \lxlinks{Pecunias expensas alicui ferre.}\par}
\par{\dictfontsize\hangindent=1.5em\hangafter=1\noindent \textit{Grand}-prévôt \textit{de l’armée}. \lxlinks{Penes quem est summa imperii custodia.}\par}
\par{\dictfontsize\hangindent=1.5em\hangafter=1\noindent Primipile \textit{ou primipilaire}, \textit{c. à d.}, \textit{1ᵉʳ centurion}; \textit{celui qui commande la 1ᵉʳᵉ compagnie des triaires. Premier capitaine}. \lxlinks{Primi pili centurio.}\par}
\par{\dictfontsize\hangindent=1.5em\hangafter=1\noindent Prince \textit{royal, impérial}. \lxlinks{Heres regni, in spem regni genitus.}\par}
\par{\dictfontsize\hangindent=1.5em\hangafter=1\noindent \textit{Droits de justice}; \textit{frais de} procédure; \textit{dépens}. \lxlinks{Lis æstimata, expensæ litis.}\par}
\par{\dictfontsize\hangindent=1.5em\hangafter=1\noindent Procession. \lxlinks{Solemnis sacrorum pompa, supplicantium agmina.}\par}
\par{\dictfontsize\hangindent=1.5em\hangafter=1\noindent \textit{Aller en} procession. \lxlinks{Deducere pie supplices.}\par}
\par{\dictfontsize\hangindent=1.5em\hangafter=1\noindent Proclamer \textit{par des hérauts envoyés dans différentes directions}. \lxlinks{Præconibus circummissis pronuntiare.}\par}
\par{\dictfontsize\hangindent=1.5em\hangafter=1\noindent \textit{Faire} profession; \textit{faire, émettre ses vœux}; \textit{se consacrer à DIEU par les vœux}. \lxlinks{Vota solemnia sacri ordinis nuncupare, DEO se solemni votorum sponsione devovere.}\par}
\par{\dictfontsize\hangindent=1.5em\hangafter=1\noindent \textit{Au} prorata (\textit{dr.}), \textit{à proportion, proportionnellement}. \lxlinks{Pro rata parte, pro cujusque facultatibus.}\par}
\par{\dictfontsize\hangindent=1.5em\hangafter=1\noindent \textit{Supprimer, enlever, refuser le droit de} protestation. \lxlinks{Intercessionem tollere.}\par}
\par{\dictfontsize\hangindent=1.5em\hangafter=1\noindent Protester. \lxlinks{Intercedere, vetare, obnuntiare, contra obtestari.}\par}
\par{\dictfontsize\hangindent=1.5em\hangafter=1\noindent Protonotaire. \lxlinks{Primarius scriba, primus eorum qui reipublicæ sunt ab epistolis.}\par}
\par{\dictfontsize\hangindent=1.5em\hangafter=1\noindent Protonotariat. \lxlinks{Primarii scribæ munus.}\par}
\par{\dictfontsize\hangindent=1.5em\hangafter=1\noindent Puits \textit{artésien}; \textit{jet d’eau}. \lxlinks{Aquæ artificiose prosilientes.}\par}
\par{\dictfontsize\hangindent=0pt\noindent\textls[120]{\scshape q}\par}
\par{\dictfontsize\hangindent=1.5em\hangafter=1\noindent Quaker. \lxlinks{Secta tremulantium.}\par}
\par{\dictfontsize\hangindent=1.5em\hangafter=1\noindent Qualités. \lxlinks{Bona quæ a natura data sunt, ingenii bona, dotes animi.}\par}
\par{\dictfontsize\hangindent=1.5em\hangafter=1\noindent Quarantaine. \lxlinks{Mora quadragesimum in diem producta, tempus valetudini experiundæ præfixum.}\par}
\par{\dictfontsize\hangindent=1.5em\hangafter=1\noindent \textit{Loger l’armée, mettre l’armée en} quartier \textit{dans les villages}. \lxlinks{Exercitum in agris collocare.}\par}
\par{\dictfontsize\hangindent=1.5em\hangafter=1\noindent \textit{Quartiers}; quartiers \textit{d’hiver}; \textit{campement}; \textit{cantonnements}. \lxlinks{Stationes militum; hiberna, stativa.}\par}
\par{\dictfontsize\hangindent=1.5em\hangafter=1\noindent \textit{Assigner les} quartiers \textit{d’hiver}. \lxlinks{Hiberna constituere.}\par}
\par{\dictfontsize\hangindent=1.5em\hangafter=1\noindent \textit{Quitter ses} quartiers \textit{d’hiver}; \textit{entrer en campagne}. \lxlinks{Copias ex hibernis excire.}\par}
\par{\dictfontsize\hangindent=1.5em\hangafter=1\noindent Quartier \textit{général du roi}. \lxlinks{Tabernaculum regis.}\par}
\par{\dictfontsize\hangindent=1.5em\hangafter=1\noindent Quartier \textit{général} (\textit{prétoire anciennement}). \lxlinks{Prætorium.}\par}
\par{\dictfontsize\hangindent=1.5em\hangafter=1\noindent Quête, (\textit{collecte ecclésiast.}). \lxlinks{Pecunia viritim collecta. Stipis voluntaria collatio.}\par}
\par{\dictfontsize\hangindent=1.5em\hangafter=1\noindent \textit{Argent quêté, produit d’une} quête. \lxlinks{Pecunia viritim, in capita collata.}\par}
\par{\dictfontsize\hangindent=1.5em\hangafter=1\noindent Quittance, \textit{acquit}; \textit{reçu}. \lxlinks{Syngrapha, chirographum, acceptilatio.}\par}
\par{\dictfontsize\hangindent=1.5em\hangafter=1\noindent Quittance \textit{générale, finale}. \lxlinks{Agnitio satisfactionis plenissima.}\par}
\par{\dictfontsize\hangindent=0pt\noindent\textls[120]{\scshape r}\par}
\par{\dictfontsize\hangindent=1.5em\hangafter=1\noindent Rançon. \lxlinks{Pretium captivorum.}\par}
\par{\dictfontsize\hangindent=1.5em\hangafter=1\noindent Rangs. \lxlinks{Ordines.}\par}
\par{\dictfontsize\hangindent=1.5em\hangafter=1\noindent \textit{Garder les} rangs. \lxlinks{Servare ordines.}\par}
\par{\dictfontsize\hangindent=1.5em\hangafter=1\noindent \textit{Rompre les} rangs. \lxlinks{Ordines laxare.}\par}
\par{\dictfontsize\hangindent=1.5em\hangafter=1\noindent Rangs \textit{serrés}; \textit{être en rangs serrés}. \lxlinks{Conserta agmina, densati ordines consistunt, inter aciem consistens miles.}\par}
\par{\dictfontsize\hangindent=1.5em\hangafter=1\noindent \textit{Lettre de} rappel; \textit{édit, décret en vertu duquel ceux qui occupent des postes militaires ou civils dans un pays étranger sont rappelés, pour motif de guerre}. \lxlinks{Edictum, decretum quo cives, milites sub extero Rege stipendia merentes domum revocantur.}\par}
\par{\dictfontsize\hangindent=1.5em\hangafter=1\noindent \textit{Référendaire}; rapporteur. \lxlinks{Cujus est in consilio publico causarum momenta exponere.}\par}
\par{\dictfontsize\hangindent=1.5em\hangafter=1\noindent Ratifier. \lxlinks{Ratum habere, auctoritate publica aliquid affirmare.}\par}
\par{\dictfontsize\hangindent=1.5em\hangafter=1\noindent Ration. \lxlinks{Demensum annonæ militaris.}\par}
\par{\dictfontsize\hangindent=1.5em\hangafter=1\noindent Ravelin; \textit{demi-lune}. \lxlinks{Portæ propugnaculum, munimentum extra mœnia.}\par}
\par{\dictfontsize\hangindent=1.5em\hangafter=1\noindent Récapituler; \textit{résumer}. \lxlinks{Summa rerum capita repetere.}\par}
\par{\dictfontsize\hangindent=1.5em\hangafter=1\noindent \textit{Assurer} réception, \textit{donner quittance de l’argent prêté à qqn}. \lxlinks{Pecuniam datam referre. Pecunias expensas alicui referre acceptas.}\par}
\par{\dictfontsize\hangindent=1.5em\hangafter=1\noindent Recette \textit{assurée}. \lxlinks{Reditus status.}\par}
\par{\dictfontsize\hangindent=1.5em\hangafter=1\noindent Recette \textit{fortuite}. \lxlinks{Reditus fortuitus.}\par}
\par{\dictfontsize\hangindent=1.5em\hangafter=1\noindent \textit{Le produit de la} recette \textit{couvre les dépenses}; \textit{les dépenses et les recettes se balancent}. \lxlinks{Ratio par est acceptorum et facti impendii.}\par}
\par{\dictfontsize\hangindent=1.5em\hangafter=1\noindent \textit{Porter en compte les} recettes. \lxlinks{Ratio accepti scribitur.}\par}
\par{\dictfontsize\hangindent=1.5em\hangafter=1\noindent \textit{Le produit de la} recette \textit{et les dépenses se balancent}. \lxlinks{Ratio par est acceptorum et datorum, ratio accepti et expensi congruit.}\par}
\par{\dictfontsize\hangindent=1.5em\hangafter=1\noindent Receveur \textit{des contributions, percepteur}. \lxlinks{Ærarii tribunus, curator census bonorum.}\par}
\par{\dictfontsize\hangindent=1.5em\hangafter=1\noindent Receveur \textit{de l’accise, des contributions indirectes}. \lxlinks{Curator census e victu persolvi soliti.}\par}
\par{\dictfontsize\hangindent=1.5em\hangafter=1\noindent \textit{Donner à quelqu’un une} récompense \textit{en argent}. \lxlinks{Præmia rei pecuniariæ alicui tribuere.}\par}
\par{\dictfontsize\hangindent=1.5em\hangafter=1\noindent Recruter, \textit{faire des levées}; \textit{enrôler}. \lxlinks{Exercitui, regionibus supplementum scribere.}\par}
\par{\dictfontsize\hangindent=1.5em\hangafter=1\noindent Recueil \textit{de notes, mémoire, mémorial}; \textit{livre de raison ou grand livre}, \textit{c. à d.}, \textit{un registre où les négociants portent tout leur compte par doit et avoir}. \lxlinks{Electorum commentarius.}\par}
\par{\dictfontsize\hangindent=1.5em\hangafter=1\noindent \textit{Revenir aux premières conditions de} reddition. \lxlinks{Ad easdem deditionis conditiones recurrere.}\par}
\par{\dictfontsize\hangindent=1.5em\hangafter=1\noindent \textit{Élever, construire des} redoutes. \lxlinks{Ad extremas fossas castella constituere.}\par}
\par{\dictfontsize\hangindent=1.5em\hangafter=1\noindent \textit{Se charger d’un rapport}; \textit{rapporter}, référer (\textit{t. de palais}). \lxlinks{Aliquid ad principis arbitrium referre.}\par}
\par{\dictfontsize\hangindent=1.5em\hangafter=1\noindent Réfléchir. \lxlinks{Rei alicujus rationem ducere.}\par}
\par{\dictfontsize\hangindent=1.5em\hangafter=1\noindent Réfuter, \textit{repousser, une accusation calomnieuse, en triompher}; \textit{convaincre de crime l’accusateur par ses propres arguments}. \lxlinks{Refellere crimen et rejicere in suum auctorem, refutare injuriam, retundere conatus calumniæ.}\par}
\par{\dictfontsize\hangindent=1.5em\hangafter=1\noindent Régale; \textit{droit régalien}. \lxlinks{Jus regi sacrum.}\par}
\par{\dictfontsize\hangindent=1.5em\hangafter=1\noindent Régent; \textit{lieutenant-général de l’empire, du royaume}; \textit{vicaire de l’empire}; \textit{interroi}. \lxlinks{Interrex.}\par}
\par{\dictfontsize\hangindent=1.5em\hangafter=1\noindent Registres \textit{publics, documents} (\textit{actes, titres, pièces justificatives}). \lxlinks{Auctoritates publicæ.}\par}
\par{\dictfontsize\hangindent=1.5em\hangafter=1\noindent Règle. \lxlinks{Regula.}\par}
\par{\dictfontsize\hangindent=1.5em\hangafter=1\noindent Religieuse. \lxlinks{Virgo DEO sacrata, quæ integrum virginitatis florem DEO solemni votorum sponsione dicavit.}\par}
\par{\dictfontsize\hangindent=1.5em\hangafter=1\noindent \textit{Jour de} repos. \lxlinks{Dies ad quietem datus.}\par}
\par{\dictfontsize\hangindent=1.5em\hangafter=1\noindent Représailles. \lxlinks{Jus talionis, pignorationes.}\par}
\par{\dictfontsize\hangindent=1.5em\hangafter=1\noindent \textit{Faire des} réquisitions \textit{de vivres, de munitions de bouche}. \lxlinks{Frumenti vecturas finitimis civitatibus describere.}\par}
\par{\dictfontsize\hangindent=1.5em\hangafter=1\noindent Résidence \textit{royale}. \lxlinks{Regia.}\par}
\par{\dictfontsize\hangindent=1.5em\hangafter=1\noindent Respect; \textit{vénération, déférence}; \textit{témoignage de respect}; \textit{hommage}. \lxlinks{Reverentia, veneratio.}\par}
\par{\dictfontsize\hangindent=1.5em\hangafter=1\noindent Résumé \textit{d’un compte; compte par bref état}. \lxlinks{Breviarium rationum.}\par}
\par{\dictfontsize\hangindent=1.5em\hangafter=1\noindent \textit{Ligne de circonvallation, ouvrages}, retranchements, \textit{lignes de défense en avant du camp}. \lxlinks{Limes externus, circuitus castrorum. Pro portis castrorum.}\par}
\par{\dictfontsize\hangindent=1.5em\hangafter=1\noindent Retranchement, \textit{tranchées, ouvrages en terre, palissades pour protéger un camp, une armée}. \lxlinks{Aggeres castrenses.}\par}
\par{\dictfontsize\hangindent=1.5em\hangafter=1\noindent \textit{Couvrir, protéger par des} retranchements \textit{le front de son armée}. \lxlinks{A fronte contra hostem fossam ducere.}\par}
\par{\dictfontsize\hangindent=1.5em\hangafter=1\noindent \textit{Collationner}, revoir \textit{un écrit, un ouvrage, un opuscule}. \lxlinks{Recognoscere exemplum, libellos.}\par}
\par{\dictfontsize\hangindent=1.5em\hangafter=1\noindent Revue, \textit{inspection, révision}. \lxlinks{Delectus militum, lustratio exercitus.}\par}
\par{\dictfontsize\hangindent=1.5em\hangafter=1\noindent \textit{Faire passer la} revue \textit{à l’armée}; \textit{passer la revue des troupes}. \lxlinks{Numerum copiarum inire, censum agere militum, lustrare, recensere exercitum.}\par}
\par{\dictfontsize\hangindent=1.5em\hangafter=1\noindent \textit{Inspecteurs}; \textit{intendants qui passent la} revue \textit{des troupes}; \textit{officiers de révision}. \lxlinks{Censores militum.}\par}
\par{\dictfontsize\hangindent=1.5em\hangafter=1\noindent \textit{Rôle}; \textit{état des hommes présents à la} revue. \lxlinks{Album militare.}\par}
\par{\dictfontsize\hangindent=1.5em\hangafter=1\noindent \textit{Ordo}; rituel. \lxlinks{Liber in quo continetur ratio Divini officii, liber in quo ratio in dispensandis salutis mysteriis tenenda proscribitur.}\par}
\par{\dictfontsize\hangindent=1.5em\hangafter=1\noindent \textit{Ordo, directoire}, rituel. \lxlinks{Liber in quo continetur ratio divini officii et salutis mysteriorum dispensandorum.}\par}
\par{\dictfontsize\hangindent=1.5em\hangafter=1\noindent Rote; \textit{sacrée congrégation}, (\textit{tribunal de la Rote}). \lxlinks{Judices sacrarum cognitionum. Sacri cognitores Romæ.}\par}
\par{\dictfontsize\hangindent=1.5em\hangafter=1\noindent \textit{Voie}, route \textit{militaire}. \lxlinks{Via militaris.}\par}
\par{\dictfontsize\hangindent=1.5em\hangafter=1\noindent \textit{Abandonner la voie, la} route \textit{militaire}. \lxlinks{Via militari decedere.}\par}
\par{\dictfontsize\hangindent=0pt\noindent\textls[120]{\scshape s}\par}
\par{\dictfontsize\hangindent=1.5em\hangafter=1\noindent \textit{Recevoir, fréquenter les} Sacrements. \lxlinks{Sacra mysteria, salutis mysteria obire.}\par}
\par{\dictfontsize\hangindent=1.5em\hangafter=1\noindent \textit{Sacristie.} \lxlinks{Aditum, locus in quo apparatus sacer sumitur, in quo occulta et recondita templi servantur.}\par}
\par{\dictfontsize\hangindent=1.5em\hangafter=1\noindent \textit{Saisir qqche}; \textit{faire la} saisie \textit{de qqche}; \textit{mettre arrêt sur qqche}; \textit{séquestrer}; \textit{mettre à la séquestration, confisquer}. \lxlinks{Manu prætoria alicujus bona detinere.}\par}
\par{\dictfontsize\hangindent=1.5em\hangafter=1\noindent Salaire; \textit{gages}; \textit{appointements, traitement}; \textit{solde} (\textit{milit.}). \lxlinks{Annua merces, stipendium.}\par}
\par{\dictfontsize\hangindent=1.5em\hangafter=1\noindent Sauf-conduit. \lxlinks{Litteræ tutum iter præstantes, litteræ auctoritatis publicæ.}\par}
\par{\dictfontsize\hangindent=1.5em\hangafter=1\noindent Sauf-conduit. \lxlinks{Fides publica.}\par}
\par{\dictfontsize\hangindent=1.5em\hangafter=1\noindent \textit{Donner un} sauf-conduit \textit{à qqn}. \lxlinks{Præstare fidem publicam.}\par}
\par{\dictfontsize\hangindent=1.5em\hangafter=1\noindent Scander \textit{un vers}. \lxlinks{Versum, carmen metiri.}\par}
\par{\dictfontsize\hangindent=1.5em\hangafter=1\noindent Sceau \textit{de la ville}. \lxlinks{Signum publicum.}\par}
\par{\dictfontsize\hangindent=1.5em\hangafter=1\noindent Séjour (\textit{t. milit.}). \lxlinks{Locus ad militum requiem.}\par}
\par{\dictfontsize\hangindent=1.5em\hangafter=1\noindent Selle. \lxlinks{Sella equestris.}\par}
\par{\dictfontsize\hangindent=1.5em\hangafter=1\noindent Sénat (\textit{France}); \textit{chambre haute, chambre des Lords} (\textit{Anglet.}); \textit{chambre des seigneurs} (\textit{Allemagne}). \lxlinks{Ordo senatorius, senatus amplissimus, patrum sanctum concilium.}\par}
\par{\dictfontsize\hangindent=1.5em\hangafter=1\noindent Sénat, \textit{ordre du sénat, conseil, les membres du conseil}. \lxlinks{Ordo senatorius.}\par}
\par{\dictfontsize\hangindent=1.5em\hangafter=1\noindent Sentinelle \textit{perdue}. \lxlinks{Conclamatæ salutis extimæ vigiliæ. Procubitores, qui infestas maxime et hosti oppositas stationes occupant.}\par}
\par{\dictfontsize\hangindent=1.5em\hangafter=1\noindent \textit{Relever qqn de} sentinelle;—\textit{relever de poste un factionnaire}. \lxlinks{Vices stationum mutare.}\par}
\par{\dictfontsize\hangindent=1.5em\hangafter=1\noindent \textit{Saisir qqche}; \textit{faire la saisie de qqche}; \textit{mettre arrêt sur qqche}; séquestrer; \textit{mettre à la séquestration, confisquer}. \lxlinks{Manu prætoria alicujus bona detinere.}\par}
\par{\dictfontsize\hangindent=1.5em\hangafter=1\noindent Sergent-major. \lxlinks{Ordinum diribitor, instructor centuriæ.}\par}
\par{\dictfontsize\hangindent=1.5em\hangafter=1\noindent Serment \textit{de bannissement, de ne pas revenir et de ne pas se venger} (\textit{féod.}). \lxlinks{Cautio de inimicitiis ponendis, de non revertendo.}\par}
\par{\dictfontsize\hangindent=1.5em\hangafter=1\noindent \textit{Se purger d’une imputation, d’une accusation par} serment. \lxlinks{Jurejurando se purgare.}\par}
\par{\dictfontsize\hangindent=1.5em\hangafter=1\noindent \textit{Vaisselle de table}, service \textit{en or et en argent}. \lxlinks{Vasa convivalia ex auro et argento.}\par}
\par{\dictfontsize\hangindent=1.5em\hangafter=1\noindent \textit{Constituer, imposer une} servitude \textit{sur un bien}; \textit{grever de charges, d’impôts un bien}. \lxlinks{Imponere servitutem fundo.}\par}
\par{\dictfontsize\hangindent=1.5em\hangafter=1\noindent \textit{Il est d’un caractère si bien tempéré, qu’il sait unir la} sévérité \textit{à l’affabilité}. \lxlinks{Est ita temperatis moderatisque moribus, ut summa severitas cum summa humanitate conjungatur.}\par}
\par{\dictfontsize\hangindent=1.5em\hangafter=1\noindent Simonie. \lxlinks{Rerum sacrarum nefaria nundinatio.}\par}
\par{\dictfontsize\hangindent=1.5em\hangafter=1\noindent \textit{Être} situé. \lxlinks{Situm esse.}\par}
\par{\dictfontsize\hangindent=1.5em\hangafter=1\noindent Sommairement, \textit{en résumé, par manière d’extrait, compendieusement, en abrégé, en raccourci}. \lxlinks{Ex singulis capitibus, per capita, summatim, per compendia.}\par}
\par{\dictfontsize\hangindent=1.5em\hangafter=1\noindent \textit{Voir à l’aide de l’addition et de la} soustraction, \textit{ce qui reste}. \lxlinks{Addendo deducendoque experiri, quid reliqui sit.}\par}
\par{\dictfontsize\hangindent=1.5em\hangafter=1\noindent Spahis. \lxlinks{Turcicus equitatus.}\par}
\par{\dictfontsize\hangindent=1.5em\hangafter=1\noindent Substitut. \lxlinks{Qui alicui ad munus obeundum adjunctus est.}\par}
\par{\dictfontsize\hangindent=1.5em\hangafter=1\noindent \textit{Cicéron ne dut ses} succès \textit{qu’à lui-même}. \lxlinks{Cicero omnia incrementa sibi debuit.}\par}
\par{\dictfontsize\hangindent=1.5em\hangafter=1\noindent \textit{Pacte de confraternité}; \textit{de} succession \textit{réciproque}. \lxlinks{Pacta gentilitia.}\par}
\par{\dictfontsize\hangindent=1.5em\hangafter=1\noindent Succursale, (\textit{église}) \textit{annexe}. \lxlinks{Templum proprio Curione sacro destitutum, templum ejusdem Parochiæ.}\par}
\par{\dictfontsize\hangindent=1.5em\hangafter=1\noindent \textit{Vers le} sud, \textit{au sud, au midi}. \lxlinks{Ad meridiem.}\par}
\par{\dictfontsize\hangindent=1.5em\hangafter=1\noindent Suffragant. \lxlinks{Vicarius Episcopi.}\par}
\par{\dictfontsize\hangindent=1.5em\hangafter=1\noindent \textit{Lieu de} supplice; \textit{gibet}. \lxlinks{Locus publicus facinorosorum supplicio deputatus; locus quo hi qui ad supplicium sunt dati, securi percutiuntur.}\par}
\par{\dictfontsize\hangindent=1.5em\hangafter=1\noindent Supérieure (\textit{d’une simple congrégation}). \lxlinks{Rectrix Virginum DEO sacrarum, quæ collegio sacrarum Virginum \textit{vel} Virginibus DEO sacratis præest, antistes sacrarum Virginum.}\par}
\par{\dictfontsize\hangindent=1.5em\hangafter=1\noindent Surplis; \textit{aube}. \lxlinks{Vestis, tunica lintea.}\par}
\par{\dictfontsize\hangindent=1.5em\hangafter=1\noindent \textit{Donner un} sursis \textit{à qqn, accorder un délai à qqn}. \lxlinks{Statuere alicui diem laxiorem.}\par}
\par{\dictfontsize\hangindent=1.5em\hangafter=1\noindent \textit{Obtenir un} sursis, \textit{un délai d’ajournement}. \lxlinks{Potestatem ampliandi nancisci.}\par}
\par{\dictfontsize\hangindent=0pt\noindent\textls[120]{\scshape t}\par}
\par{\dictfontsize\hangindent=1.5em\hangafter=1\noindent \textit{Musique en tête}, tambour \textit{battant}. \lxlinks{Tympana resonantia, tympanorum sonus.}\par}
\par{\dictfontsize\hangindent=1.5em\hangafter=1\noindent Taxe; \textit{taux, prix}; \textit{estimation}. \lxlinks{Rei æstimatæ pretium.}\par}
\par{\dictfontsize\hangindent=1.5em\hangafter=1\noindent \textit{Chanter le} Te Deum. \lxlinks{Hymnum Ambrosianum cantare.}\par}
\par{\dictfontsize\hangindent=1.5em\hangafter=1\noindent Teneur \textit{de livres}. \lxlinks{Actor summarum, rationum magister.}\par}
\par{\dictfontsize\hangindent=1.5em\hangafter=1\noindent \textit{Second}, témoin (\textit{d’un duel}). \lxlinks{Arbiter singularis certaminis.}\par}
\par{\dictfontsize\hangindent=1.5em\hangafter=1\noindent \textit{Prendre son} temps (\textit{familier}). \lxlinks{Tempus aucupari.}\par}
\par{\dictfontsize\hangindent=1.5em\hangafter=1\noindent \textit{Dans les} termes \textit{mêmes prescrits par Nevius}. \lxlinks{In ea ipsa verba quæ Nævius præibat.}\par}
\par{\dictfontsize\hangindent=1.5em\hangafter=1\noindent Testament. \lxlinks{Testamentum; supremæ morientis litteræ.}\par}
\par{\dictfontsize\hangindent=1.5em\hangafter=1\noindent Texte (\textit{d’un auteur, d’une loi}). \lxlinks{Verba auctoris \textit{seu} legis.}\par}
\par{\dictfontsize\hangindent=1.5em\hangafter=1\noindent Tirer \textit{du fusil, du pistolet, du revolver}. \lxlinks{Glandes fundere.}\par}
\par{\dictfontsize\hangindent=1.5em\hangafter=1\noindent Tomber. \lxlinks{Labi.}\par}
\par{\dictfontsize\hangindent=1.5em\hangafter=1\noindent Tracer \textit{de nouvelles fortifications}. \lxlinks{Muris locum destinare. Designare mœnia, fossas.}\par}
\par{\dictfontsize\hangindent=1.5em\hangafter=1\noindent Traînards. \lxlinks{Milites itinere confecti, moratores, invalidi milites, qui agmen sequi non possunt.}\par}
\par{\dictfontsize\hangindent=1.5em\hangafter=1\noindent \textit{Payer par} traite; \textit{faire tirer une traite sur son compte}. \lxlinks{Pecuniam permutatione trajicere.}\par}
\par{\dictfontsize\hangindent=1.5em\hangafter=1\noindent Remplir les \{ clauses \} d’un traité. \{ articles \} \lxlinks{In pactione manere.}\par}
\par{\dictfontsize\hangindent=1.5em\hangafter=1\noindent Traiteur; \textit{restaurateur}. \lxlinks{Structor conviviorum publicus.}\par}
\par{\dictfontsize\hangindent=1.5em\hangafter=1\noindent Tranchées. \lxlinks{Agger militaris, fossa e terræ aggestu in alterum labrum altior ac munitior.}\par}
\par{\dictfontsize\hangindent=1.5em\hangafter=1\noindent Tranchée; \textit{approches}. \lxlinks{Accessus operum, viæ obsidionales, accessus militum, ductus obliqui, appropinquationes operum militarium.}\par}
\par{\dictfontsize\hangindent=1.5em\hangafter=1\noindent \textit{Ouvrir, conduire des} tranchées \textit{transversales, élever des traverses}. \lxlinks{Transversam obducere.}\par}
\par{\dictfontsize\hangindent=1.5em\hangafter=1\noindent \textit{Retranchement}, tranchée. \lxlinks{Lorica valli inferioris, ambulacrum valli, succinctum valli.}\par}
\par{\dictfontsize\hangindent=1.5em\hangafter=1\noindent \textit{Ouvrir, conduire des} tranchées \textit{en ligne droite}. \lxlinks{Fossam diversis lateribus ducere.}\par}
\par{\dictfontsize\hangindent=1.5em\hangafter=1\noindent \textit{Les soldats sont dans l’eau, dans les} tranchées. \lxlinks{Milites in appropinquantibus paludibus et restagnantibus aquis pene innatant.}\par}
\par{\dictfontsize\hangindent=1.5em\hangafter=1\noindent \textit{Détruire les} tranchées \textit{et les gabions}; \textit{détruire les retranchements}. \lxlinks{Apparatum et munitiones disjicere.}\par}
\par{\dictfontsize\hangindent=1.5em\hangafter=1\noindent \textit{Ouvrir les} tranchées; \textit{s’approcher des remparts en ouvrant des tranchées}. \lxlinks{Vineas agere, adductibus ad urbem, propius muros accedere.}\par}
\par{\dictfontsize\hangindent=1.5em\hangafter=1\noindent Trésor \textit{du roi}; \textit{cour des comptes}. \lxlinks{Ærarium regium, tribunal redituum principis.}\par}
\par{\dictfontsize\hangindent=1.5em\hangafter=1\noindent Trésor \textit{public}. \lxlinks{Ærarium commune.}\par}
\par{\dictfontsize\hangindent=1.5em\hangafter=1\noindent Trésor \textit{privé} (\textit{d’un roi, d’un prince}), \textit{cassette} (\textit{d’un roi}), \textit{etc.} \lxlinks{Ærarium sanctius.}\par}
\par{\dictfontsize\hangindent=1.5em\hangafter=1\noindent \textit{Épuiser le} trésor. \lxlinks{Ærarium exhaurire.}\par}
\par{\dictfontsize\hangindent=1.5em\hangafter=1\noindent \textit{Verser de l’argent au} trésor; \textit{faire un versement de fonds au trésor}. \lxlinks{In ærarium referre pecuniam.}\par}
\par{\dictfontsize\hangindent=1.5em\hangafter=1\noindent \textit{Épuisement incroyable du} trésor \textit{public}. \lxlinks{Incredibiles angustiæ pecuniæ publicæ.}\par}
\par{\dictfontsize\hangindent=1.5em\hangafter=1\noindent Trésorier \textit{militaire}. \lxlinks{Quæstor, minister impensæ bellicæ, tribunus ærarius.}\par}
\par{\dictfontsize\hangindent=1.5em\hangafter=1\noindent \textit{Officier} trésorier; \textit{tribun du trésor}; \textit{payeur de l’armée}. \lxlinks{Minister impensæ bellicæ, quæstor ærarii bellici tribunus ærarius.}\par}
\par{\dictfontsize\hangindent=1.5em\hangafter=1\noindent \textit{Fermer les} tribunaux. \lxlinks{Justitium indicere.}\par}
\par{\dictfontsize\hangindent=1.5em\hangafter=1\noindent \textit{Rouvrir les} tribunaux. \lxlinks{Justitium remittere, laxare, solvere.}\par}
\par{\dictfontsize\hangindent=1.5em\hangafter=1\noindent Tripliques (\textit{t. de palais}); \textit{réponses à des dupliques, lesquelles sont dirigées elles-mêmes contre les répliques}. \lxlinks{Secunda actoris defensio.}\par}
\par{\dictfontsize\hangindent=1.5em\hangafter=1\noindent Troupes \textit{auxiliaires}. \lxlinks{Auxilia, auxiliares milites, copiæ subsidiariæ, catervæ conductitiæ.}\par}
\par{\dictfontsize\hangindent=1.5em\hangafter=1\noindent Troupes \textit{d’élite}. \lxlinks{Lectus exercitus.}\par}
\par{\dictfontsize\hangindent=1.5em\hangafter=1\noindent \textit{Bohémien}; \textit{Égyptien}; tsigane. \lxlinks{Ægyptii errones.}\par}
\par{\dictfontsize\hangindent=0pt\noindent\textls[120]{\scshape u}\par}
\par{\dictfontsize\hangindent=1.5em\hangafter=1\noindent \textit{Équipement}; uniforme \textit{d’un soldat}. \lxlinks{Vestitus militaris.}\par}
\par{\dictfontsize\hangindent=1.5em\hangafter=1\noindent Usurper \textit{sur...., empiéter sur les droits de qqn}. \lxlinks{Jus alicujus convellere. In jure suo aliquem interpellare.}\par}
\par{\dictfontsize\hangindent=0pt\noindent\textls[120]{\scshape v}\par}
\par{\dictfontsize\hangindent=1.5em\hangafter=1\noindent \textit{Annoncer, notifier les} vacations (\textit{t. de palais}). \lxlinks{Justitium indicere.}\par}
\par{\dictfontsize\hangindent=1.5em\hangafter=1\noindent Vaisseau \textit{de guerre}. \lxlinks{Navis bellica, navis præsidiaria, longa.}\par}
\par{\dictfontsize\hangindent=1.5em\hangafter=1\noindent \textit{Lancer un} vaisseau. \lxlinks{Ex navalibus navem deducere.}\par}
\par{\dictfontsize\hangindent=1.5em\hangafter=1\noindent \textit{Bateau à} vapeur. \lxlinks{Pyroscapha.}\par}
\par{\dictfontsize\hangindent=1.5em\hangafter=1\noindent Vassal. \lxlinks{Cliens beneficiarius, fiduciarius.}\par}
\par{\dictfontsize\hangindent=1.5em\hangafter=1\noindent \textit{Être le} vassal \textit{de qqn}. \lxlinks{Jure beneficiario alicui esse obnoxium.}\par}
\par{\dictfontsize\hangindent=1.5em\hangafter=1\noindent \textit{Poser des} ventouses. \lxlinks{Cucurbitulis vitreis sanguinem extrahere.}\par}
\par{\dictfontsize\hangindent=1.5em\hangafter=1\noindent Verbaliser, \textit{enregistrer}. \lxlinks{In adversaria, acta, commentarios publicos referre.}\par}
\par{\dictfontsize\hangindent=1.5em\hangafter=1\noindent \textit{Mettre une maison en} vente. \lxlinks{Ædes proscribere.}\par}
\par{\dictfontsize\hangindent=1.5em\hangafter=1\noindent Vicaire. \lxlinks{Qui sacri Curionis explet vicem, pastoris vicarius.}\par}
\par{\dictfontsize\hangindent=1.5em\hangafter=1\noindent Ville \textit{marchande ou commerçante}. \lxlinks{Forum rerum venalium quod in hac illave urbe præstatur.}\par}
\par{\dictfontsize\hangindent=1.5em\hangafter=1\noindent Visiter \textit{les provinces}. \lxlinks{Proficisci in provincias visendi causa.}\par}
\par{\dictfontsize\hangindent=1.5em\hangafter=1\noindent Vivres, \textit{comestibles, provisions}. \lxlinks{Frumentum, res frumentaria, commeatus; alimenta.}\par}
\par{\dictfontsize\hangindent=1.5em\hangafter=1\noindent \textit{Faire des réquisitions de} vivres, \textit{en nature}. \lxlinks{Vecturas frumenti finitimis urbibus describere.}\par}
\par{\dictfontsize\hangindent=1.5em\hangafter=1\noindent \textit{Fournir les} vivres, \textit{les munitions de bouche, les provisions}; \textit{entretenir l’abondance des vivres}. \lxlinks{Commeatum supportare, sustinere.}\par}
\par{\dictfontsize\hangindent=1.5em\hangafter=1\noindent \textit{Mettre à la} voile. \lxlinks{Classe proficisci.}\par}
\par{\dictfontsize\hangindent=1.5em\hangafter=1\noindent Voix; \textit{ton, son} (\textit{de la voix}). \lxlinks{Vox, vocis sonus.}\par}
\par{\dictfontsize\hangindent=1.5em\hangafter=1\noindent Volontaires (\textit{guer.}). \lxlinks{Milites voluntarii, honorarii.}\par}
\par{\dictfontsize\hangindent=1.5em\hangafter=1\noindent Voter. \lxlinks{Suffragium ferre, inire; in suffragium ire; sententiam dicere, expromere.}\par}
\par{\dictfontsize\hangindent=1.5em\hangafter=1\noindent \textit{Faire} voter. \lxlinks{In suffragium mittere; ad suffragium vocare.}\par}
\par{\dictfontsize\hangindent=0pt\noindent\textls[120]{\scshape y}\par}
\par{\dictfontsize\hangindent=1.5em\hangafter=1\noindent Yacht; \textit{aviso}; \textit{une barque d’avis} (\textit{marine}). \lxlinks{Navis tabellaria, navis velox.}\par}
\par{\dictfontsize\hangindent=1.5em\hangafter=1\noindent Yacht, \textit{bâtiment léger, petit navire à un pont}. \lxlinks{Navis actuaria, liburnica.}\par}
\par{\dictfontsize\hangindent=0pt\noindent\textls[120]{\scshape z}\par}
\par{\dictfontsize\hangindent=1.5em\hangafter=1\noindent \textit{Fourrure de} zibeline. \lxlinks{Pellis mustelarum scythicarum.}\par}
